% Lesson 26 — Grammar: Revise clauses — Anglais Tle A — Cours signé OnBuch+
% Programme officiel MINESEC. Compiler avec : tectonic EN26-grammar-revise-clauses.tex
\newcommand{\DOCMATIERE}{Anglais}
\newcommand{\DOCNIVEAU}{Tle A}
\newcommand{\DOCMODULE}{Chapter 3 : Environment, Well-being and Health}
\newcommand{\DOCLECON}{EN26}
\newcommand{\DOCTITRE}{Lesson 26 — Grammar: Revise clauses}
% OnBuch+ — préambule « cours signés » ENRICHI (compilé avec tectonic / XeLaTeX)
% Système visuel « Pop » d'OnBuch+ : fond crème (jamais blanc), bordures encre
% épaisses, ombres « dures » décalées, titres Archivo Black en capitales.
% Version MATHÉMATIQUES : graphes (pgfplots/tikz), boîtes pédagogiques
% (définition, propriété, méthode, attention, le savais-tu, exercice résolu,
% exercice, TP, bilan), page de garde et signature OnBuch+ ancrée en bas.
%
% Macros attendues dans main.tex AVANT \input{preamble} :
%   \DOCMATIERE  \DOCNIVEAU  \DOCMODULE  \DOCLECON (numéro)  \DOCTITRE
\documentclass[11pt]{article}

\usepackage{fontspec}
\usepackage[english,french]{babel} % français = langue principale ; l'anglais via \eng{..} / textebox
\usepackage[a4paper,margin=2cm,top=3.4cm,bottom=2.8cm,headheight=1.4cm,footskip=1.3cm]{geometry}
\usepackage{amsmath,amssymb,amsfonts}
\usepackage{mathtools} % \xmapsto, \coloneqq... souvent utilisés par les modèles
\usepackage{cancel} % simplifications barrées (bilans de forces)
% \abs{z} (module, valeur absolue) et \norm{u} (norme), utilisés par les modèles
\providecommand{\abs}{}\let\abs\relax\DeclarePairedDelimiter{\abs}{\lvert}{\rvert}
\providecommand{\norm}{}\let\norm\relax\DeclarePairedDelimiter{\norm}{\lVert}{\rVert}
\usepackage[table]{xcolor} % [table] : fournit aussi \rowcolors (alternance de lignes)
\usepackage{pagecolor}
\usepackage{tikz}
\usetikzlibrary{arrows.meta,positioning,calc,shapes.geometric,shapes.misc,shapes.arrows,decorations.pathmorphing,decorations.pathreplacing,decorations.markings,patterns,shadows,mindmap,trees,fit,backgrounds,matrix,intersections,angles,quotes}
\usepackage{pgfplots}
\usepackage[version=4]{mhchem} % équations nucléaires \ce{^{235}_{92}U}
\usepackage[european,siunitx]{circuitikz} % schémas électriques
\pgfplotsset{compat=1.18}
\usepgfplotslibrary{fillbetween,groupplots} % aires entre courbes (calcul intégral)
\usepackage{enumitem}
\usepackage{fancyhdr}
% [most] charge toutes les bibliothèques tcolorbox (listings, minted, raster,
% documentation...) et gonfle inutilement la mémoire TeX sur un document long
% (~300 boîtes) : on ne charge que ce qui est réellement utilisé.
\usepackage[skins,breakable]{tcolorbox}
\usepackage{graphicx}
\usepackage{colortbl}
\usepackage{booktabs}
\usepackage{tabularx}
\usepackage{diagbox} % case d'en-tête diagonale des tableaux à double entrée
% Colonnes extensibles centrée / gauche / droite (C, L, R), souvent utilisées par les modèles
\newcolumntype{C}{>{\centering\arraybackslash}X}
\newcolumntype{L}{>{\raggedright\arraybackslash}X}
\newcolumntype{R}{>{\raggedleft\arraybackslash}X}
\usepackage{array}
\usepackage{multicol}
\usepackage{multirow}
\usepackage{siunitx}
% parse-numbers=false : accepte aussi \SI{2,5 \times 10^{-3}}{...}
\sisetup{locale=FR,output-decimal-marker={,},group-minimum-digits=4,parse-numbers=false}
\DeclareSIUnit\ohm{\ensuremath{\symup{\Omega}}} % U+2126 absent de la police
\DeclareSIUnit\division{div} % réglages d'oscilloscope (ms/div, V/div)
\DeclareSIUnit\tr{tr} % tours (vitesse de rotation, ex. \SI{1500}{\tr\per\minute})
\DeclareSIUnit\molar{M} % concentration molaire (ex. \SI{3}{\molar}, \SI{1}{\micro\molar})
\DeclareSIUnit\an{an} % année (le modèle confond parfois avec \year, un compteur TeX) — ex. \SI{5}{\kilo\meter\per\an}
\DeclareSIUnit\FCFA{FCFA} % franc CFA (ex. \SI{500}{\FCFA\per\kilo\gram})
\DeclareSIUnit\degreeBrix{\textdegree Brix} % teneur en sucre d'un sirop/jus (réfractométrie)
\DeclareSIUnit\month{mois} % le modèle utilise parfois \month, un compteur TeX, comme unité de durée
\usepackage{qrcode}
% \degree : commande fréquemment utilisée par les modèles pour un angle ou une
% température en dehors de \SI{}{} (non fournie par les paquets chargés).
\providecommand{\degree}{^{\circ}}
% \qed (fin de démonstration), utilisé par les modèles sans amsthm
\providecommand{\qed}{\hfill$\square$}
% \barre : double barre des valeurs interdites dans un tableau de variations (tkz-tab)
\providecommand{\barre}{\ensuremath{\|}}
% environnement proof (amsthm non chargé) utilisé par les modèles
\ifdefined\proof\else\newenvironment{proof}[1][Démonstration]{\par\noindent\textit{#1.}\ }{\qed\par}\fi
\usepackage{lastpage}
\usepackage{hyperref}
\hypersetup{hidelinks}

% ============================================================
% Palette « Pop » OnBuch+ — mêmes hex que la classe P (plus_ui.dart)
% ============================================================
\definecolor{poporange}{HTML}{F59321}
\definecolor{popdark}{HTML}{E07A0C}
\definecolor{popcream}{HTML}{FFF6E8}
\definecolor{popink}{HTML}{1C1714}
\definecolor{poppurple}{HTML}{7A5AE0}
\definecolor{popgreen}{HTML}{2E9E6A}
\definecolor{poppink}{HTML}{E0457B}
\definecolor{popblue}{HTML}{2F7DE1}
\definecolor{popgold}{HTML}{C98A12}
\definecolor{popmuted}{HTML}{8A7F76}
\definecolor{popline}{HTML}{EADFD2}
\definecolor{popgray}{HTML}{8A7F76}
\colorlet{popgrey}{popgray}
% Alias tolérants (le modèle nomme parfois les couleurs différemment)
\colorlet{popwhite}{white}
\colorlet{popblack}{popink}
\colorlet{popred}{poppink}
\colorlet{popyellow}{popgold}
% Teintes claires (fonds de tableaux/encadrés) — le modèle nomme parfois
% les couleurs avec un suffixe L (ex. popblueL) sans les avoir définies.
\colorlet{poporangeL}{poporange!20}
\colorlet{popdarkL}{popdark!20}
\colorlet{poppurpleL}{poppurple!20}
\colorlet{popgreenL}{popgreen!20}
\colorlet{poppinkL}{poppink!20}
\colorlet{popblueL}{popblue!20}
\colorlet{popgoldL}{popgold!20}
\colorlet{popredL}{popred!20}
\colorlet{popgrayL}{popgray!20}
% Teintes claires pour fonds de tableaux / zones
\colorlet{popblueL}{popblue!10!white}
\colorlet{poporangeL}{poporange!14!white}
\colorlet{popgreenL}{popgreen!12!white}
\colorlet{poppurpleL}{poppurple!12!white}
\colorlet{poppinkL}{poppink!12!white}
\colorlet{popgoldL}{popgold!14!white}

\pagecolor{popcream}

% ============================================================
% Typographie
% ============================================================
\defaultfontfeatures{Ligatures=TeX}
\setmainfont{PlusJakartaSans-Medium.ttf}[Path=../../../fonts/,BoldFont=PlusJakartaSans-Bold.ttf]
% Symboles, API et emojis absents de la police principale : repli sur DejaVu Sans (caractères actifs)
\newfontface\symfont{DejaVuSans.ttf}[Path=../../../fonts/]
\makeatletter
\newcommand\symactive[1]{\catcode#1=13 \begingroup\lccode`\~=#1 \lowercase{\endgroup\def~}{{\symfont\char#1}}}
\newcommand\symblank[1]{\catcode#1=13 \begingroup\lccode`\~=#1 \lowercase{\endgroup\def~}{}}
\newcommand\symhyph[1]{\catcode#1=13 \begingroup\lccode`\~=#1 \lowercase{\endgroup\def~}{\mbox{-}}}
\newcount\symc
\newcommand\symrange[2]{\symc=#1 \@whilenum\symc<#2 \do{\expandafter\symactive\expandafter{\the\symc}\advance\symc by 1 }}
\newcommand\symblankrange[2]{\symc=#1 \@whilenum\symc<#2 \do{\expandafter\symblank\expandafter{\the\symc}\advance\symc by 1 }}
\makeatother
\symrange{"0250}{"02FF}\symactive{"2713}\symactive{"2714}\symactive{"2717}\symactive{"2718}\symactive{"2610}\symactive{"2611}\symactive{"2612}
\symhyph{"2011}\symhyph{"2E1A}\symblankrange{"1F300}{"1FAFF}\symblank{"FE0F}
% Maths en sans-serif (Fira Math) avec chiffres et lettres droites de Plus
% Jakarta, pour que les formules se fondent dans le texte.
\usepackage[math-style=ISO,bold-style=ISO]{unicode-math}
\setmathfont{FiraMath-Regular.otf}[Path=../../../fonts/]
\setmathfont{PlusJakartaSans-Medium.ttf}[Path=../../../fonts/,range={up/num,up/{Latin,latin},\mathup}]
\usepackage{icomma} % virgule décimale sans espace en mode math
% Caractères mathématiques Unicode absents des polices de texte (Plus Jakarta,
% Archivo Black) : rendus par leur équivalent mathématique, en texte comme en
% formule — sinon ils s'affichent en carré vide (ex. « Arithmétique dans ℤ »).
\usepackage{newunicodechar}
\newunicodechar{ℝ}{\ensuremath{\mathbb{R}}}
\newunicodechar{ℤ}{\ensuremath{\mathbb{Z}}}
\newunicodechar{ℂ}{\ensuremath{\mathbb{C}}}
\newunicodechar{ℕ}{\ensuremath{\mathbb{N}}}
\newunicodechar{ℚ}{\ensuremath{\mathbb{Q}}}
\newunicodechar{𝔻}{\ensuremath{\mathbb{D}}}
\newunicodechar{α}{\ensuremath{\alpha}}
\newunicodechar{β}{\ensuremath{\beta}}
\newunicodechar{γ}{\ensuremath{\gamma}}
\newunicodechar{δ}{\ensuremath{\delta}}
\newunicodechar{ε}{\ensuremath{\varepsilon}}
\newunicodechar{θ}{\ensuremath{\theta}}
\newunicodechar{λ}{\ensuremath{\lambda}}
\newunicodechar{μ}{\ensuremath{\mu}}
\newunicodechar{σ}{\ensuremath{\sigma}}
\newunicodechar{τ}{\ensuremath{\tau}}
\newunicodechar{φ}{\ensuremath{\varphi}}
\newunicodechar{ψ}{\ensuremath{\psi}}
\newunicodechar{ω}{\ensuremath{\omega}}
\newunicodechar{Σ}{\ensuremath{\Sigma}}
% majuscules produites par \MakeUppercase dans les titres (α-aminés -> Α)
\newunicodechar{Α}{\ensuremath{\alpha}}
\newunicodechar{Β}{\ensuremath{\beta}}
\newunicodechar{Γ}{\ensuremath{\Gamma}}
\newunicodechar{Δ}{\ensuremath{\Delta}}
\newunicodechar{Ε}{\ensuremath{\varepsilon}}
\newunicodechar{Θ}{\ensuremath{\Theta}}
\newunicodechar{Λ}{\ensuremath{\Lambda}}
\newunicodechar{Μ}{\ensuremath{\mu}}
\newunicodechar{Τ}{\ensuremath{\tau}}
\newunicodechar{Φ}{\ensuremath{\Phi}}
\newunicodechar{Ψ}{\ensuremath{\Psi}}
\newunicodechar{Ω}{\ensuremath{\Omega}}
\newunicodechar{↦}{\ensuremath{\mapsto}}
\newunicodechar{∈}{\ensuremath{\in}}
\newunicodechar{∉}{\ensuremath{\notin}}
\newunicodechar{∧}{\ensuremath{\wedge}}
\newunicodechar{⇔}{\ensuremath{\Leftrightarrow}}
\newunicodechar{⟺}{\ensuremath{\Longleftrightarrow}}
\newunicodechar{⇒}{\ensuremath{\Rightarrow}}
\newunicodechar{⟹}{\ensuremath{\Longrightarrow}}
\newunicodechar{∘}{\ensuremath{\circ}}
\newunicodechar{∪}{\ensuremath{\cup}}
\newunicodechar{∩}{\ensuremath{\cap}}
\newunicodechar{≡}{\ensuremath{\equiv}}
\newunicodechar{⊂}{\ensuremath{\subset}}
\newunicodechar{⊥}{\ensuremath{\perp}}
\newunicodechar{∃}{\ensuremath{\exists}}
\newunicodechar{∀}{\ensuremath{\forall}}
\newunicodechar{∛}{\ensuremath{\sqrt[3]{\,}}}
\newunicodechar{∜}{\ensuremath{\sqrt[4]{\,}}}
\newunicodechar{⃗}{\ensuremath{{}^{\to}}}
\newunicodechar{ⁿ}{\ifmmode^{n}\else\textsuperscript{\ensuremath{n}}\fi}
\newunicodechar{ˣ}{\ifmmode^{x}\else\textsuperscript{\ensuremath{x}}\fi}
\newunicodechar{ᵇ}{\ifmmode^{b}\else\textsuperscript{\ensuremath{b}}\fi}
\newunicodechar{ʳ}{\ifmmode^{r}\else\textsuperscript{\ensuremath{r}}\fi}
\newunicodechar{ᵘ}{\ifmmode^{u}\else\textsuperscript{\ensuremath{u}}\fi}
\newunicodechar{ʸ}{\ifmmode^{y}\else\textsuperscript{\ensuremath{y}}\fi}
\newunicodechar{ᵃ}{\ifmmode^{a}\else\textsuperscript{\ensuremath{a}}\fi}
\newunicodechar{ᵉ}{\ifmmode^{e}\else\textsuperscript{\ensuremath{e}}\fi}
\newunicodechar{ᵏ}{\ifmmode^{k}\else\textsuperscript{\ensuremath{k}}\fi}
\newunicodechar{ᵖ}{\ifmmode^{p}\else\textsuperscript{\ensuremath{p}}\fi}
\newunicodechar{ᴺ}{\ifmmode^{N}\else\textsuperscript{\ensuremath{N}}\fi}
\newunicodechar{ᶜ}{\ifmmode^{c}\else\textsuperscript{\ensuremath{c}}\fi}
\newunicodechar{ʰ}{\ifmmode^{h}\else\textsuperscript{\ensuremath{h}}\fi}
\newunicodechar{ᵠ}{\ifmmode^{q}\else\textsuperscript{\ensuremath{q}}\fi}
\newunicodechar{ₙ}{\ifmmode_{n}\else\textsubscript{\ensuremath{n}}\fi}
\newunicodechar{ₐ}{\ifmmode_{a}\else\textsubscript{\ensuremath{a}}\fi}
\newunicodechar{ᵢ}{\ifmmode_{i}\else\textsubscript{\ensuremath{i}}\fi}
\newunicodechar{ᵣ}{\ifmmode_{r}\else\textsubscript{\ensuremath{r}}\fi}
\newunicodechar{ₖ}{\ifmmode_{k}\else\textsubscript{\ensuremath{k}}\fi}
\newunicodechar{ᵦ}{\ifmmode_{\beta}\else\textsubscript{\ensuremath{\beta}}\fi}
\newunicodechar{ₓ}{\ifmmode_{x}\else\textsubscript{\ensuremath{x}}\fi}
\newfontfamily\popdisplay{ArchivoBlack-Regular.ttf}[Path=../../../fonts/]
\newfontfamily\poplogo{SpaceGrotesk-Bold.ttf}[Path=../../../fonts/]
\newfontfamily\popsemi{PlusJakartaSans-SemiBold.ttf}[Path=../../../fonts/]
\newfontfamily\popxbold{PlusJakartaSans-ExtraBold.ttf}[Path=../../../fonts/]
\newfontfamily\popxboldls{PlusJakartaSans-ExtraBold.ttf}[Path=../../../fonts/,LetterSpace=3.0]

\setlist[enumerate]{leftmargin=1.7em,itemsep=5pt,topsep=4pt}
\setlist[itemize]{leftmargin=1.7em,itemsep=4pt,topsep=4pt,label={\color{poporange}\textbullet}}
\setlength{\parindent}{0pt}
\setlength{\parskip}{5pt}
\renewcommand{\arraystretch}{1.25}

% pgfplots : style Pop par défaut
\pgfplotsset{
  every axis/.append style={axis line style={popink,line width=1pt},tick style={popink},
    grid style={popline,line width=0.5pt},label style={font=\small},tick label style={font=\footnotesize}},
  popcurve/.style={poporange,line width=1.8pt,no markers,smooth},
}
% Même style disponible dans un \draw TikZ simple (hors pgfplots)
\tikzset{popcurve/.style={poporange,line width=1.8pt}}
\tikzset{popgrid/.style={popline,very thin}}
\colorlet{popcurve}{poporange} % le nom du style est parfois utilisé comme couleur

% ============================================================
% Pill / squiggle
% ============================================================
\newcommand{\poppill}[3]{%
  \tikz[baseline=-0.55ex]\node[draw=popink,line width=1.2pt,rounded corners=2.6pt,%
    fill=#2,inner xsep=6.5pt,inner ysep=3pt]{%
    \popxboldls\fontsize{7}{7}\selectfont\color{#3}\MakeUppercase{#1}};%
}
\newcommand{\popsquiggle}[1]{%
  \tikz[baseline=-0.4ex]{
    \draw[#1,line width=2.6pt,line cap=round]
      plot [smooth] coordinates {(0,0.09) (0.34,0.01) (0.68,0.21) (1.02,0.01) (1.36,0.10)};
  }%
}
% Mot-clé surligné (marqueur orange sous le texte)
\newcommand{\cle}[1]{{\popxbold\color{popdark}#1}}

% ============================================================
% En-tête / pied
% ============================================================
\pagestyle{fancy}
\fancyhf{}
\renewcommand{\headrulewidth}{0pt}
\renewcommand{\footrulewidth}{0pt}
\lhead{\raisebox{-0.15ex}{\poplogo\fontsize{13}{13}\selectfont\color{popink}OnBuch\color{poporange}+}}
\rhead{\poppill{\DOCMATIERE\ \textbullet\ \DOCNIVEAU\ \textbullet\ Leçon \DOCLECON}{white}{popink}}
\lfoot{\popxbold\fontsize{7.6}{7.6}\selectfont\color{popmuted}\MakeUppercase{Cours signé OnBuch+}\ \textbullet\ {\popsemi\DOCTITRE}}
\rfoot{\tikz[baseline=-0.6ex]\node[rounded corners=8.5pt,fill=popink,minimum height=17pt,minimum width=17pt,inner xsep=5pt,inner sep=0]{\popxbold\fontsize{8}{8}\selectfont\color{popcream}\thepage\,/\,\pageref*{LastPage}};}

% ============================================================
% Titres
% ============================================================
\newcommand{\chaptitre}[2]{%
  \begin{tcolorbox}[enhanced,colback=poporange,colframe=popink,boxrule=2.6pt,arc=11pt,
      drop shadow=popink,left=15pt,right=15pt,top=12pt,bottom=11pt,before skip=3pt,after skip=18pt]
    \poppill{#2}{popink}{popcream}\par\vspace{9pt}
    {\raggedright\hyphenpenalty=10000\exhyphenpenalty=10000\popdisplay\fontsize{22}{24}\selectfont\color{popcream}\MakeUppercase{#1}\par}
  \end{tcolorbox}
}
% Section numérotée : pastille orange avec le numéro + titre Archivo
\newcounter{popsec}
\newcommand{\coursec}[1]{%
  \stepcounter{popsec}\setcounter{popsub}{0}%
  \par\vspace{16pt}\noindent
  \begin{minipage}{\linewidth}
  \tikz[baseline=-0.7ex]\node[draw=popink,line width=1.6pt,fill=poporange,rounded corners=4pt,minimum size=22pt,inner sep=0]{\popdisplay\fontsize{12}{12}\selectfont\color{popcream}\thepopsec};%
  \hspace{8pt}{\popdisplay\fontsize{15}{17}\selectfont\color{popink}\MakeUppercase{#1}}\par
  \vspace{4pt}\noindent{\color{poporange}\rule{36pt}{3.6pt}}
  \end{minipage}\par\nopagebreak\vspace{8pt}
}
\newcounter{popsub}
\newcommand{\courssub}[1]{%
  \stepcounter{popsub}%
  \par\vspace{9pt}\noindent{\popxbold\fontsize{11.5}{13}\selectfont\color{popink}{\color{poporange}\thepopsec.\thepopsub}\hspace{6pt}#1}\par\nopagebreak\vspace{4pt}
}

% ============================================================
% Boîtes pédagogiques — même recette : fond blanc, bordure épaisse colorée,
% ombre dure encre, badge plein. Titre optionnel [#1].
% ============================================================
\tcbset{popbase/.style={breakable,enhanced,colback=white,boxrule=2.3pt,arc=9pt,
  shadow={3pt}{-3pt}{0pt}{popink},left=14pt,right=14pt,top=11pt,bottom=11pt,before skip=11pt,after skip=15pt,
  fontupper=\popsemi\fontsize{10.6}{14.5}\selectfont\color{popink},
  fontlower=\popsemi\fontsize{10.6}{14.5}\selectfont\color{popink}}}
% Coche dessinée (le glyphe ✓ n'existe pas dans les polices du document)
\newcommand{\popcheck}[1]{\tikz[baseline=-0.5ex]\draw[#1,line width=1.6pt,line cap=round,line join=round](-0.13,0.0)--(-0.04,-0.09)--(0.13,0.1);}
\providecommand{\coche}{\popcheck{popgreen}}
\providecommand{\resultat}[1]{\par\smallskip\noindent\fcolorbox{popgreen}{popgreenL}{\parbox{\dimexpr\linewidth-2\fboxsep-2\fboxrule\relax}{\textbf{Résultat.} #1}}\par\smallskip}
\newcommand{\popboxhead}[3]{\poppill{#1}{#2}{white}\ifx\relax#3\relax\else\hspace{6pt}{\popxbold\fontsize{10.6}{12}\selectfont\color{popink}#3}\fi\par\vspace{7pt}}

\newtcolorbox{aretenir}[1][]{popbase,colframe=popgreen,before upper={\popboxhead{À retenir}{popgreen}{#1}}}
% Texte en anglais (document de lecture, dialogue, extrait) : cadre
% d'étude. Un titre optionnel (source, auteur) s'affiche dans l'en-tête.
\newtcolorbox{textebox}[1][]{popbase,colframe=popblue,before upper={\popboxhead{Text}{popblue}{#1}\begin{otherlanguage}{english}},after upper={\end{otherlanguage}}}
% Marques ✓ / ✗ (forme correcte / fautive) souvent utilisées par les modèles
\providecommand{\cross}{\textcolor{poppink}{\ensuremath{\times}}}
\providecommand{\xmark}{\cross}\providecommand{\crossmark}{\cross}\providecommand{\cmark}{\ensuremath{\checkmark}}
% Ligne de réponse / justification à compléter (inventée par les modèles)
\providecommand{\justification}[1]{\par\noindent\textit{Justification~:} \dotfill\par}
% \textipa (package tipa non chargé) : la police ne couvre pas l'API -> simple passage
\providecommand{\textipa}[1]{#1}
% Soulignement (ulem non chargé) : \uline, \ul
\providecommand{\uline}[1]{\underline{#1}}\providecommand{\ul}[1]{\underline{#1}}
% \glossaire{\item ...} inventée par les modèles : liste de glossaire
\providecommand{\glossaire}[1]{\begin{itemize}#1\end{itemize}}
% \alert (beamer) inventée par les modèles : texte mis en évidence
\providecommand{\alert}[1]{\textbf{\textcolor{poppink}{#1}}}
% variantes ASCII d'ordinaux écrites en macro
\providecommand{\iere}{\textsuperscript{re}}\providecommand{\ieme}{\textsuperscript{e}}\providecommand{\ere}{\textsuperscript{re}}\providecommand{\eme}{\textsuperscript{e}}\providecommand{\er}{\textsuperscript{er}}\providecommand{\ie}{\textsuperscript{e}}
% Ordinaux français écrits en macro par les modèles : 1\ière, 3\ième
\newcommand{\ière}{\textsuperscript{re}}\newcommand{\ième}{\textsuperscript{e}}
% \exemple (inventée par les modèles) : étiquette « Exemple : »
\providecommand{\exemple}{\textbf{Exemple~:}\ }
% \square utilisable aussi hors mode math (cases à cocher)
\AtBeginDocument{\let\squareMath\square\renewcommand{\square}{\ensuremath{\squareMath}}}
% Mot ou phrase en anglais à l'intérieur d'un paragraphe français
\newcommand{\eng}[1]{\ifmmode\text{\textit{#1}}\else\begingroup\selectlanguage{english}\itshape #1\endgroup\fi}
\let\esp\eng % alias toléré
% flèches d'intonation inventées par les modèles
\providecommand{\textsearrow}{}\renewcommand{\textsearrow}{\ensuremath{\searrow}}\providecommand{\textnearrow}{}\renewcommand{\textnearrow}{\ensuremath{\nearrow}}
% \fr{..} : mot français (inventé par les modèles)
\providecommand{\fr}[1]{\textit{#1}}
\newtcolorbox{exemplebox}[1][]{popbase,colframe=poporange,before upper={\popboxhead{Exemple}{poporange}{#1}}}
\newtcolorbox{definition}[1][]{popbase,colframe=popblue,before upper={\popboxhead{Définition}{popblue}{#1}}}
\newtcolorbox{propriete}[1][]{popbase,colframe=poppurple,before upper={\popboxhead{Propriété}{poppurple}{#1}}}
\newtcolorbox{methode}[1][]{popbase,colframe=popgold,before upper={\popboxhead{Méthode}{popgold}{#1}}}
\newtcolorbox{attention}[1][]{popbase,colframe=poppink,before upper={\popboxhead{Attention}{poppink}{#1}}}
\newtcolorbox{experience}[1][]{popbase,colframe=popblue,colback=popblueL,before upper={\popboxhead{Expérience}{popblue}{#1}}}
% « Le savais-tu ? » : carte violette pleine (contexte, culture, Cameroun)
\newtcolorbox{savaistu}[1][]{popbase,colback=poppurple,colframe=popink,
  fontupper=\popsemi\fontsize{10.4}{14.2}\selectfont\color{white},
  before upper={\poppill{Le savais-tu\,?}{white}{poppurple}\ifx\relax#1\relax\else\hspace{6pt}{\popxbold\color{white}#1}\fi\par\vspace{7pt}}}
% Exercice résolu : énoncé en haut, \tcblower puis la solution sur fond crème
\newcounter{popexo}\newcounter{popexr}
\newtcolorbox{exoresolu}[1][]{popbase,colframe=poporange,colbacklower=popcream,
  segmentation style={popink,line width=1.2pt,dashed},
  before upper={\stepcounter{popexr}\popboxhead{Exercice résolu \thepopexr}{poporange}{#1}},
  before lower={\poppill{Solution}{popink}{popcream}\par\vspace{6pt}}}
% Exercice (non résolu en place) — la correction est dans la partie Corrigés
\newtcolorbox{exercice}[1][]{popbase,colframe=popink,
  before upper={\stepcounter{popexo}\popboxhead{Exercice \thepopexo}{popink}{#1}}}
% Corrigé
\newtcolorbox{corrige}[1][]{popbase,colframe=popgreen,colback=popgreenL,
  before upper={\popboxhead{Corrigé}{popgreen}{#1}}}
% Objectifs / prérequis / bilan
\newtcolorbox{objectifs}{popbase,colframe=popink,colback=poporangeL,
  before upper={\popboxhead{Objectifs de la leçon}{popink}{}}}
\newtcolorbox{prerequis}{popbase,colframe=popink,colback=popblueL,
  before upper={\popboxhead{Prérequis}{popblue}{}}}
\newtcolorbox{bilan}{popbase,colframe=popink,colback=white,boxrule=2.8pt,
  before upper={\popboxhead{Fiche bilan}{poporange}{L'essentiel à connaître pour l'examen}}}
% Figure encadrée avec légende
\newtcolorbox{popfigure}[1][]{enhanced,breakable=false,colback=white,colframe=popline,boxrule=1.4pt,arc=8pt,
  left=10pt,right=10pt,top=10pt,bottom=8pt,before skip=10pt,after skip=12pt,halign=center,
  lower separated=false,halign lower=center,
  fontlower=\popsemi\fontsize{9}{11.5}\selectfont\color{popmuted},
  #1}
\newcommand{\legende}[1]{\par\vspace{6pt}{\popsemi\fontsize{9}{11.5}\selectfont\color{popmuted}#1}}

% ============================================================
% Signature OnBuch+
% ============================================================
\newcommand{\poplogoinline}[1]{{\poplogo\fontsize{#1}{#1}\selectfont\color{popink}OnBuch\color{poporange}+}}
\newcommand{\signatureonbuch}{%
  \begin{tcolorbox}[enhanced,colback=popink,colframe=popink,boxrule=2.2pt,arc=10pt,
      drop shadow=poporange,left=14pt,right=14pt,top=10pt,bottom=10pt,before skip=0pt,after skip=0pt]
    \tikz[baseline=-0.7ex]\node[circle,fill=popgreen,draw=popcream,line width=1.3pt,minimum size=22pt,inner sep=0]{\popcheck{white}};%
    \hspace{9pt}\begin{minipage}[c]{0.62\linewidth}
      {\popxbold\fontsize{12}{13}\selectfont\color{popcream}Signé\ \textbullet\ {\poplogo OnBuch\color{poporange}+}}\par\vspace{2pt}
      {\raggedright\popsemi\fontsize{8}{10.5}\selectfont\color{popline}\MakeUppercase{Équipe pédagogique OnBuch+}\\\MakeUppercase{Conforme au programme officiel MINESEC}\par}
    \end{minipage}\hfill
    {\poplogo\fontsize{22}{22}\selectfont\color{popcream}OnBuch\color{poporange}+}
  \end{tcolorbox}%
}

% ============================================================
% Page de garde
% #1 titre · #2 matière · #3 niveau · #4 module · #5 n° leçon · #6 durée ·
% #7 description / objectifs (texte ou itemize)
% ============================================================
\newcommand{\pagegarde}[7]{%
  \thispagestyle{empty}
  \newgeometry{a4paper,margin=1.7cm,top=1.6cm,bottom=1.4cm}%
  \begin{tikzpicture}[remember picture,overlay]
    \fill[poporange!18] ([xshift=-3.2cm,yshift=-3.6cm]current page.north east) circle (4.6cm);
    \fill[poppurple!14] ([xshift=2.4cm,yshift=7.6cm]current page.south west) circle (3.8cm);
  \end{tikzpicture}%
  {\poplogo\fontsize{30}{30}\selectfont\color{popink}OnBuch\color{poporange}+}\hfill
  \raisebox{0.6ex}{\poppill{Cours signé \textbullet\ Premium}{popink}{popcream}}\par
  \vspace{1pt}\popsquiggle{poppurple}\par
  \vspace{8pt}
  \begin{tcolorbox}[enhanced,colback=poporange,colframe=popink,boxrule=2.8pt,arc=13pt,
      drop shadow=popink,left=18pt,right=18pt,top=12pt,bottom=13pt,after skip=10pt]
    \poppill{#2 \textbullet\ #3}{popink}{popcream}\hspace{5pt}\poppill{#4}{white}{popink}\par\vspace{8pt}
    {\popdisplay\fontsize{14}{14}\selectfont\color{popink}LEÇON #5}\par\vspace{4pt}
    {\raggedright\hyphenpenalty=10000\exhyphenpenalty=10000\popdisplay\fontsize{25}{27}\selectfont\color{popcream}\MakeUppercase{#1}\par}
    \vspace{8pt}
    {\popxbold\fontsize{9.5}{9.5}\selectfont\color{popink}\MakeUppercase{Durée conseillée : #6}}
  \end{tcolorbox}
  \begin{tcolorbox}[enhanced,colback=white,colframe=popink,boxrule=2.2pt,arc=10pt,
      drop shadow=popink,left=16pt,right=16pt,top=10pt,bottom=10pt,after skip=8pt]
    \poppill{Dans cette leçon}{poppurple}{white}\par\vspace{5pt}
    \popsemi\fontsize{9.3}{12.6}\selectfont\color{popink}#7
  \end{tcolorbox}
  \noindent\begin{minipage}[t]{0.487\linewidth}
    \begin{tcolorbox}[enhanced,colback=popgreen,colframe=popink,boxrule=2.2pt,arc=10pt,
        drop shadow=popink,left=11pt,right=11pt,top=10pt,bottom=10pt,height=3.1cm,valign=center]
      \tcbox[colback=white,colframe=popink,boxrule=1.3pt,arc=5pt,boxsep=0pt,nobeforeafter,
        left=3pt,right=3pt,top=3pt,bottom=3pt,valign=center]{\qrcode[height=1.9cm]{https://play.google.com/store/apps/details?id=com.luvvix.onbuch&hl=fr}}%
      \hspace{8pt}\begin{minipage}[c]{0.5\linewidth}
        {\popdisplay\fontsize{11}{12.5}\selectfont\color{white}\MakeUppercase{Télécharge OnBuch}}\par\vspace{4pt}
        {\raggedright\popsemi\fontsize{8.4}{11}\selectfont\color{white}Gratuit \textbullet\ Google Play\\Tuteur IA, annales, concours, résultats\par}
      \end{minipage}
    \end{tcolorbox}
  \end{minipage}\hfill
  \begin{minipage}[t]{0.487\linewidth}
    \begin{tcolorbox}[enhanced,colback=poppurple,colframe=popink,boxrule=2.2pt,arc=10pt,
        drop shadow=popink,left=11pt,right=11pt,top=10pt,bottom=10pt,height=3.1cm,valign=center]
      \tikz[baseline=-0.7ex]\node[circle,fill=white,draw=popink,line width=1.3pt,minimum size=1.6cm,inner sep=0]{\popdisplay\fontsize{24}{24}\selectfont\color{poppurple}+};%
      \hspace{8pt}\begin{minipage}[c]{0.6\linewidth}
        {\popdisplay\fontsize{11}{12.5}\selectfont\color{white}\MakeUppercase{Passe à OnBuch+}}\par\vspace{4pt}
        {\raggedright\popsemi\fontsize{8.4}{11}\selectfont\color{white}Tous les cours signés, Tuteur IA illimité, fiches PDF — onglet OnBuch+ de l'app\par}
      \end{minipage}
    \end{tcolorbox}
  \end{minipage}
  \vspace{8pt}
  \popcontenu
  \vfill
  \signatureonbuch
  \restoregeometry
  \newpage
}

% Rangée « Ce cours contient » (page de garde)
\newcommand{\poptile}[3]{% #1 couleur #2 gros texte #3 légende
  \begin{tcolorbox}[enhanced,colback=#1,colframe=popink,boxrule=2pt,arc=9pt,shadow={2.5pt}{-2.5pt}{0pt}{popink},
      left=6pt,right=6pt,top=9pt,bottom=9pt,halign=center,valign=center,height=1.75cm,width=\linewidth,nobeforeafter]
    {\popdisplay\fontsize{12.5}{13}\selectfont\color{white}\MakeUppercase{#2}}\par\vspace{3pt}
    {\centering\popsemi\fontsize{7.8}{9.5}\selectfont\color{white}#3\par}
  \end{tcolorbox}}
\newcommand{\popcontenu}{%
  \par\noindent{\popxboldls\fontsize{7.5}{7.5}\selectfont\color{popmuted}CE COURS SIGNÉ CONTIENT}\par\vspace{7pt}
  \noindent\begin{minipage}[t]{0.235\linewidth}\poptile{popblue}{Cours}{complet, illustré, pas à pas}\end{minipage}\hfill
  \begin{minipage}[t]{0.235\linewidth}\poptile{poporange}{Méthodes}{et exercices résolus}\end{minipage}\hfill
  \begin{minipage}[t]{0.235\linewidth}\poptile{poppink}{Exercices}{type examen + corrigés}\end{minipage}\hfill
  \begin{minipage}[t]{0.235\linewidth}\poptile{popgreen}{Fiche bilan}{l'essentiel à retenir}\end{minipage}\par}

% Sommaire
\newtcolorbox{sommaire}{enhanced,colback=white,colframe=popink,boxrule=2.2pt,arc=10pt,
  drop shadow=popink,left=18pt,right=18pt,top=13pt,bottom=13pt,before skip=6pt,after skip=20pt,
  before upper={\poppill{Sommaire}{poppurple}{white}\par\vspace{9pt}\popsemi\fontsize{10.6}{14.5}\selectfont\color{popink}}}

% Signature de fin de cours — poussée tout en bas de la dernière page
\newcommand{\findecours}{%
  \par\vfill\nopagebreak
  \begin{center}\popsquiggle{poporange}\end{center}\vspace{4pt}
  \signatureonbuch
}
\AtBeginDocument{\def\llbracket{\mathopen{[\mkern-3.5mu[}}\def\rrbracket{\mathclose{]\mkern-3.5mu]}}} % intervalles d'entiers (glyphe ⟦ absent de la police)
% Glyphes absents de Fira Math : redéfinis à partir de glyphes disponibles
\AtBeginDocument{%
  \def\setminus{\mathbin{\backslash}}\let\smallsetminus\setminus
  \def\vdots{\mathord{\vbox{\baselineskip3.2pt\lineskiplimit0pt\kern4pt\hbox{.}\hbox{.}\hbox{.}}}}%
  \def\ddots{\mathinner{\mkern1mu\raise6.5pt\hbox{.}\mkern2mu\raise3.6pt\hbox{.}\mkern2mu\raise0.7pt\hbox{.}\mkern1mu}}%
  \def\triangle{\mathord{\scalebox{0.9}{$\symup{\Delta}$}}}%
  \def\nabla{\mathord{\raisebox{\depth}{\scalebox{1}[-1]{$\symup{\Delta}$}}}}%
  \def\checkmark{\popcheck{popdark}}%
  \def\star{\mathbin{\ast}}\def\bigstar{\mathord{\ast}}%
  \def\diamond{\mathbin{\scalebox{0.6}{$\Diamond$}}}%
  \def\bigcup{\mathop{\mathchoice{\raisebox{-0.3em}{\scalebox{1.7}{$\cup$}}}{\scalebox{1.25}{$\cup$}}{\cup}{\cup}}\displaylimits}%
  \def\bigcap{\mathop{\mathchoice{\raisebox{-0.3em}{\scalebox{1.7}{$\cap$}}}{\scalebox{1.25}{$\cap$}}{\cap}{\cap}}\displaylimits}%
  \def\lesseqqgtr{\mathrel{\lessgtr}}\def\bigoplus{\mathop{\oplus}\displaylimits}%
}
\providecommand{\ssi}{si et seulement si} % abréviation fréquente
\AtBeginDocument{\providecommand{\wideparen}{\overparen}} % arc de cercle

%% FIGFIX-BEGIN v3 — correctifs globaux des figures (docs/onbuch-generation/tools/figfix)
\usepackage{adjustbox}\usepackage{etoolbox}
% toute figure TikZ/pgfplots plus large que la ligne est réduite à la largeur disponible
\BeforeBeginEnvironment{tikzpicture}{\begin{adjustbox}{max width=\linewidth}}
\AfterEndEnvironment{tikzpicture}{\end{adjustbox}}
% légendes pgfplots lisibles par-dessus les courbes
\pgfplotsset{every axis/.append style={legend style={fill=white,fill opacity=0.92,text opacity=1,draw=black!15,inner sep=3pt}}}
% étiquettes d'arêtes (syntaxe edge["..."]) avec fond blanc
\tikzset{every edge quotes/.append style={fill=white,inner sep=1.2pt}}
% bulles de cartes mentales : texte à la ligne dans la bulle, bulle un peu plus grande
\tikzset{every concept/.append style={text width=2.0cm,align=center,minimum size=2.3cm,inner sep=2pt}}
%% FIGFIX-END
\begin{document}

\pagegarde{Lesson 26 — Grammar: Revise clauses}{Anglais}{Tle A}{Chapter 3 : Environment, Well-being and Health}{EN26}{2 h}{Cette leçon de grammaire révise les trois grands types de propositions anglaises (relative clauses, noun clauses, adverbial clauses) : leur formation, leurs emplois et leurs pièges. Elle consolide les acquis de Première et prépare les élèves de Terminale A à manier des phrases complexes à l'écrit comme à l'oral du Bac.
\vspace{4pt}
\begin{multicols}{2}\small
\begin{itemize}
\item À la fin de cette leçon, je sais distinguer une proposition principale d'une proposition subordonnée.
\item À la fin de cette leçon, je sais identifier les trois types de clauses : relatives, nominales (noun clauses) et circonstancielles (adverbial clauses).
\item À la fin de cette leçon, je sais employer correctement les pronoms relatifs who, which, that, whose, where, when.
\item À la fin de cette leçon, je sais distinguer une relative déterminative (sans virgules) d'une relative explicative (entre virgules).
\item À la fin de cette leçon, je sais introduire une noun clause avec that, if, whether ou un mot interrogatif.
\item À la fin de cette leçon, je sais construire des adverbial clauses de temps, cause, but, concession et condition avec la bonne concordance des temps.
\end{itemize}\end{multicols}}

\begin{sommaire}
\begin{multicols}{2}
\begin{enumerate}
\item Observation et notion de clause
\item Relative clauses : formation et pronoms relatifs
\item Relatives déterminatives et explicatives
\item Noun clauses : les subordonnées compléments
\item Adverbial clauses : temps, cause, but, concession, condition
\item Comparaison français/anglais et entraînement au Bac
\item Activité d'intégration
\item Méthodes et exercices résolus
\item Exercices
\item Corrigés des exercices
\item Fiche bilan
\end{enumerate}
\end{multicols}
\end{sommaire}

\begin{objectifs}
\begin{itemize}
\item À la fin de cette leçon, je sais distinguer une proposition principale d'une proposition subordonnée.
\item À la fin de cette leçon, je sais identifier les trois types de clauses : relatives, nominales (noun clauses) et circonstancielles (adverbial clauses).
\item À la fin de cette leçon, je sais employer correctement les pronoms relatifs who, which, that, whose, where, when.
\item À la fin de cette leçon, je sais distinguer une relative déterminative (sans virgules) d'une relative explicative (entre virgules).
\item À la fin de cette leçon, je sais introduire une noun clause avec that, if, whether ou un mot interrogatif.
\item À la fin de cette leçon, je sais construire des adverbial clauses de temps, cause, but, concession et condition avec la bonne concordance des temps.
\item À la fin de cette leçon, je sais éviter les erreurs fréquentes des francophones (double sujet, that après préposition, futur après when/if).
\item À la fin de cette leçon, je sais combiner des phrases simples en phrases complexes pour enrichir mes productions écrites.
\end{itemize}
\end{objectifs}

\begin{prerequis}
\begin{itemize}
\item La phrase simple anglaise : sujet + verbe + complément (classe de Seconde/Première)
\item Les pronoms relatifs de base who, which, that (Première)
\item Les connecteurs logiques : because, although, if, when (Première)
\item La concordance des temps dans les phrases conditionnelles (Première)
\item Le vocabulaire du chapitre 3 : environment, well-being and health (leçons précédentes)
\end{itemize}
\end{prerequis}

\newpage

\coursec{Observation et notion de clause}

\courssub{Situation d'accroche}

Pendant la campagne de sensibilisation sur l'environnement au lycée de Buea, une élève déclare devant toute la classe : \eng{The trees which are planted near the school will protect us from erosion.} « Les arbres qui sont plantés près de l'école nous protégeront de l'érosion. » Son camarade répond aussitôt : \eng{I believe that reforestation is the solution because our hills are losing their soil.} « Je crois que le reboisement est la solution parce que nos collines perdent leur sol. »

Sans le savoir, ces deux élèves viennent d'utiliser trois types de propositions (\cle{clauses}) étudiés depuis la classe de Première : une proposition relative (\eng{which are planted near the school}), une proposition complétive (\eng{that reforestation is the solution}) et une proposition circonstancielle de cause (\eng{because our hills are losing their soil}).

\textbf{Problématique :} qu'est-ce qu'une clause ? Comment distinguer une proposition principale d'une proposition subordonnée, et comment reconnaître les trois grandes familles de subordonnées en anglais ? Au Bac, la maîtrise des clauses est indispensable pour comprendre les textes de compréhension écrite et pour produire, à l'expression écrite, des phrases complexes correctes.

\courssub{Lecture et observation}

\begin{textebox}[An environmental warning — texte support]
The forest which surrounds our village is disappearing. Every year, farmers cut down more trees because they need land for their crops. Scientists say that deforestation is dangerous: when the trees disappear, the rain washes away the soil, and the rivers become muddy. If we do not act now, our children will suffer. The chief, who has lived here for sixty years, remembers that the hills were green when he was a boy. He believes that we can save our land if everyone plants at least one tree every year. Although the task is difficult, the villagers have decided to start a reforestation programme next month.
\end{textebox}

\textbf{Glossaire :} \emph{to surround} (v., entourer) ; \emph{to cut down} (v., abattre) ; \emph{crops} (n., cultures) ; \emph{to wash away} (v., emporter, éroder) ; \emph{muddy} (adj., boueux) ; \emph{chief} (n., chef du village) ; \emph{task} (n., tâche).

\textbf{Questions de compréhension :}
\begin{enumerate}
\item What is happening to the forest around the village?
\item Why do farmers cut down trees?
\item What does the chief remember about the hills?
\item What have the villagers decided to do?
\end{enumerate}

\begin{experience}[Activité d'observation — soulignons les verbes]
Relevez dans le texte tous les \cle{verbes conjugués} et soulignez-les. Puis entourez les mots qui les introduisent : \eng{which}, \eng{because}, \eng{that}, \eng{when}, \eng{if}, \eng{who}, \eng{although}. Que remarquez-vous ? Chaque verbe conjugué correspond à une \emph{clause}, et certains verbes sont précédés d'un \emph{mot de liaison} (\cle{linking word}) qui les rattache à une autre partie de la phrase.
\end{experience}

\courssub{Qu'est-ce qu'une clause ?}

\begin{definition}[Clause]
Une \cle{clause} (proposition) est un groupe de mots organisé autour d'un \cle{verbe conjugué} et contenant un \cle{sujet} (exprimé ou sous-entendu). Une phrase peut contenir une seule clause (phrase simple) ou plusieurs clauses (phrase complexe).
\begin{itemize}
\item \eng{The forest is disappearing.} = une seule clause (sujet \eng{the forest} + verbe \eng{is disappearing}).
\item \eng{The forest which surrounds our village is disappearing.} = deux clauses : \eng{the forest ... is disappearing} et \eng{which surrounds our village}.
\end{itemize}
\end{definition}

\begin{attention}[Clause n'est pas « phrase »]
Le mot anglais \eng{clause} est un faux ami partiel : il ne signifie pas « clause » au sens juridique, mais \emph{proposition}. Une \eng{sentence} (phrase) peut contenir plusieurs \eng{clauses}. De même, un groupe de mots sans verbe conjugué n'est pas une clause : \eng{in the morning}, \eng{a very tall tree} sont de simples groupes (\cle{phrases} au sens grammatical anglais), pas des clauses.
\end{attention}

\courssub{Main clause et subordinate clause}

\begin{definition}[Main clause / subordinate clause]
La \cle{main clause} (proposition principale) est autonome : elle peut exister seule et avoir un sens complet. La \cle{subordinate clause} (proposition subordonnée) est dépendante : introduite par un mot de liaison, elle ne peut pas exister seule et complète la principale.
\begin{itemize}
\item \eng{Our children will suffer.} (main clause — phrase complète, autonome)
\item \eng{If we do not act now, our children will suffer.} (\eng{if we do not act now} = subordinate clause ; seule, elle reste « en suspens »).
\end{itemize}
\end{definition}

\begin{propriete}[La subordonnée dépend toujours d'une principale]
Une \cle{subordinate clause} commence toujours par un \cle{mot de liaison} (\eng{who}, \eng{which}, \eng{that}, \eng{because}, \eng{if}, \eng{when}, \eng{although}...) et dépend d'une \cle{main clause}. Elle ne peut jamais constituer une phrase complète à elle seule. Écrire \eng{Because our hills are losing their soil.} comme une phrase isolée est une faute (\emph{fragment de phrase}) : il manque la principale.
\end{propriete}

\begin{exemplebox}[Observation guidée]
\eng{The man who lives here is a doctor.} = « L'homme qui vit ici est médecin. »
\begin{itemize}
\item Main clause : \eng{the man ... is a doctor} (autonome).
\item Subordinate clause : \eng{who lives here} (introduite par \eng{who}, ne peut pas exister seule).
\end{itemize}
\eng{I know that he is right.} = « Je sais qu'il a raison. »
\begin{itemize}
\item Main clause : \eng{I know}.
\item Subordinate clause : \eng{that he is right} (complète le verbe \eng{know}).
\end{itemize}
\end{exemplebox}

\courssub{Les trois types de subordonnées}

\begin{definition}[Les trois familles de subordonnées]
\begin{enumerate}
\item La \cle{relative clause} (proposition relative) : introduite par \eng{who}, \eng{which}, \eng{that}, \eng{whose}, \eng{where} ; elle décrit ou précise un nom. \eng{The chief who has lived here for sixty years remembers...}
\item La \cle{noun clause} (proposition complétive) : introduite par \eng{that}, \eng{whether}, \eng{if}, \eng{what}, \eng{how} ; elle fonctionne comme un nom (complément du verbe). \eng{Scientists say that deforestation is dangerous.}
\item La \cle{adverbial clause} (proposition circonstancielle) : introduite par \eng{because}, \eng{if}, \eng{when}, \eng{although}, \eng{so that} ; elle joue le rôle d'un adverbe (temps, cause, but, concession, condition). \eng{If we do not act now, our children will suffer.}
\end{enumerate}
\end{definition}

\begin{center}
\begin{tabularx}{\linewidth}{|X|X|X|}
\hline
\rowcolor{popblueL} Type de clause & Mots de liaison typiques & Exemple \\
\hline
Relative clause & who, which, that, whose, where & \eng{The trees which are planted near the school...} \\
\hline
Noun clause & that, whether, if, what, how & \eng{I believe that reforestation is the solution.} \\
\hline
Adverbial clause & because, if, when, although, so that & \eng{...because our hills are losing their soil.} \\
\hline
\end{tabularx}
\end{center}

\begin{popfigure}
\begin{tikzpicture}[level distance=1.6cm,
  level 1/.style={sibling distance=5.2cm},
  level 2/.style={sibling distance=1.9cm},
  every node/.style={draw, rounded corners, align=center, font=\small, inner sep=3pt}]
\node[fill=popblue!30, align=center]{CLAUSE\\ sujet + verbe conjugué}
  child { node[fill=popgreen!30, align=center]{MAIN CLAUSE\\ autonome} }
  child { node[fill=poporange!30, align=center]{SUBORDINATE CLAUSE\\ dépendante}
    child { node[fill=poporange!15, align=center]{relative\\ who, which} }
    child { node[fill=poporange!15, align=center]{noun\\ that, what} }
    child { node[fill=poporange!15, align=center]{adverbial\\ if, because} }
  };
\end{tikzpicture}
\legende{Les deux grandes catégories de clauses et les trois types de subordonnées.}
\end{popfigure}

\begin{popfigure}
\begin{tikzpicture}[mindmap, grow cyclic,
  every node/.style={concept, align=center, font=\small},
  root concept/.append style={fill=popblue!40, font=\small\bfseries},
  level 1/.append style={level distance=3.4cm, sibling angle=120},
  level 2/.append style={level distance=2.2cm, font=\scriptsize}]
\node {CLAUSES}
  child[concept color=popgreen!40] { node[align=center]{Relative\\ clauses}
    child { node {who / which} }
    child { node {that / whose} }
    child { node {where} }
  }
  child[concept color=poporange!40] { node[align=center]{Noun\\ clauses}
    child { node {that} }
    child { node {whether / if} }
    child { node {what / how} }
  }
  child[concept color=poppurple!35] { node[align=center]{Adverbial\\ clauses}
    child { node {because / if} }
    child { node {when / while} }
    child { node {although / so that} }
  };
\end{tikzpicture}
\legende{Carte mentale : les trois types de subordonnées et leurs mots de liaison typiques.}
\end{popfigure}

\courssub{Méthode d'identification et test d'autonomie}

\begin{methode}[Identifier les clauses d'une phrase en quatre étapes]
\begin{enumerate}
\item \textbf{Souligner} tous les verbes conjugués de la phrase : chacun correspond à une clause.
\item \textbf{Chercher}, devant chaque verbe, un mot de liaison (\eng{who}, \eng{which}, \eng{that}, \eng{because}, \eng{if}, \eng{when}, \eng{although}...) : s'il y en a un, la clause est subordonnée.
\item \textbf{Appliquer le test d'autonomie} : lire la clause seule. Si elle a un sens complet, c'est la main clause ; si elle reste « en suspens », c'est une subordinate clause.
\item \textbf{Classer} la subordonnée : décrit-elle un nom (relative) ? sert-elle de complément au verbe (noun clause) ? exprime-t-elle temps, cause, but, condition ou concession (adverbial clause) ?
\end{enumerate}
\end{methode}

\begin{exoresolu}[Analyse d'une phrase complexe]
Énoncé : identifiez les clauses de la phrase \eng{The chief believes that we can save our land if everyone plants a tree.} et précisez le type de chaque subordonnée.
\tcblower
Solution détaillée :
\begin{enumerate}
\item Verbes conjugués : \eng{believes}, \eng{can save}, \eng{plants} : trois clauses.
\item Mots de liaison : \eng{that} (devant \eng{we can save}) et \eng{if} (devant \eng{everyone plants}).
\item Test d'autonomie : \eng{the chief believes} peut exister seule $\rightarrow$ main clause. \eng{That we can save our land} et \eng{if everyone plants a tree} ne peuvent pas exister seules $\rightarrow$ subordinate clauses.
\item Classement : \eng{that we can save our land} est complément du verbe \eng{believes} $\rightarrow$ \textbf{noun clause}. \eng{If everyone plants a tree} exprime une condition $\rightarrow$ \textbf{adverbial clause} (de condition).
\end{enumerate}
Traduction : « Le chef croit que nous pouvons sauver notre terre si chacun plante un arbre. »
\end{exoresolu}

\begin{attention}[Pièges fréquents des francophones]
\begin{itemize}
\item Ne pas confondre \eng{that} relatif et \eng{that} complétif : dans \eng{the forest that surrounds us}, \eng{that} est sujet du verbe (relative) ; dans \eng{I know that he is right}, \eng{that} introduit simplement la complétive et ne joue aucun rôle dans sa clause.
\item En anglais, on ne met \textbf{jamais} de virgule avant \eng{that} dans une complétive : \eng{I know, that he is right} est fautif.
\item Une subordonnée seule n'est pas une phrase : évitez les fragments du type \eng{Because they need land.} sans principale.
\item Le mot de liaison est souvent obligatoire en anglais là où le français le répète ou le garde : \eng{I think (that) he is right} — \eng{that} peut être omis à l'oral, mais jamais le sujet qui suit.
\end{itemize}
\end{attention}

\begin{exercice}[À vous de jouer]
Pour chaque phrase, soulignez les verbes conjugués, encadrez les mots de liaison et indiquez le type de chaque subordonnée (relative, noun, adverbial) :
\begin{enumerate}
\item \eng{The river which crosses Douala is polluted.}
\item \eng{Doctors say that clean water is essential for health.}
\item \eng{When the rains come, the farmers plant their seeds.}
\item \eng{The nurse who works at the health centre is very kind.}
\item \eng{Although the forest is far, the students will visit it because they want to learn.}
\end{enumerate}
\end{exercice}

\begin{corrige}[Exercice « À vous de jouer »]
\begin{enumerate}
\item \eng{is polluted} (main clause : \eng{the river ... is polluted}) ; \eng{which crosses Douala} : relative clause (mot de liaison \eng{which}).
\item \eng{say} (main clause) ; \eng{that clean water is essential for health} : noun clause, complément de \eng{say}.
\item \eng{the farmers plant their seeds} (main clause) ; \eng{when the rains come} : adverbial clause de temps.
\item \eng{is very kind} (main clause) ; \eng{who works at the health centre} : relative clause.
\item \eng{the students will visit it} (main clause) ; \eng{although the forest is far} : adverbial clause de concession ; \eng{because they want to learn} : adverbial clause de cause.
\end{enumerate}
\end{corrige}

\begin{aretenir}[L'essentiel de la section]
\begin{itemize}
\item Une \cle{clause} = sujet + verbe conjugué ; une phrase peut en contenir plusieurs.
\item La \cle{main clause} est autonome ; la \cle{subordinate clause}, introduite par un mot de liaison, ne peut pas exister seule (test d'autonomie).
\item Trois types de subordonnées : \cle{relative} (\eng{who}, \eng{which}, \eng{that}), \cle{noun} (\eng{that}, \eng{whether}, \eng{what}), \cle{adverbial} (\eng{because}, \eng{if}, \eng{when}, \eng{although}).
\item Méthode : souligner les verbes, repérer les mots de liaison, tester l'autonomie, classer.
\end{itemize}
\end{aretenir}

\coursec{Relative clauses : formation et pronoms relatifs}

Dans la section précédente, nous avons vu qu'une \cle{clause} est un groupe de mots organisé autour d'un verbe conjugué, et qu'une phrase complexe combine une clause principale et une ou plusieurs clauses subordonnées. Nous allons maintenant étudier la première grande famille de subordonnées : les \cle{relative clauses} (propositions subordonnées relatives), qui servent à donner une information sur un nom. C'est l'une des structures les plus fréquentes de l'anglais écrit — et l'une des mieux notées au Bac lorsqu'elle est maîtrisée.

\courssub{Observation : à quoi sert une relative clause ?}

\begin{textebox}[Reading text — “A Clinic in the Hills” (original text)]
Last month, I visited a health centre which stands on a hill near Bamenda. The nurse who welcomed us has worked there for fifteen years. She showed us the room where vaccines are kept and introduced us to a young doctor whose dedication impressed everyone. The medicines that the centre receives come from Douala, but the road, which is very steep, makes transport difficult. The day when the new generator arrived was a day of celebration: it was the reason why the clinic could finally store vaccines safely. The villagers, most of whom walk several kilometres to reach the centre, are grateful to the people whose generosity made this project possible.
\end{textebox}

\textbf{Glossaire :} \eng{health centre} (n. m.) : centre de santé ; \eng{steep} (adj.) : escarpé ; \eng{generator} (n. m.) : groupe électrogène ; \eng{to store} (v.) : conserver ; \eng{grateful} (adj.) : reconnaissant ; \eng{dedication} (n. f.) : dévouement.

\textbf{Questions d'observation :}
\begin{enumerate}
\item Soulignez dans le texte tous les mots : \eng{who, which, whose, where, when, why, that}. Quel nom chacun d'eux suit-il directement ?
\item Dans \eng{“The nurse who welcomed us”}, le mot \eng{who} remplace-t-il un sujet ou un complément ?
\item Traduisez : \eng{“the room where vaccines are kept”} et \eng{“a young doctor whose dedication impressed everyone”}.
\end{enumerate}

\begin{definition}[Relative clause et antécédent]
Une \cle{relative clause} est une proposition subordonnée qui apporte une information sur un nom appelé \cle{antécédent}, placé \textbf{juste avant} elle. Elle est introduite par un \cle{pronom relatif} (\eng{who, which, that, whom, whose}) ou un \cle{adverbe relatif} (\eng{where, when, why}). Elle correspond en français aux propositions en \emph{qui, que, dont, où, lequel, auquel, duquel}.
\end{definition}

\begin{exemplebox}[Observer avant de conclure]
\eng{The nurse who treated me was kind.} « L'infirmière qui m'a soignée était gentille. »

\eng{This is the village where I was born.} « Voici le village où je suis né. »

\eng{The girl whose father is a doctor won the prize.} « La fille dont le père est médecin a gagné le prix. »
\end{exemplebox}

On observe que le pronom relatif anglais se choisit selon \textbf{la nature de l'antécédent} (personne, chose, lieu, moment) et selon \textbf{sa fonction} dans la relative (sujet, complément, possession).

\courssub{Les pronoms relatifs : who, which, that, whom, whose}

\begin{propriete}[Choix du pronom relatif]
\begin{itemize}
\item \cle{who} : antécédent = \textbf{personne}, fonction = \textbf{sujet}. \eng{The doctor who lives here is kind.} « Le médecin qui habite ici est gentil. »
\item \cle{which} : antécédent = \textbf{chose ou animal}, fonction = sujet ou complément. \eng{The vaccine which protects children is free.} « Le vaccin qui protège les enfants est gratuit. »
\item \cle{that} : peut remplacer \eng{who} ou \eng{which} en \textbf{registre courant} (jamais après une préposition ni après une virgule). \eng{The nurse that I saw was busy.} « L'infirmière que j'ai vue était occupée. »
\item \cle{whom} : antécédent = personne, fonction = \textbf{complément} (registre soutenu ; souvent remplacé par \eng{who} ou omis à l'oral). \eng{The man whom I met is a surgeon.} « L'homme que j'ai rencontré est chirurgien. »
\item \cle{whose} : \textbf{possession}, pour les personnes \textbf{et} les choses (= \emph{dont, duquel}). \eng{The patient whose temperature is high must rest.} « Le patient dont la température est élevée doit se reposer. »
\end{itemize}
\end{propriete}

\begin{propriete}[Les adverbes relatifs : where, when, why]
\begin{itemize}
\item \cle{where} : pour les \textbf{lieux} (= \emph{où, dans lequel}). \eng{The hospital where she works is modern.} « L'hôpital où elle travaille est moderne. »
\item \cle{when} : pour les \textbf{moments} (= \emph{où, auquel}). \eng{I remember the day when the clinic opened.} « Je me souviens du jour où le centre de santé a ouvert. »
\item \cle{why} : uniquement après \eng{the reason}. \eng{That is the reason why he stopped smoking.} « C'est la raison pour laquelle il a arrêté de fumer. »
\end{itemize}
\end{propriete}

\begin{center}
\begin{tabularx}{\linewidth}{|X|X|X|X|}
\hline
\rowcolor{popblueL} Antécédent & Sujet & Complément & Possession \\
\hline
Personne & who (that) & whom / who (that) & whose \\
\hline
Chose / animal & which (that) & which (that) & whose / of which \\
\hline
Lieu & \multicolumn{3}{l|}{where (= in/at which)} \\
\hline
Moment & \multicolumn{3}{l|}{when (= on/at which)} \\
\hline
Raison & \multicolumn{3}{l|}{why (après \eng{the reason})} \\
\hline
\end{tabularx}
\end{center}

\begin{popfigure}
\begin{tikzpicture}[
  box/.style={draw=popblue, fill=popblueL, rounded corners=2pt, align=center, font=\small, inner sep=4pt},
  head/.style={draw=popdark, fill=popblue, text=white, rounded corners=2pt, align=center, font=\small\bfseries, inner sep=4pt}
]
\node[head] (a) at (0,0) {Antécédent};
\node[head] (s) at (3.4,0) {Sujet};
\node[head] (c) at (6.8,0) {Complément};
\node[head] (p) at (10.2,0) {Possession};
\node[box] (a1) at (0,-1.1) {personne};
\node[box, align=center] (s1) at (3.4,-1.1){\textbf{who}\\(that)};
\node[box, align=center] (c1) at (6.8,-1.1){\textbf{whom} / who\\(that)};
\node[box] (p1) at (10.2,-1.1) {\textbf{whose}};
\node[box] (a2) at (0,-2.4) {chose / animal};
\node[box, align=center] (s2) at (3.4,-2.4){\textbf{which}\\(that)};
\node[box, align=center] (c2) at (6.8,-2.4){\textbf{which}\\(that)};
\node[box, align=center] (p2) at (10.2,-2.4){\textbf{whose}\\of which};
\node[box, fill=popgreenL, draw=popgreen] (l) at (1.7,-3.7) {lieu : \textbf{where}};
\node[box, fill=poporangeL, draw=poporange] (t) at (5.1,-3.7) {temps : \textbf{when}};
\node[box, fill=poppinkL, draw=poppink] (r) at (8.5,-3.7) {the reason \textbf{why}};
\end{tikzpicture}
\legende{Tableau des pronoms relatifs anglais selon l'antécédent et la fonction.}
\end{popfigure}

\courssub{L'omission du pronom relatif complément}

\begin{propriete}[Omission possible : relatif complément uniquement]
Quand le pronom relatif est \textbf{complément} (et non sujet), on peut l'omettre en style courant :

\eng{The book (that) I bought is interesting.} « Le livre que j'ai acheté est intéressant. »

\eng{The nurse (whom) we thanked smiled.} « L'infirmière que nous avons remerciée a souri. »

En revanche, quand le relatif est \textbf{sujet}, il est \textbf{obligatoire} : \eng{The boy who came first won the race.} « Le garçon qui est arrivé premier a gagné la course. » — on ne peut jamais dire \eng{The boy came first won...}.
\end{propriete}

\textbf{Astuce de reconnaissance :} si deux sujets se suivent (\eng{the book / I}), le relatif est complément et peut tomber ; si le relatif est immédiatement suivi d'un verbe (\eng{who came}), il est sujet et doit rester.

\courssub{Préposition + relatif : registre soutenu et registre courant}

\begin{propriete}[Deux constructions possibles]
\begin{itemize}
\item \textbf{Registre soutenu} : la préposition précède le relatif, et l'on utilise obligatoirement \eng{which} (choses) ou \eng{whom} (personnes) : \eng{the house in which I live} ; \eng{the colleague with whom I work}.
\item \textbf{Registre courant} : la préposition reste en fin de relative, et le relatif (\eng{that}) est souvent omis : \eng{the house (that) I live in} ; \eng{the colleague (that) I work with}.
\end{itemize}
\end{propriete}

\begin{exemplebox}[Les deux registres côte à côte]
\eng{The hospital in which my mother works is far.} = \eng{The hospital (that) my mother works in is far.} « L'hôpital dans lequel ma mère travaille est loin. »

\eng{The doctor to whom I spoke was reassuring.} = \eng{The doctor (that) I spoke to was reassuring.} « Le médecin à qui j'ai parlé était rassurant. »
\end{exemplebox}

\begin{attention}[Après une préposition : jamais “that” !]
Après une préposition placée devant le relatif, on utilise \textbf{uniquement} \eng{which} ou \eng{whom} :

✗ \eng{the pen with that I write} → ✓ \eng{the pen with which I write} « le stylo avec lequel j'écris »

✗ \eng{the friend to that I talked} → ✓ \eng{the friend to whom I talked} « l'ami à qui j'ai parlé »
\end{attention}

\begin{attention}[Ne jamais répéter le sujet ou le complément]
Le pronom relatif \textbf{remplace} déjà le nom : il ne faut jamais ajouter un pronom personnel après lui.

✗ \eng{The man who he came yesterday is a doctor.} → ✓ \eng{The man who came yesterday is a doctor.} « L'homme qui est venu hier est médecin. »

✗ \eng{The film which I saw it was boring.} → ✓ \eng{The film which I saw was boring.} « Le film que j'ai vu était ennuyeux. »

C'est l'erreur la plus fréquente des francophones, car en français on dit « l'homme qui, lui, est venu… » dans la langue familière.
\end{attention}

\begin{center}
\begin{tabularx}{\linewidth}{|X|X|X|}
\hline
\rowcolor{popblueL} Français & English & Remarque \\
\hline
l'infirmière \emph{qui} soigne & the nurse \textbf{who} treats & personne sujet \\
\hline
le vaccin \emph{qui} protège & the vaccine \textbf{which} protects & chose sujet \\
\hline
le livre \emph{que} j'ai lu & the book (\textbf{that}) I read & complément : omission possible \\
\hline
la fille \emph{dont} le père & the girl \textbf{whose} father & possession, personnes et choses \\
\hline
le village \emph{où} je suis né & the village \textbf{where} I was born & lieu \\
\hline
le jour \emph{où} il est parti & the day \textbf{when} he left & moment \\
\hline
la raison \emph{pour laquelle} & the reason \textbf{why} & après \eng{the reason} seulement \\
\hline
l'ami \emph{avec qui} je sors & the friend \textbf{with whom} I go out & soutenu ; courant : \eng{the friend I go out with} \\
\hline
\end{tabularx}
\end{center}

\begin{exoresolu}[Reliez les deux phrases par un pronom relatif]
Énoncé — Combine each pair of sentences using a relative pronoun (who, which, whose, where, when, that) :
\begin{enumerate}
\item The woman is a midwife. She helped my sister.
\item I bought a book. It is about healthy food.
\item That is the school. I studied there.
\item The boy won the prize. His mother teaches biology.
\item I will never forget the year. We moved to Yaoundé then.
\end{enumerate}
\tcblower
Solution détaillée :
\begin{enumerate}
\item \eng{The woman who helped my sister is a midwife.} — antécédent \eng{the woman} (personne), relatif \textbf{sujet} du verbe \eng{helped} → \eng{who} obligatoire.
\item \eng{The book (that/which) I bought is about healthy food.} — antécédent chose, relatif \textbf{complément} de \eng{bought} → \eng{which/that}, omission possible.
\item \eng{That is the school where I studied.} — antécédent lieu → \eng{where} (= \eng{in which}).
\item \eng{The boy whose mother teaches biology won the prize.} — possession (« sa mère ») → \eng{whose}.
\item \eng{I will never forget the year when we moved to Yaoundé.} — antécédent moment → \eng{when} (= \eng{in which}).
\end{enumerate}
\end{exoresolu}

\begin{exercice}[À vous de jouer]
\begin{enumerate}
\item Complétez avec \eng{who, which, whose, where, when} ou \eng{why} : \eng{The farmer \ldots\ fields were flooded lost his crops.} / \eng{The market \ldots\ my mother sells fruit is near the stadium.} / \eng{Give me one reason \ldots\ I should believe you.}
\item Mettez au style soutenu : \eng{The girl (that) I go to school with is my cousin.}
\item Corrigez : \eng{The students who they passed the exam celebrated.}
\end{enumerate}
\end{exercice}

\begin{corrige}[Exercice « À vous de jouer »]
1. \eng{whose} (possession) ; \eng{where} (lieu) ; \eng{why} (après \eng{reason}). 2. \eng{The girl with whom I go to school is my cousin.} (préposition \eng{with} devant + \eng{whom}). 3. \eng{The students who passed the exam celebrated.} — supprimer \eng{they}, car \eng{who} est déjà le sujet.
\end{corrige}

\begin{savaistu}[“That”, le champion de l'oral]
En anglais parlé, \eng{that} est de loin le pronom relatif le plus fréquent : les Anglophones disent naturellement \eng{“the man that lives next door”} ou \eng{“the book that I read”}. Mais attention : au Bac écrit, savoir varier \eng{who}, \eng{which}, \eng{whose}, \eng{where} et employer correctement \eng{whom} dans une lettre formelle montre une réelle maîtrise de la langue — et c'est exactement ce que les correcteurs récompensent. Réservez donc \eng{that} aux dialogues, et brillez à l'écrit avec toute la gamme des relatifs !
\end{savaistu}

\begin{aretenir}[L'essentiel de la section]
\begin{itemize}
\item La relative clause suit immédiatement son antécédent et le décrit.
\item Personnes : \eng{who} (sujet), \eng{whom} (complément soutenu) ; choses : \eng{which} ; \eng{that} remplace \eng{who/which} en style courant, jamais après préposition ni virgule.
\item \eng{whose} = possession (personnes et choses) ; \eng{where} = lieu ; \eng{when} = moment ; \eng{why} après \eng{the reason}.
\item Relatif complément → omission possible ; relatif sujet → jamais d'omission, jamais de pronom redoublé.
\end{itemize}
\end{aretenir}

\coursec{Relatives déterminatives et explicatives}

Dans la section précédente, nous avons appris à former une relative clause et à choisir le bon pronom relatif (\eng{who}, \eng{which}, \eng{that}, \eng{whose}, \eng{whom}). Mais il reste une question essentielle, que l'examinateur du Bac surveille de très près : faut-il mettre des \cle{virgules} autour de la relative ou non ? La réponse change le sens de la phrase. C'est toute la différence entre la relative \cle{déterminative} (\eng{defining relative clause}) et la relative \cle{explicative} (\eng{non-defining relative clause}).

\courssub{Observation : deux phrases, deux sens}

\begin{textebox}[Observing commas in a school report]
Read these two sentences written by a headmistress in Buea about her students.

\eng{“The students who cheated in the exam were punished.”}

\eng{“The students, who cheated in the exam, were punished.”}

Both sentences contain the same words. Only the commas change. Yet the meaning is completely different. In the first sentence, only \emph{some} students cheated — those who cheated were punished, the others were not. In the second sentence, \emph{all} the students cheated, and all of them were punished. The commas decide how many students are concerned.
\end{textebox}

\textbf{Glossaire}

\begin{center}
\begin{tabularx}{\linewidth}{|X|X|X|}\hline
\rowcolor{popblueL} Mot anglais & Nature & Traduction \\ \hline
headmistress & nom & directrice (d'école) \\ \hline
to cheat & verbe & tricher \\ \hline
to punish & verbe & punir \\ \hline
concerned & adjectif & concerné \\ \hline
\end{tabularx}
\end{center}

\textbf{Questions de compréhension}

\begin{enumerate}
\item In the first sentence, were all the students punished?
\item In the second sentence, how many students cheated?
\item What single element changes the meaning between the two sentences?
\end{enumerate}

\textbf{Observation grammaticale.} Dans la première phrase, la relative \eng{who cheated in the exam} est \emph{indispensable} : sans elle, on ne sait pas de quels élèves il s'agit. Elle \emph{détermine}, elle \emph{identifie} le nom. Dans la seconde, la même relative n'apporte qu'une information \emph{supplémentaire}, entre parenthèses en quelque sorte : on sait déjà de qui l'on parle (\eng{the students}, tous les élèves). C'est cette différence de fonction qui commande la ponctuation.

\courssub{La relative déterminative (defining relative clause)}

\begin{definition}[Relative déterminative]
Une \cle{relative déterminative} (\eng{defining} ou \eng{restrictive relative clause}) identifie précisément la personne ou la chose dont on parle. Elle est \textbf{indispensable au sens} de la phrase : si on la supprime, la phrase devient vague ou change de sens. Elle s'écrit \textbf{sans virgules}.
\end{definition}

\begin{exemplebox}[Relatives déterminatives]
\eng{“Students who work hard succeed.”} « Les élèves qui travaillent dur réussissent. » (seulement ceux qui travaillent dur, pas les autres)

\eng{“The trees which were planted last year have grown.”} « Les arbres qui ont été plantés l'an dernier ont poussé. » (seulement ceux-là, pas tous les arbres)

\eng{“The girl that I met at the market is my cousin.”} « La fille que j'ai rencontrée au marché est ma cousine. »

\eng{“This is the school where my father studied.”} « Voici l'école où mon père a étudié. »
\end{exemplebox}

Dans une relative déterminative, on peut :
\begin{itemize}
\item utiliser \eng{who}, \eng{which} ou \eng{that} (\eng{that} est même préféré dans la langue courante) ;
\item \textbf{omettre} le pronom relatif quand il est complément d'objet : \eng{“The book (that) I bought is interesting.”} « Le livre que j'ai acheté est intéressant. »
\end{itemize}

\courssub{La relative explicative (non-defining relative clause)}

\begin{definition}[Relative explicative]
Une \cle{relative explicative} (\eng{non-defining} ou \eng{non-restrictive relative clause}) donne une information \textbf{supplémentaire}, non essentielle, sur une personne ou une chose déjà identifiée (nom propre, possesseur, nom unique). Elle est placée \textbf{entre deux virgules} (ou entre une virgule et un point en fin de phrase).
\end{definition}

\begin{exemplebox}[Relatives explicatives]
\eng{“My father, who is a farmer, works hard.”} « Mon père, qui est agriculteur, travaille dur. » (je n'ai qu'un père : la relative ajoute un détail)

\eng{“The trees, which were planted last year, have grown.”} « Les arbres, qui d'ailleurs ont été plantés l'an dernier, ont poussé. » (tous les arbres, sans exception)

\eng{“Mount Cameroon, which is 4 040 metres high, attracts many tourists.”} « Le mont Cameroun, qui culmine à 4 040 mètres, attire de nombreux touristes. »

\eng{“Our English teacher, whom we all admire, won a national prize.”} « Notre professeur d'anglais, que nous admirons tous, a remporté un prix national. »
\end{exemplebox}

\begin{propriete}[Les trois interdits de la relative explicative]
Dans une relative \textbf{explicative} (entre virgules) :
\begin{enumerate}
\item \eng{that} est \textbf{interdit} : on utilise obligatoirement \eng{who} (personnes) ou \eng{which} (choses). On écrit \eng{“My uncle, who lives in Bamenda, …”} et jamais \eng{“My uncle, that lives in Bamenda, …”}.
\item L'\textbf{omission} du pronom relatif est \textbf{interdite}, même quand il est complément : \eng{“This novel, which I finished yesterday, is excellent.”} (on ne peut pas supprimer \eng{which}).
\item Les \textbf{deux virgules} sont \textbf{obligatoires} : une avant et une après la relative (sauf si la phrase se termine, auquel cas un point remplace la seconde virgule).
\end{enumerate}
\end{propriete}

\begin{propriete}[Tableau récapitulatif : déterminative ou explicative ?]
\begin{center}
\begin{tabularx}{\linewidth}{|X|X|X|}\hline
\rowcolor{popblueL} Critère & Déterminative (defining) & Explicative (non-defining) \\ \hline
Fonction & identifie, indispensable & ajoute un détail, supprimable \\ \hline
Virgules & aucune & deux virgules obligatoires \\ \hline
\eng{that} & possible (et fréquent) & interdit \\ \hline
Omission du relatif & possible (relatif objet) & interdite \\ \hline
Nom visé & nom commun indéterminé & nom propre, possesseur, réalité unique \\ \hline
Sens si on la supprime & phrase incomplète ou ambiguë & phrase toujours correcte \\ \hline
\end{tabularx}
\end{center}
\end{propriete}

\courssub{L'effet de sens : la virgule qui change la famille}

C'est le point le plus subtil — et le plus souvent interrogé au Bac. La présence ou l'absence des virgules modifie le \emph{nombre} de personnes ou de choses concernées.

\begin{exemplebox}[Un seul frère ou plusieurs ?]
\eng{“My brother who lives in Douala is a teacher.”} « Mon frère qui vit à Douala est enseignant. » — Sens : j'ai \textbf{plusieurs frères}, et je précise duquel je parle (celui de Douala, pas celui de Yaoundé).

\eng{“My brother, who lives in Douala, is a teacher.”} « Mon frère, qui vit à Douala, est enseignant. » — Sens : je n'ai qu'\textbf{un seul frère} ; le fait qu'il vive à Douala est une simple information en plus.
\end{exemplebox}

\begin{popfigure}
\begin{center}
\begin{tikzpicture}[font=\small]
\node[draw=popblue, fill=popblueL, rounded corners, align=center, text width=6.4cm, inner sep=6pt] (a) at (0,0) {\textbf{My brother who lives}\\ \textbf{in Douala is a teacher.}\\[2pt] \emph{pas de virgules}};
\node[below=2pt of a, align=center, text width=6.4cm, text=popdark] {« Mon frère qui vit à Douala... »\\ j'ai \textbf{plusieurs} frères};
\node[draw=poporange, fill=poporangeL, rounded corners, align=center, text width=6.4cm, inner sep=6pt] (b) at (8.6,0) {\textbf{My brother, who lives}\\ \textbf{in Douala, is a teacher.}\\[2pt] \emph{deux virgules}};
\node[below=2pt of b, align=center, text width=6.4cm, text=popdark] {« Mon frère, qui vit à Douala,... »\\ j'ai \textbf{un seul} frère};
\node[draw=popgreen, fill=popgreenL, rounded corners, align=center, text width=4.6cm, inner sep=5pt] at (4.3,-3.6) {Mêmes mots,\\ sens différent !};
\draw[-{Stealth}, thick, popblue] (a.south) -- ++(0,-0.4) -| (3.2,-3.0);
\draw[-{Stealth}, thick, poporange] (b.south) -- ++(0,-0.4) -| (5.4,-3.0);
\end{tikzpicture}
\end{center}
\legende{La virgule comme interrupteur de sens : déterminative (restriction) contre explicative (simple ajout).}
\end{popfigure}

\courssub{Comparaison avec le français}

Le français connaît aussi la distinction, mais de façon moins systématique et surtout moins sanctionnée à l'écrit. En français, « les élèves qui ont triché ont été punis » et « les élèves, qui ont triché, ont été punis » existent tous les deux, mais l'usage des virgules y est plus souple. En anglais, la règle est \textbf{stricte} : la ponctuation fait partie de la grammaire, pas seulement du style.

\begin{center}
\begin{tabularx}{\linewidth}{|X|X|X|}\hline
\rowcolor{popblueL} Français & English & Remarque \\ \hline
Les élèves qui travaillent réussissent. & Students who work hard succeed. & déterminative : pas de virgules dans les deux langues \\ \hline
Mon père, qui est agriculteur, travaille dur. & My father, who is a farmer, works hard. & explicative : virgules dans les deux langues \\ \hline
Le livre que j'ai lu est bon. & The book (that) I read is good. & en anglais on peut omettre le relatif objet ; en français « que » est obligatoire \\ \hline
Mon oncle, qui vit à Bamenda, ... & My uncle, who lives in Bamenda, ... & jamais \eng{that} après une virgule \\ \hline
\end{tabularx}
\end{center}

\begin{attention}[Erreurs fréquentes des francophones]
\begin{itemize}
\item \textbf{Oublier la deuxième virgule.} Beaucoup d'élèves écrivent \eng{“My father, who is a farmer works hard.”} Faux ! La virgule se met \textbf{avant ET après} la relative explicative : \eng{“My father, who is a farmer, works hard.”} Pensez à des parenthèses : on n'ouvre jamais une parenthèse sans la fermer.
\item \textbf{Utiliser \eng{that} entre virgules.} \eng{“Yaoundé, that is the capital, …”} est incorrect : après une virgule, c'est toujours \eng{who} ou \eng{which}.
\item \textbf{Omettre le relatif dans une explicative.} \eng{“This novel, I finished yesterday, is excellent.”} est impossible : il faut \eng{“This novel, which I finished yesterday, is excellent.”}
\item \textbf{Mettre des virgules à une déterminative.} \eng{“Students, who work hard, succeed.”} signifie que \emph{tous} les élèves travaillent dur — ce n'est probablement pas ce que vous vouliez dire !
\end{itemize}
\end{attention}

\begin{methode}[Comment choisir : le test de suppression]
Pour décider si la relative est déterminative ou explicative, posez-vous une seule question :
\begin{enumerate}
\item \textbf{Supprimez mentalement la relative.} La phrase garde-t-elle un sens complet et clair ?
\item Si \textbf{oui}, l'information est supplémentaire : c'est une relative \textbf{explicative} → mettez \textbf{deux virgules}, utilisez \eng{who}/\eng{which}, jamais \eng{that}.
\item Si \textbf{non} (la phrase devient vague : on ne sait plus de qui ou de quoi on parle), c'est une relative \textbf{déterminative} → \textbf{pas de virgules}, \eng{that} possible, omission possible si le relatif est objet.
\item Vérifiez le nom visé : un \textbf{nom propre} (\eng{Mount Cameroon}), un \textbf{possesseur} (\eng{my mother}) ou une \textbf{réalité unique} (\eng{the sun}) appelle presque toujours une explicative.
\end{enumerate}
\end{methode}

\begin{exoresolu}[Déterminative ou explicative ? Ponctuez et justifiez]
Énoncé. Ponctuez les phrases suivantes (ajoutez des virgules si nécessaire) et remplacez le relatif si besoin.
\begin{enumerate}
\item \eng{My grandmother that lives in the village grows vegetables.}
\item \eng{The candidates who arrived late were not allowed in.}
\item \eng{The Wouri River that flows through Douala is very wide.}
\end{enumerate}
\tcblower
Solution détaillée.
\begin{enumerate}
\item \eng{“My grandmother, who lives in the village, grows vegetables.”} — On n'a qu'une seule grand-mère (dans ce contexte) : la relative est une simple précision → explicative, deux virgules, et \eng{that} devient obligatoirement \eng{who}. « Ma grand-mère, qui vit au village, cultive des légumes. »
\item \eng{“The candidates who arrived late were not allowed in.”} — Sans la relative, « les candidats n'ont pas été admis » ne dit plus lesquels : la relative identifie → déterminative, aucune virgule. « Les candidats arrivés en retard n'ont pas été autorisés à entrer. »
\item \eng{“The Wouri River, which flows through Douala, is very wide.”} — Nom propre, réalité unique → explicative : virgules et \eng{which} à la place de \eng{that}. « Le fleuve Wouri, qui traverse Douala, est très large. »
\end{enumerate}
\end{exoresolu}

\begin{exercice}[À vous de jouer]
\begin{enumerate}
\item Ponctuez : \eng{“Our headmaster who comes from Bamenda speaks three languages.”}
\item Choisissez le bon relatif : \eng{“Limbe, (that / which) is a coastal town, has a beautiful botanic garden.”}
\item Traduisez en anglais : « Mon frère, qui habite à Yaoundé, est médecin. » (un seul frère) puis « Mon frère qui habite à Yaoundé est médecin. » (plusieurs frères).
\item Supprimez le relatif si possible : \eng{“The mangoes (that) she sold were sweet.”} / \eng{“These mangoes, which she sold cheaply, were sweet.”}
\end{enumerate}
\end{exercice}

\begin{corrige}[Exercice « À vous de jouer »]
\begin{enumerate}
\item \eng{“Our headmaster, who comes from Bamenda, speaks three languages.”} — un seul proviseur : explicative, deux virgules.
\item \eng{“Limbe, which is a coastal town, has a beautiful botanic garden.”} — \eng{that} est interdit entre virgules.
\item \eng{“My brother, who lives in Yaoundé, is a doctor.”} (un seul frère) ; \eng{“My brother who lives in Yaoundé is a doctor.”} (plusieurs frères).
\item \eng{“The mangoes she sold were sweet.”} (omission possible : déterminative, relatif objet) ; dans la seconde phrase, \eng{which} est \textbf{obligatoire} : explicative.
\end{enumerate}
\end{corrige}

\begin{savaistu}[La virgule, arme du rédacteur au Bac]
Dans les épreuves d'expression écrite, les correcteurs valorisent les candidats qui enrichissent leurs descriptions avec des relatives explicatives : \eng{“The Mfoundi, which crosses Yaoundé, is often polluted.”} ou \eng{“My grandmother, who is seventy, still farms her cocoa field.”} Une seule relative explicative bien ponctuée montre que vous maîtrisez une nuance que peu de candidats contrôlent. C'est un détail qui peut transformer une copie moyenne en excellente copie.
\end{savaistu}

\begin{aretenir}[L'essentiel de la section]
\begin{itemize}
\item Relative \textbf{déterminative} : indispensable, \textbf{sans virgules}, \eng{that} possible, omission possible (relatif objet).
\item Relative \textbf{explicative} : information en plus, \textbf{entre deux virgules}, \eng{that} interdit, omission interdite.
\item Test de suppression : si la phrase survit sans la relative → virgules.
\item La virgule change le sens : \eng{“My brother who lives in Douala”} (plusieurs frères) contre \eng{“My brother, who lives in Douala,”} (un seul).
\end{itemize}
\end{aretenir}

\coursec{Noun clauses : les subordonnées compléments}

Dans la section précédente, nous avons vu comment les \emph{relative clauses} (\eng{the man who..., the book which...}) se greffent sur un nom pour le préciser. Mais une subordonnée peut faire bien plus : elle peut, à elle seule, \cle{tenir la place d'un nom entier} dans la phrase. C'est exactement ce que fait la \emph{noun clause} : là où le français dit « Je sais \textbf{la vérité} », l'anglais peut dire \eng{I know \textbf{the truth}}… ou remplacer ce nom par toute une proposition : \eng{I know \textbf{that he is right}}. Cette section vous apprend à reconnaître, construire et manipuler ces subordonnées compléments, indispensables au Bac (compréhension de texte, style indirect, expression écrite).

\courssub{Observation : quand une phrase entière joue le rôle d'un nom}

\begin{textebox}[A radio interview about health in Yaoundé]
\textbf{Journalist:} Doctor, many people believe \textbf{that city air is more dangerous than before}. Do you agree?\\
\textbf{Dr. Mbarga:} I think \textbf{that pollution is a serious problem}, especially near busy markets. But I don't know \textbf{if people realise how much harm it does}. What worries me is \textbf{that most patients come too late}.\\
\textbf{Journalist:} Can you tell us \textbf{what the government should do}?\\
\textbf{Dr. Mbarga:} \textbf{That the authorities act quickly} is essential. I also hope \textbf{that schools will teach children how to protect their health}. Whether we like it or not, our lifestyle is changing, and I believe \textbf{that prevention is better than cure}.
\end{textebox}

\textbf{Glossaire :} \emph{harm} (n.m.) : mal, dommage ; \emph{to realise} (v.) : se rendre compte ; \emph{essential} (adj.) : essentiel ; \emph{prevention} (n.f.) : prévention ; \emph{cure} (n.m.) : remède.

\textbf{Observons.} Dans ce dialogue, chaque segment en gras pourrait être remplacé par un simple nom :

\begin{center}
\begin{tabularx}{\linewidth}{|X|X|}\hline
\rowcolor{popblueL} Phrase complète & Avec un nom simple \\ \hline
I believe \textbf{that pollution is a serious problem}. & I believe \textbf{it}. \\ \hline
I don't know \textbf{if people realise it}. & I don't know \textbf{the answer}. \\ \hline
\textbf{That the authorities act quickly} is essential. & \textbf{Quick action} is essential. \\ \hline
\end{tabularx}
\end{center}

\begin{definition}[Noun clause (subordonnée complément)]
Une \cle{noun clause} est une proposition subordonnée qui joue le rôle d'un \textbf{nom} (ou d'un groupe nominal) dans la phrase principale. Elle peut être :
\begin{itemize}
\item \textbf{complément d'objet} du verbe principal : \eng{I know that he is right.} (« Je sais qu'il a raison. ») — la subordonnée répond à la question « Je sais \emph{quoi} ? » ;
\item \textbf{sujet} du verbe principal : \eng{What he said is true.} (« Ce qu'il a dit est vrai. ») — la subordonnée répond à « \emph{Qu'est-ce qui} est vrai ? » ;
\item \textbf{attribut} après \emph{be} : \eng{The problem is that nobody listens.} (« Le problème, c'est que personne n'écoute. »).
\end{itemize}
\end{definition}

\courssub{Les trois portes d'entrée : that, if/whether, wh- words}

Une noun clause est toujours introduite par l'un de ces trois types de mots :

\begin{popfigure}
\begin{tikzpicture}[
  grow=down,
  level distance=1.6cm,
  every node/.style={align=center, font=\small},
  root/.style={rectangle, rounded corners, draw=popdark, fill=popblue, text=white, font=\bfseries\small, inner sep=6pt},
  branche/.style={rectangle, rounded corners, draw=popdark, fill=poporangeL, inner sep=5pt},
  ex/.style={rectangle, rounded corners, draw=popline, fill=popcream, inner sep=5pt, font=\itshape\small},
  edge from parent/.style={draw=popmuted, thick}
]
\node[root] {NOUN CLAUSE}
  child { node[branche, align=center]{\textbf{that}\\ (déclaration)}
    child { node[ex, align=center]{I believe that\\ pollution is dangerous.} } }
  child { node[branche, align=center]{\textbf{if / whether}\\ (question fermée)}
    child { node[ex, align=center]{He asked if\\ I was ready.} } }
  child { node[branche, align=center]{\textbf{wh- words}\\ (what, where, why...)}
    child { node[ex, align=center]{I don't know\\ where he lives.} } };
\end{tikzpicture}
\legende{Les trois types d'introducteurs d'une noun clause}
\end{popfigure}

\courssub{Les noun clauses introduites par « that »}

\begin{propriete}[That après les verbes d'opinion et de déclaration]
Après les verbes qui expriment une \textbf{opinion}, une \textbf{croyance} ou une \textbf{déclaration} — \eng{think, believe, say, know, hope, feel, suppose, admit, promise} — la subordonnée complément est introduite par \cle{that} :
\begin{center}
\textbf{sujet + verbe + (that) + sujet + verbe}
\end{center}
\eng{I believe that pollution is dangerous.} « Je crois que la pollution est dangereuse. »

\textbf{Omission de that :} à l'oral et dans le style courant, \eng{that} est très souvent omis, exactement comme « que » ne peut jamais l'être en français :
\eng{I know (that) he is right.} « Je sais qu'il a raison. »

\textbf{Attention :} à l'écrit formel (essai du Bac), il est plus élégant de conserver \eng{that}, surtout si la subordonnée est longue.

\textbf{Note (subjonctif) :} Après des adjectifs comme \eng{essential, important, vital, necessary} (souvent dans la structure \eng{It is... that...} ou en sujet), le verbe de la \eng{that}-clause peut prendre la forme du \emph{mandative subjunctive} (base verbale sans \eng{s} à la 3\textsuperscript{e} personne) : \eng{It is essential that the authorities \textbf{act} quickly.}
\end{propriete}

\begin{exemplebox}[That : exemples traduits]
\begin{itemize}
\item \eng{She said that she was tired.} « Elle a dit qu'elle était fatiguée. »
\item \eng{We hope that the rains will come soon.} « Nous espérons que les pluies arriveront bientôt. »
\item \eng{The doctor admitted that he had made a mistake.} « Le médecin a admis qu'il avait fait une erreur. »
\item \eng{I feel (that) something is wrong.} « Je sens que quelque chose ne va pas. »
\item \eng{Scientists say that the climate is changing.} « Les scientifiques disent que le climat change. »
\end{itemize}
\end{exemplebox}

\begin{attention}[That ne se traduit pas toujours par « que »... et vice versa]
Deux pièges symétriques :
\begin{enumerate}
\item En français, « que » est \textbf{obligatoire} ; en anglais, \eng{that} est \textbf{souvent omis}. Ne vous étonnez donc pas de lire \eng{I think he's ill} sans aucun mot de liaison.
\item Ne confondez pas \eng{that} conjonction (qui introduit une noun clause et peut être omis) et \eng{that} pronom relatif (vu à la section 3 : \eng{the book that I bought}, où \eng{that} est sujet ou objet du verbe de la subordonnée et ne peut pas disparaître quand il est sujet).
\end{enumerate}
Test rapide : si vous pouvez supprimer \eng{that} sans casser la phrase, c'est une conjonction ; sinon, c'est un relatif.
\end{attention}

\courssub{Les noun clauses introduites par « if » et « whether »}

Quand on rapporte une \textbf{question fermée} (dont la réponse est \emph{oui} ou \emph{non}), il n'y a pas de mot interrogatif à reprendre : on utilise \cle{if} ou \cle{whether} (« si »).

\begin{propriete}[If / whether : la question fermée rapportée]
\begin{center}
Question directe : \eng{“Are you ready?”} $\rightarrow$ Question indirecte : \eng{He asked \textbf{if/whether} I was ready.}
\end{center}
\begin{itemize}
\item \eng{if} est le plus courant à l'oral : \eng{Do you know if the bus has left?} « Sais-tu si le car est parti ? »
\item \eng{whether} est plus formel et est \textbf{obligatoire} devant \eng{or not} placé juste après, et après une préposition : \eng{Whether it rains or not, we will go.} « Qu'il pleuve ou non, nous irons. »
\item Comme le « si » français, \eng{if/whether} n'appelle \textbf{jamais d'inversion} : la subordonnée garde l'ordre sujet + verbe.
\end{itemize}
\end{propriete}

\begin{exemplebox}[If / whether : exemples traduits]
\begin{itemize}
\item \eng{He asked if I was ready.} « Il m'a demandé si j'étais prêt. »
\item \eng{I wonder whether she will come.} « Je me demande si elle viendra. »
\item \eng{Whether it rains or not, we will go.} « Qu'il pleuve ou non, nous irons. »
\item \eng{Nobody knows if the hospital is open on Sundays.} « Personne ne sait si l'hôpital est ouvert le dimanche. »
\end{itemize}
\end{exemplebox}

\begin{attention}[If « si » ou if « si » ?]
Ne confondez pas \eng{if} de la noun clause (« si » interrogatif, réponse oui/non) et \eng{if} de la condition (adverbial clause, section 5). Test : remplacez par \eng{whether}. Si la phrase reste correcte, c'est une noun clause : \eng{I don't know \textbf{whether} he will come.} ✓ Si la phrase devient bizarre, c'est une condition : \eng{\textbf{If} it rains, I will stay home} (\eng{whether} impossible ici).
\end{attention}

\courssub{Les noun clauses introduites par les mots interrogatifs (wh- words)}

Pour rapporter une \textbf{question ouverte} (avec mot interrogatif), on reprend ce mot — \eng{what, where, when, how, why, who, which, whose} — suivi de l'ordre \textbf{sujet + verbe}.

\begin{propriete}[L'ordre des mots : jamais d'inversion en question indirecte]
Dans une question \textbf{directe}, l'anglais inverse sujet et auxiliaire : \eng{Where \textbf{is the hospital}?} Mais dès que la question devient une noun clause (question indirecte), on revient à l'ordre normal \textbf{sujet + verbe}, et le verbe ne prend plus d'auxiliaire \eng{do/does/did} :
\begin{center}
\eng{Do you know \textbf{where the hospital is}?} ✓\qquad \eng{...where \textbf{is the hospital}?} ✗
\end{center}
\begin{center}
\begin{tabular}{|l|l|}\hline
\rowcolor{popblueL} Question directe (inversion) & Question indirecte (sujet + verbe) \\ \hline
Where does he live? & I don't know where he lives. \\ \hline
What does she want? & Tell me what she wants. \\ \hline
Why did he leave? & I wonder why he left. \\ \hline
When will the exam start? & Ask him when the exam will start. \\ \hline
\end{tabular}
\end{center}
Remarquez : \eng{does he live} $\rightarrow$ \eng{he live\textbf{s}} ; \eng{did he leave} $\rightarrow$ \eng{he \textbf{left}}. L'auxiliaire disparaît, mais la marque de la 3\textsuperscript{e} personne ou du passé revient sur le verbe.
\end{propriete}

\begin{exemplebox}[Wh- words : exemples traduits]
\begin{itemize}
\item \eng{I don't know what he wants.} « Je ne sais pas ce qu'il veut. »
\item \eng{Can you tell me where the post office is?} « Pouvez-vous me dire où est la poste ? »
\item \eng{She explained how the machine works.} « Elle a expliqué comment la machine fonctionne. »
\item \eng{I can't remember when the meeting starts.} « Je ne me souviens pas quand la réunion commence. »
\item \eng{Do you understand why forests must be protected?} « Comprenez-vous pourquoi les forêts doivent être protégées ? »
\end{itemize}
\end{exemplebox}

\begin{attention}[L'erreur n°1 des francophones : garder l'inversion]
Le français dit « Il m'a demandé où j'étais » sans inversion, mais beaucoup d'élèves calquent la question directe anglaise :
\begin{itemize}
\item ✗ \eng{He asked where was I.} \quad ✓ \eng{He asked where I was.}
\item ✗ \eng{I don't know what does he want.} \quad ✓ \eng{I don't know what he wants.}
\item ✗ \eng{Tell me when does the bus leave.} \quad ✓ \eng{Tell me when the bus leaves.}
\end{itemize}
Astuce : après le mot interrogatif en question indirecte, imaginez une phrase affirmative ordinaire (\eng{he lives, I was, the bus leaves}).
\end{attention}

\begin{attention}[What (sans antécédent) ≠ which (relatif, avec antécédent)]
\eng{What} signifie « ce qui / ce que » et \textbf{contient son propre antécédent} : il n'est jamais précédé d'un nom. \eng{Which}, pronom relatif (section 2), se réfère à un nom placé avant lui :
\begin{itemize}
\item \eng{\textbf{What} he said surprised me.} « Ce qu'il a dit m'a surpris. » (noun clause sujet ; pas de nom avant \eng{what})
\item \eng{The \textbf{words which} he said surprised me.} « Les mots qu'il a dits m'ont surpris. » (relative : \eng{which} renvoie à \eng{words})
\end{itemize}
On ne dit donc jamais ✗ \eng{the thing what...} ; dites \eng{the thing that/which...} ou simplement \eng{what...}.
\end{attention}

\courssub{La noun clause en position de sujet}

La subordonnée complément peut aussi occuper la place de \textbf{sujet}, avant le verbe principal :

\begin{exemplebox}[Noun clause sujet]
\begin{itemize}
\item \eng{\textbf{What he said} is true.} « Ce qu'il a dit est vrai. »
\item \eng{\textbf{That she passed the exam} surprised everyone.} « Le fait qu'elle ait réussi l'examen a surpris tout le monde. »
\item \eng{\textbf{Whether he will come} is still uncertain.} « On ne sait pas encore s'il viendra. »
\item \eng{\textbf{How we treat our environment} determines our future.} « La façon dont nous traitons notre environnement détermine notre avenir. »
\end{itemize}
\textbf{Remarque :} une noun clause sujet avec \eng{that} est soutenue ; à l'écrit, on lui préfère souvent la tournure avec \eng{it} anticipé : \eng{\textbf{It} surprised everyone \textbf{that she passed the exam}.} — structure très utile dans les essais du Bac.
\end{exemplebox}

\courssub{Tableau récapitulatif et comparaison avec le français}

\begin{center}
\begin{tabularx}{\linewidth}{|l|X|X|}\hline
\rowcolor{popblueL} Introducteur & Emploi & Exemple traduit \\ \hline
\eng{that} (omissible) & déclaration, opinion & \eng{I think (that) he is ill.} « Je pense qu'il est malade. » \\ \hline
\eng{if / whether} & question fermée rapportée & \eng{He asked if I was ready.} « Il a demandé si j'étais prêt. » \\ \hline
\eng{what} & « ce qui / ce que » & \eng{I know what you mean.} « Je sais ce que tu veux dire. » \\ \hline
\eng{where / when / how / why} & question ouverte rapportée & \eng{Tell me where you live.} « Dis-moi où tu habites. » \\ \hline
\eng{who / whose / which} & identité, choix & \eng{I don't know who called.} « Je ne sais pas qui a appelé. » \\ \hline
\end{tabularx}
\end{center}

\begin{center}
\begin{tabularx}{\linewidth}{|X|X|X|}\hline
\rowcolor{popgreenL} Français & English & Remarque \\ \hline
Je sais \textbf{qu}'il a raison. & I know \textbf{(that)} he is right. & « que » obligatoire ; \eng{that} omissible \\ \hline
Il demande \textbf{si} tu viens. & He asks \textbf{if/whether} you are coming. & pas d'inversion après \eng{if} \\ \hline
Dis-moi \textbf{où il habite}. & Tell me \textbf{where he lives}. & sujet + verbe, pas de \eng{does} \\ \hline
\textbf{Ce qu'}il dit est vrai. & \textbf{What} he says is true. & \eng{what} = « ce qui/ce que », sans antécédent \\ \hline
Qu'il pleuve \textbf{ou non}... & Whether it rains \textbf{or not}... & \eng{whether... or not}, tournure soutenue \\ \hline
\end{tabularx}
\end{center}

\courssub{Entraînement}

\begin{exoresolu}[Transformez en question indirecte]
Énoncé — Mettez chaque question directe dans une noun clause après l'amorce donnée.
\begin{enumerate}
\item “Where is the market?” — Can you tell me...?
\item “Does the clinic open on Saturdays?” — Do you know...?
\item “Why did she leave early?” — I don't understand...
\item “What time does the plane land in Douala?” — Could you tell me...?
\end{enumerate}
\tcblower
Solution détaillée :
\begin{enumerate}
\item \eng{Can you tell me \textbf{where the market is}?} — inversion supprimée : \eng{is the market} $\rightarrow$ \eng{the market is}.
\item \eng{Do you know \textbf{if the clinic opens on Saturdays}?} — question fermée : on introduit \eng{if} ; l'auxiliaire \eng{does} disparaît et le -s de la 3\textsuperscript{e} personne revient : \eng{opens}.
\item \eng{I don't understand \textbf{why she left early}.} — \eng{did she leave} $\rightarrow$ \eng{she left} (le passé est porté par \eng{left}).
\item \eng{Could you tell me \textbf{what time the plane lands in Douala}?} — \eng{does the plane land} $\rightarrow$ \eng{the plane lands}.
\end{enumerate}
Règle appliquée : mot interrogatif (ou \eng{if}) + sujet + verbe conjugué, sans auxiliaire \eng{do/does/did}.
\end{exoresolu}

\begin{exercice}[À vous de jouer]
\begin{enumerate}
\item Complétez avec \eng{that, if, whether, what, where, why} ou \eng{how} :
\begin{enumerate}
\item I believe \dots\ exercise is good for health.
\item He asked me \dots\ I had ever visited Bamenda.
\item \dots\ she said made everybody laugh.
\item Nobody knows \dots\ the meeting will take place.
\item \dots\ we like it or not, the climate is changing.
\end{enumerate}
\item Corrigez l'erreur dans chaque phrase :
\begin{enumerate}
\item I don't know where does he live.
\item She asked what was I doing.
\item Tell me when will the bus arrive.
\end{enumerate}
\item Traduisez : « Je ne sais pas ce qu'il veut. » / « Elle a dit qu'elle était fatiguée. » / « Sais-tu si l'hôpital est ouvert ? »
\end{enumerate}
\end{exercice}

\begin{corrige}[Exercice « À vous de jouer »]
\textbf{1.} a) \eng{that} (omissible) ; b) \eng{if/whether} ; c) \eng{What} (noun clause sujet, « ce qu' ») ; d) \eng{where} ; e) \eng{Whether} (obligatoire devant \eng{or not}).
\textbf{2.} a) \eng{I don't know where he lives.} ; b) \eng{She asked what I was doing.} ; c) \eng{Tell me when the bus will arrive.} — dans les trois cas, on supprime l'inversion et l'auxiliaire de la question directe.
\textbf{3.} \eng{I don't know what he wants.} / \eng{She said (that) she was tired.} / \eng{Do you know if the hospital is open?}
\end{corrige}

\begin{savaistu}[« Whether... or not » : la tournure qui impressionne au Bac]
Dans les essais d'argumentation du Bac, la structure \eng{whether... or not} est très appréciée des correcteurs car elle introduit élégamment une alternative ou une concession : \eng{Whether we like it or not, climate change is real.} « Que cela nous plaise ou non, le changement climatique est réel. » Autre modèle utile : \eng{Whether rich or poor, everyone needs clean water.} Placez-la en début de paragraphe pour ouvrir un développement équilibré.
\end{savaistu}

\begin{aretenir}[L'essentiel sur les noun clauses]
\begin{itemize}
\item Une noun clause = une proposition qui tient la place d'un \textbf{nom} (sujet, objet, attribut).
\item Trois introducteurs : \eng{that} (déclaration, omissible), \eng{if/whether} (question fermée), les \eng{wh- words} (question ouverte).
\item Règle d'or : dans la subordonnée, toujours \textbf{sujet + verbe} — jamais d'inversion, jamais d'auxiliaire \eng{do/does/did}.
\item \eng{what} = « ce qui/ce que » sans antécédent ; \eng{which} = relatif avec antécédent.
\item En sujet : \eng{What he said is true.} / \eng{That she passed surprised everyone.} (ou \eng{It surprised everyone that...}).
\end{itemize}
\end{aretenir}

\coursec{Adverbial clauses : temps, cause, but, concession, condition}

Dans la section précédente, nous avons étudié les \emph{noun clauses}, ces subordonnées qui se comportent comme un nom : elles sont sujet ou complément du verbe principal (\eng{What he said surprised me}). Nous passons maintenant à la troisième grande famille de propositions subordonnées : les \emph{adverbial clauses}, c'est-à-dire les \cle{subordonnées circonstancielles}. Comme leur nom l'indique, elles jouent le rôle d'un \textbf{adverbe} ou d'un \textbf{complément circonstanciel} : elles précisent \emph{quand}, \emph{pourquoi}, \emph{dans quel but}, \emph{malgré quoi} ou \emph{à quelle condition} l'action principale se déroule. C'est un point de grammaire capital pour le Bac, car il permet de construire des phrases complexes, riches et nuancées — exactement ce que les correcteurs attendent dans une \emph{composition}.

\courssub{Observation : à la découverte des circonstancielles}

Lisons d'abord ce court texte sur un projet environnemental dans le Nord-Ouest du Cameroun, puis observons les phrases en gras.

\begin{textebox}[A Green Project in Bamenda]
\textbf{When the rainy season begins}, the students of Government High School plant trees around the campus. \textbf{Because the soil is poor on the hill}, they add natural fertiliser before putting the young plants in the ground. They work every Saturday morning \textbf{so that the school stays green and shady}. \textbf{Although the work is hard}, nobody complains: everyone knows that trees protect the environment. Their teacher often tells them: “\textbf{If you protect nature, nature will protect you}.” The students have decided to continue the project \textbf{until every corner of the hill is covered with trees}.
\end{textebox}

\textbf{Glossaire :} \emph{soil} (n.) : sol ; \emph{fertiliser} (n.) : engrais ; \emph{shady} (adj.) : ombragé ; \emph{to complain} (v.) : se plaindre ; \emph{to protect} (v.) : protéger.

\textbf{Questions de compréhension :}
\begin{enumerate}
\item When do the students plant trees?
\item Why do they add fertiliser?
\item What is the purpose of their work?
\item What does the teacher say about nature?
\end{enumerate}

\begin{definition}[Adverbial clause]
Une \cle{adverbial clause} (subordonnée circonstancielle) est une proposition subordonnée introduite par une \textbf{conjonction de subordination} (\eng{when, because, so that, although, if...}). Elle joue le même rôle qu'un adverbe : elle répond aux questions \emph{when?} (quand), \emph{why?} (pourquoi), \emph{what for?} (dans quel but), \emph{in spite of what?} (malgré quoi), \emph{on what condition?} (à quelle condition). Elle contient toujours son propre sujet et son propre verbe conjugué.
\end{definition}

\courssub{Les cinq grandes relations logiques}

\begin{exemplebox}[Un exemple pour chaque relation]
\begin{itemize}
\item \textbf{Temps} : \eng{When the rain stops, we will plant the trees.} « Quand la pluie s'arrêtera, nous planterons les arbres. »
\item \textbf{Cause} : \eng{Because the soil is poor, farmers use fertiliser.} « Parce que le sol est pauvre, les agriculteurs utilisent de l'engrais. »
\item \textbf{But} : \eng{They recycle so that the town stays clean.} « Ils recyclent afin que la ville reste propre. »
\item \textbf{Concession} : \eng{Although he is poor, he is happy.} « Bien qu'il soit pauvre, il est heureux. »
\item \textbf{Condition} : \eng{Unless we act now, the forest will disappear.} « À moins que nous n'agissions maintenant, la forêt disparaîtra. »
\end{itemize}
\end{exemplebox}

\begin{center}
\begin{tabularx}{\linewidth}{|l|X|X|}
\hline
\rowcolor{popblueL} \textbf{Relation} & \textbf{Connecteurs} & \textbf{Exemple} \\
\hline
Temps (time) & when, while, as soon as, before, after, until & \eng{When the rain stops, we will plant.} \\
\hline
Cause (reason) & because, since, as & \eng{Since the soil is poor, they use fertiliser.} \\
\hline
But (purpose) & to, in order to (+ BV), so that (+ proposition) & \eng{They recycle to keep the town clean.} \\
\hline
Concession & although, even though, whereas & \eng{Although it rained, they worked.} \\
\hline
Condition & if, unless, provided that & \eng{If we act now, we will save the forest.} \\
\hline
\end{tabularx}
\end{center}

\begin{popfigure}
\begin{center}
\begin{tikzpicture}[node distance=0.5cm]
\node[draw=popblue, fill=popblueL, rounded corners, align=center, minimum width=3.2cm, minimum height=1.1cm] (time) {\textbf{Time}\\ when, while, as soon as};
\node[draw=popgreen, fill=popgreenL, rounded corners, align=center, minimum width=3.2cm, minimum height=1.1cm, right=of time] (reason) {\textbf{Reason}\\ because, since, as};
\node[draw=popgold, fill=popgoldL, rounded corners, align=center, minimum width=3.2cm, minimum height=1.1cm, right=of reason] (purpose) {\textbf{Purpose}\\ to, so that};
\node[draw=poppurple, fill=poppurpleL, rounded corners, align=center, minimum width=3.2cm, minimum height=1.1cm, below=0.7cm of time] (conc) {\textbf{Concession}\\ although, even though};
\node[draw=poporange, fill=poporangeL, rounded corners, align=center, minimum width=3.2cm, minimum height=1.1cm, right=of conc] (cond) {\textbf{Condition}\\ if, unless, provided that};
\node[draw=popdark, fill=popcream, rounded corners, align=center, minimum width=3.2cm, minimum height=1.1cm, right=of cond] (main) {\textbf{Main clause}\\ the action itself};
\draw[-{Stealth}, thick, popmuted] (time) -- (reason);
\draw[-{Stealth}, thick, popmuted] (reason) -- (purpose);
\draw[-{Stealth}, thick, popmuted] (conc) -- (cond);
\draw[-{Stealth}, thick, popmuted] (cond) -- (main);
\end{tikzpicture}
\end{center}
\legende{Les cinq relations logiques exprimées par les subordonnées circonstancielles}
\end{popfigure}

\courssub{Étude détaillée de chaque relation}

\textbf{1. Les circonstancielles de temps (time clauses).}\par
Elles situent l'action dans le temps. Les connecteurs essentiels sont : \eng{when} (quand), \eng{while} (pendant que), \eng{as soon as} (dès que), \eng{before} (avant que), \eng{after} (après que), \eng{until} (jusqu'à ce que).

\eng{While the doctor was speaking, the nurses took notes.} « Pendant que le médecin parlait, les infirmières prenaient des notes. »

\eng{We waited until the ambulance arrived.} « Nous avons attendu jusqu'à ce que l'ambulance arrive. »

\textbf{2. Les circonstancielles de cause (reason clauses).}\par
Elles répondent à la question \emph{why?}. On utilise \eng{because} (parce que, la cause la plus directe), \eng{since} et \eng{as} (puisque, comme — cause présentée comme évidente, souvent en tête de phrase).

\eng{As the hospital was full, they sent the patients to Douala.} « Comme l'hôpital était plein, ils ont envoyé les patients à Douala. »

\textbf{3. Les circonstancielles de but (purpose clauses).}\par
Attention à la construction : \eng{to} et \eng{in order to} sont suivis de la \textbf{base verbale} (même sujet dans les deux propositions), tandis que \eng{so that} est suivi d'une \textbf{proposition complète} (sujet différent possible).

\eng{She eats fruit to stay healthy.} « Elle mange des fruits pour rester en bonne santé. »

\eng{The nurse speaks slowly so that the old man can understand.} « L'infirmière parle lentement pour que le vieil homme puisse comprendre. »

\textbf{4. Les circonstancielles de concession.}\par
Elles expriment une opposition : \eng{although} et \eng{even though} (bien que + indicatif en anglais !), \eng{whereas} (alors que, pour opposer deux faits).

\eng{Whereas the city is noisy, the village is peaceful.} « Alors que la ville est bruyante, le village est paisible. »

\textbf{5. Les circonstancielles de condition.}\par
\eng{if} (si), \eng{unless} (à moins que = \emph{if... not}), \eng{provided that} (à condition que).

\eng{Provided that you take your medicine, you will recover quickly.} « À condition que tu prennes tes médicaments, tu guériras vite. »

\courssub{La règle d'or : pas de futur après when, if, as soon as...}

\begin{propriete}[Pas de futur dans la circonstancielle]
Quand la proposition principale est au \textbf{futur} (\eng{will + BV}), la subordonnée introduite par \cle{when}, \cle{if}, \cle{as soon as}, \cle{before}, \cle{after}, \cle{until} ou \cle{unless} se met au \textbf{présent simple} (ou au present perfect), jamais au futur.

\eng{If it rains tomorrow, we will stay home.} « S'il pleut demain, nous resterons à la maison. »

\eng{When I arrive in Buea, I will call you.} « Quand j'arriverai à Buea, je t'appellerai. »
\end{propriete}

\begin{popfigure}
\begin{center}
\begin{tikzpicture}
\draw[-{Stealth}, very thick, popdark] (0,0) -- (10,0);
\fill[popblue] (3,0) circle (0.12);
\node[above, align=center, popblue] at (3,0.15) {\textbf{Subordonnée}\\ \eng{When + present simple}};
\fill[poporange] (7,0) circle (0.12);
\node[above, align=center, poporange] at (7,0.15) {\textbf{Principale}\\ \eng{will + base verbale}};
\node[below, popmuted] at (0.2,0) {\footnotesize past};
\node[below, popmuted] at (9.8,0) {\footnotesize future};
\draw[-{Stealth}, thick, popgreen] (3,-0.7) -- (7,-0.7);
\node[below, popgreen, align=center] at (5,-0.75) {\eng{When he comes, I will tell him.}};
\end{tikzpicture}
\end{center}
\legende{La subordonnée au présent précède l'action future de la principale}
\end{popfigure}

\begin{exemplebox}[Exemples traduits]
\begin{itemize}
\item \eng{Although it was raining, they continued planting.} « Bien qu'il pleuvait, ils ont continué à planter. »
\item \eng{Unless you work hard, you will fail.} « À moins que tu ne travailles dur, tu échoueras. »
\item \eng{As soon as the doctor arrives, we will start.} « Dès que le médecin arrivera, nous commencerons. »
\item \eng{Before you leave, wash your hands.} « Avant que tu partes, lave-toi les mains. »
\end{itemize}
\end{exemplebox}

\begin{attention}[Le piège du futur après if et when]
Les francophones calquent le futur français (« quand il \emph{viendra} ») et écrivent :

✗ \eng{If it will rain, we will stay home.} → ✓ \eng{If it rains, we will stay home.}

✗ \eng{When he will come, I will tell him.} → ✓ \eng{When he comes, I will tell him.}

En anglais, le futur ne se met \textbf{que} dans la principale. Retenez : \emph{“If or when, no will within!”}
\end{attention}

\begin{attention}[Jamais although + but dans la même phrase]
En français on serait tenté de doubler le connecteur, mais l'anglais n'en admet qu'un seul :

✗ \eng{Although he is sick, but he works.} → ✓ \eng{Although he is sick, he works.} ou ✓ \eng{He is sick, but he works.}

De même, n'écrivez jamais ✗ \eng{Because he was tired, so he slept.} → ✓ \eng{Because he was tired, he slept.}
\end{attention}

\courssub{Position de la subordonnée et ponctuation}

\begin{propriete}[La règle de la virgule]
La circonstancielle a deux positions possibles :
\begin{itemize}
\item \textbf{En tête de phrase} : elle est suivie d'une \textbf{virgule}. \eng{Although he is poor, he is happy.}
\item \textbf{Après la principale} : il n'y a \textbf{pas de virgule}. \eng{He is happy although he is poor.}
\end{itemize}
Le sens ne change pas ; seule l'insistance change (la première information est mise en relief).
\end{propriete}

\begin{methode}[Construire une phrase complexe en quatre étapes]
\begin{enumerate}
\item \textbf{Choisir la relation logique} : est-ce une cause, un but, une concession, une condition, un repère de temps ?
\item \textbf{Choisir le connecteur} adapté : \eng{because} pour la cause, \eng{so that} pour le but, \eng{although} pour la concession, \eng{if} / \eng{unless} pour la condition.
\item \textbf{Vérifier la concordance des temps} : si la principale est au futur, la circonstancielle de temps ou de condition se met au présent simple.
\item \textbf{Placer la virgule} si et seulement si la subordonnée ouvre la phrase.
\end{enumerate}
\end{methode}

\begin{exoresolu}[Relier les phrases]
Reliez les deux phrases avec le connecteur indiqué et faites les transformations nécessaires.

1. It was raining. They continued planting. (\emph{although})

2. You do not work hard. You will fail. (\emph{unless})

3. The doctor will arrive. We will start immediately. (\emph{as soon as})

4. The soil is poor. Farmers use fertiliser. (\emph{because})
\tcblower
\textbf{Solutions :}

1. \eng{Although it was raining, they continued planting.} — Concession : \eng{although} en tête, donc virgule ; un seul connecteur (pas de \eng{but}).

2. \eng{Unless you work hard, you will fail.} — \eng{unless} = \eng{if you do not work hard} ; la négation disparaît car \eng{unless} la contient déjà ; présent simple après \eng{unless}, futur dans la principale.

3. \eng{As soon as the doctor arrives, we will start.} — Pas de futur après \eng{as soon as} : ✗ \eng{will arrive} → ✓ \eng{arrives}.

4. \eng{Because the soil is poor, farmers use fertiliser.} — Cause en tête de phrase, donc virgule après la subordonnée.
\end{exoresolu}

\begin{exercice}[À vous de jouer]
1. Complétez avec \eng{when, because, so that, although, unless} :

a) \eng{... we protect the forest, many animals will lose their homes.} (condition négative)

b) \eng{They planted trees ... the school would stay shady.}

c) \eng{... she was tired, she finished her work at the health centre.}

2. Corrigez les erreurs :

a) \eng{If it will rain tomorrow, the match will be cancelled.}

b) \eng{Although he is rich, but he is not generous.}

3. Traduisez : « Dès que la pluie s'arrêtera, nous sortirons. »
\end{exercice}

\begin{corrige}[Exercice : À vous de jouer]
1. a) \eng{Unless we protect the forest, many animals will lose their homes.} — Condition négative : \eng{unless} = \eng{if we do not protect}. Présent simple après \eng{unless}.

b) \eng{so that} (but, suivi d'une proposition complète avec sujet \eng{the school}).

c) \eng{Although} (concession, opposition entre fatigue et action accomplie).

2. a) ✗ \eng{will rain} → ✓ \eng{If it rains tomorrow, the match will be cancelled.} (Règle d'or : pas de \eng{will} après \eng{if}).

b) Supprimer \eng{but} : ✓ \eng{Although he is rich, he is not generous.} (Un seul connecteur de concession).

3. \eng{As soon as the rain stops, we will go out.} — Présent simple \eng{stops} après \eng{as soon as}, futur \eng{will go} dans la principale.
\end{corrige}

\begin{aretenir}[L'essentiel de la section]
\begin{itemize}
\item Une \cle{adverbial clause} précise les circonstances de l'action : temps (\eng{when, while, as soon as, until}), cause (\eng{because, since, as}), but (\eng{to, in order to, so that}), concession (\eng{although, even though, whereas}), condition (\eng{if, unless, provided that}).
\item \textbf{Règle d'or} : pas de \eng{will} après \eng{when, if, as soon as, before, after, until, unless} — on utilise le présent simple.
\item Un seul connecteur par phrase : jamais \eng{although... but}, jamais \eng{because... so}.
\item Subordonnée en tête → virgule ; subordonnée en fin → pas de virgule.
\end{itemize}
\end{aretenir}

\coursec{Comparaison français/anglais et entraînement au Bac}

Dans la section précédente, nous avons étudié les \emph{adverbial clauses} (subordonnées circonstancielles de temps, cause, but, concession et condition) et leurs connecteurs : \eng{when}, \eng{because}, \eng{so that}, \eng{although}, \eng{if}... Vous savez désormais construire les trois grandes familles de clauses : \emph{relative}, \emph{noun} et \emph{adverbial clauses}. Il reste une dernière étape, décisive pour la réussite au Baccalauréat : comparer systématiquement le fonctionnement du français et de l'anglais, car c'est précisément là que se situent les erreurs des candidats francophones, puis s'entraîner aux exercices types de l'examen.

\courssub{Observer d'abord : quatre pièges de traduction}

\begin{textebox}[A letter to the editor — The Green Voice, Douala]
Dear Editor,

Although our city is growing fast, the air we breathe is getting dirtier every day. What worries me most is the number of old taxis whose engines produce thick black smoke. The mayor, whose office is near the central market, promised last year that he would act before the situation became worse. However, nothing has changed. The doctor I spoke to yesterday told me that respiratory diseases are increasing among children who live near busy roads. If the authorities do not react soon, what we fear today will become a tragedy tomorrow. I believe that every citizen should write to the newspapers so that our leaders finally understand what is at stake.

Yours faithfully,

Marlyse N., Form Five student
\end{textebox}

\textbf{Glossaire.} \emph{to breathe} (v.) : respirer ; \emph{engine} (n.) : moteur ; \emph{thick} (adj.) : épais ; \emph{respiratory disease} (n.) : maladie respiratoire ; \emph{to be at stake} (expr.) : être en jeu ; \emph{to react} (v.) : réagir.

\textbf{Observation.} Relevez dans la lettre : (1) \eng{Although our city is growing} — le français dirait « bien que notre ville \emph{grandisse} » avec un subjonctif ; (2) \eng{taxis whose engines produce smoke} — le français dirait « des taxis \emph{dont} les moteurs... » ; (3) \eng{What worries me most} — le français dirait « \emph{ce qui} m'inquiète le plus » ; (4) \eng{The doctor I spoke to} — le français dirait « le médecin \emph{à qui} j'ai parlé » ou « \emph{dont} j'ai parlé ». Quatre structures françaises, quatre solutions anglaises différentes. C'est tout l'objet de cette section.

\courssub{Comparaison 1 : le subjonctif français n'existe pas en anglais}

\begin{propriete}[Pas de subjonctif après although, before, so that]
Le français impose le subjonctif après \emph{bien que}, \emph{avant que}, \emph{pour que}, \emph{à condition que}. L'anglais, lui, utilise l'\cle{indicatif} (présent, passé) après \eng{although}, \eng{before}, \eng{if}, ou la structure \cle{should + base verbale} dans un style soutenu. Il ne faut donc jamais chercher à « traduire » le subjonctif : on emploie simplement le temps logique de l'indicatif.
\end{propriete}

\begin{exemplebox}[Exemples traduits]
\eng{Although he is sick, he works.} « Bien qu'il \emph{soit} malade, il travaille. »

\eng{Although she was tired, she finished her essay.} « Bien qu'elle \emph{fût} fatiguée, elle a terminé sa rédaction. »

\eng{I will leave before the rain starts.} « Je partirai avant que la pluie ne \emph{commence}. »

\eng{The doctor explained the risks so that everyone should understand.} « Le médecin a expliqué les risques pour que tout le monde \emph{comprenne}. » (style soutenu ; en langue courante : \eng{so that everyone could understand}).
\end{exemplebox}

\begin{attention}[Ne jamais traduire le subjonctif]
✗ \eng{Although he be rich...} est archaïque ou incorrect en anglais moderne. ✓ \eng{Although he is rich, he lives simply.} De même, après \eng{if} au sens conditionnel, on n'utilise jamais le futur : ✗ \eng{If it will rain} → ✓ \eng{If it rains}. Le français non plus ne dit pas « si il pleuvra » : les deux langues interdisent le futur après \emph{si} conditionnel.
\end{attention}

\courssub{Comparaison 2 : « dont » n'est pas un mot anglais}

Le français possède un pronom relatif tout prêt, \emph{dont}, qui n'a \textbf{aucun équivalent unique} en anglais. Selon le sens, « dont » se rend de trois façons :

\begin{center}
\begin{tabularx}{\linewidth}{|X|X|X|}
\hline
\rowcolor{popblueL} Français & English & Remarque \\
\hline
la maison \emph{dont} le toit est rouge & the house \textbf{whose} roof is red & possession : \eng{whose} + nom \\
\hline
l'homme \emph{dont} je parle & the man I am talking \textbf{about} & verbe + préposition : la préposition va à la fin \\
\hline
l'homme \emph{dont} je parle (soutenu) & the man \textbf{about whom} I am talking & registre formel : préposition + \eng{whom} \\
\hline
le bruit \emph{dont} il se plaint & the noise he complains \textbf{about} & idem : préposition rejetée à la fin \\
\hline
le livre \emph{dont} la couverture est bleue & the book \textbf{whose} cover is blue & possession \\
\hline
\end{tabularx}
\end{center}

\begin{exemplebox}[Exemple traduit]
\eng{The house whose roof is red belongs to my uncle.} « La maison \emph{dont} le toit est rouge appartient à mon oncle. »
\end{exemplebox}

\begin{attention}[L'erreur n° 1 des francophones : traduire « dont » mot à mot]
✗ \eng{The man whose I talk} n'existe pas. \eng{Whose} doit être \textbf{suivi d'un nom} (possession). Pour « parler de quelqu'un », la préposition \eng{about} est indispensable : ✓ \eng{The man I am talking about} ou, au registre soutenu, ✓ \eng{The man about whom I am talking}. Règle pratique : demandez-vous « de quoi s'agit-il ? » — si le verbe français se construit avec \emph{de} (parler de, se souvenir de, avoir besoin de), gardez la préposition anglaise correspondante (\eng{about}, \eng{of}, \eng{for}) et placez-la en fin de clause ou devant \eng{whom}/\eng{which}.
\end{attention}

\courssub{Comparaison 3 : « ce qui / ce que » = what, jamais which}

\begin{propriete}[What, le relatif sans antécédent]
\emph{Ce qui} (sujet) et \emph{ce que} (COD) se traduisent par \cle{what}, qui introduit une \emph{noun clause}. \eng{What} signifie « la chose que/qui » : il contient son propre antécédent. On ne peut donc jamais le faire précéder d'un nom, et on ne peut pas utiliser \eng{which} à sa place.
\end{propriete}

\begin{exemplebox}[Exemples traduits]
\eng{Here is what I think.} « Voici \emph{ce que} je pense. »

\eng{What worries me is the pollution of the river.} « \emph{Ce qui} m'inquiète, c'est la pollution du fleuve. »

\eng{I don't understand what he said.} « Je ne comprends pas \emph{ce qu'}il a dit. »

\eng{What she cooks is always delicious.} « \emph{Ce qu'}elle cuisine est toujours délicieux. »
\end{exemplebox}

\begin{attention}[What ou which ?]
✗ \eng{Here is which I think.} ✓ \eng{Here is what I think.} \eng{Which} exige un antécédent exprimé : \eng{He bought a bike, which surprised me.} (« il a acheté un vélo, \emph{ce qui} m'a surpris » — ici \eng{which} reprend toute la phrase précédente, avec une virgule). Mais sans antécédent, c'est toujours \eng{what}.
\end{attention}

\courssub{Comparaison 4 : l'anglais interdit la reprise du sujet par un pronom}

\begin{propriete}[Un seul sujet par clause]
Le français courant reprend souvent le sujet par un pronom : « \emph{Mon frère, il} est venu », « \emph{les gens que j'ai vus, ils} étaient malades ». L'anglais l'interdit absolument : le pronom relatif ou le nom suffit. ✗ \eng{My brother he came.} ✓ \eng{My brother came.} ✗ \eng{The people who I saw they were sick.} ✓ \eng{The people who I saw were sick.}
\end{propriete}

\begin{propriete}[La phrase complexe équilibrée au Bac]
Une phrase complexe bien construite = \cle{1 main clause} + \cle{1 ou 2 subordinate clauses maximum}. Au-delà, la phrase devient lourde, confuse et risquée grammaticalement. Préférez deux phrases claires à une phrase-monstre : \eng{Although the factory pollutes the river, the government has not closed it, because it employs two hundred workers.} (1 principale + 2 subordonnées : parfait).
\end{propriete}

\begin{popfigure}
\begin{center}
\begin{tikzpicture}[scale=1.0, every node/.style={font=\small}]
\node[draw=popblue, fill=popblueL, rounded corners, align=center, minimum width=3.2cm] (fr) at (0,0) {\textbf{Français}};
\node[draw=popgreen, fill=popgreenL, rounded corners, align=center, minimum width=3.2cm] (en) at (8,0) {\textbf{English}};
\node[align=center] (a1) at (0,-1.4) {dont (possession)};
\node[align=center] (b1) at (8,-1.4) {whose + nom};
\node[align=center] (a2) at (0,-2.6) {dont je parle};
\node[align=center] (b2) at (8,-2.6) {I am talking about};
\node[align=center] (a3) at (0,-3.8) {ce qui / ce que};
\node[align=center] (b3) at (8,-3.8) {what};
\node[align=center] (a4) at (0,-5.0) {bien que + subjonctif};
\node[align=center] (b4) at (8,-5.0) {although + indicatif};
\node[align=center] (a5) at (0,-6.2) {si + futur : interdit};
\node[align=center] (b5) at (8,-6.2) {if + will : interdit};
\draw[-{Stealth}, thick, poporange] (a1) -- (b1);
\draw[-{Stealth}, thick, poporange] (a2) -- (b2);
\draw[-{Stealth}, thick, poporange] (a3) -- (b3);
\draw[-{Stealth}, thick, poporange] (a4) -- (b4);
\draw[-{Stealth}, thick, poppurple] (a5) -- (b5);
\end{tikzpicture}
\end{center}
\legende{Tableau de correspondances français/anglais : les quatre faux amis grammaticaux à maîtriser pour le Bac.}
\end{popfigure}

\begin{savaistu}[Le critère caché des correcteurs]
Dans l'épreuve d'\emph{essay writing} du Baccalauréat, les grilles de correction valorisent explicitement le critère \emph{range of structures} (variété des structures). Un candidat qui aligne uniquement des phrases simples (\eng{The air is dirty. People are sick.}) perd des points, même sans faute. Un candidat qui emploie correctement une relative (\eng{the smoke which comes from factories}), une noun clause (\eng{what the government should do}) et une adverbial clause (\eng{although progress has been made}) démontre sa maîtrise de la langue et est récompensé. Trois subordonnées bien placées dans un essay peuvent faire gagner plusieurs points.
\end{savaistu}

\courssub{Méthode : l'exercice de transformation au Bac}

\begin{methode}[Combiner deux phrases simples en une phrase complexe]
\begin{enumerate}
\item \textbf{Repérer l'élément commun} aux deux phrases (un nom, une idée répétée).
\item \textbf{Choisir le bon connecteur} : pronom relatif (\eng{who}, \eng{which}, \eng{whose}, \eng{that}) si l'élément commun est un nom ; conjonction (\eng{because}, \eng{although}, \eng{when}, \eng{if}) si le lien est logique (cause, concession, temps, condition) ; \eng{what} si le français dirait « ce qui/ce que ».
\item \textbf{Remplacer} l'élément répété par le relatif ou supprimer la répétition.
\item \textbf{Vérifier} : virgules (relative explicative = virgules ; déterminative = pas de virgules), concordance des temps, et un seul sujet par clause.
\item \textbf{Relire} la phrase à voix haute : si elle contient plus de deux subordonnées, la couper.
\end{enumerate}
\end{methode}

\begin{exoresolu}[Transformation type Bac]
Énoncé — Combine each pair of sentences into one complex sentence.

(a) \eng{The river is polluted. It crosses our village.}

(b) \eng{The nurse works in Buea. Her clinic helps hundreds of children.}

(c) \eng{He smokes. He knows it is dangerous.}
\tcblower
Solution détaillée :

(a) Élément commun : \eng{the river} → relatif \eng{which} (chose, clause déterminative, sans virgule). ✓ \eng{The river which crosses our village is polluted.} (« La rivière qui traverse notre village est polluée. »)

(b) Élément commun : \eng{the nurse / her} → possession → \eng{whose}. ✓ \eng{The nurse whose clinic helps hundreds of children works in Buea.} (« L'infirmière \emph{dont} la clinique aide des centaines d'enfants travaille à Buea. »)

(c) Lien logique de concession : ✓ \eng{Although he knows it is dangerous, he smokes.} (« Bien qu'il \emph{sache} que c'est dangereux, il fume. » — indicatif anglais, pas de subjonctif.)
\end{exoresolu}

\courssub{Stratégie de lecture : découper les longues phrases du Bac}

\begin{methode}[Identifier les clauses dans un texte difficile]
\begin{enumerate}
\item Soulignez les \textbf{verbes conjugués} : chaque verbe = une clause.
\item Entourez les \textbf{connecteurs} (\eng{who}, \eng{which}, \eng{that}, \eng{whose}, \eng{what}, \eng{although}, \eng{because}, \eng{if}, \eng{when}...) : ils annoncent les subordonnées.
\item Identifiez la \textbf{main clause} : c'est elle qui porte le sens principal ; lisez-la d'abord seule.
\item Réintégrez ensuite les subordonnées une par une en posant la question : qui ? quoi ? pourquoi ? quand ?
\end{enumerate}
Exemple : \eng{The report, which was published last week, shows that the air in Douala, although it has improved, remains dangerous for children who have asthma.} → 5 verbes conjugués, 5 clauses ; main clause : \eng{The report shows that the air in Douala remains dangerous}.
\end{methode}

\courssub{Entraînement sur le thème environnement et santé}

\begin{exercice}[Exercice A — Phrases à trous]
Complétez avec \eng{who}, \eng{which}, \eng{whose}, \eng{what}, \eng{although} ou \eng{because}.

1. The students \dots planted trees were congratulated.

2. \dots he was ill, he took the exam.

3. I don't know \dots causes this disease.

4. The village \dots water is polluted needs help.

5. She stayed at home \dots she had a fever.
\end{exercice}

\begin{exercice}[Exercice B — Transformation]
Combinez les paires de phrases en une phrase complexe.

1. The factory pollutes the air. It was built in 2010.

2. The doctor is famous. His research saved many lives.

3. The product is expensive. Many people buy it. (concession)
\end{exercice}

\begin{exercice}[Exercice C — Correction d'erreurs]
Corrigez chaque phrase.

1. \eng{The man whose I spoke is a doctor.}

2. \eng{Although he is sick, but he works.}

3. \eng{Here is which I think about pollution.}

4. \eng{My cousin she lives in Bamenda.}

5. \eng{If it will rain, the match will be cancelled.}
\end{exercice}

\begin{corrige}[Exercices A, B et C]
\textbf{A.} 1. \eng{who} (personnes) ; 2. \eng{Although} (concession, jamais avec \eng{but} dans la même phrase) ; 3. \eng{what} (« ce qui ») ; 4. \eng{whose} (possession : l'eau du village) ; 5. \eng{because} (cause).

\textbf{B.} 1. \eng{The factory which was built in 2010 pollutes the air.} 2. \eng{The doctor whose research saved many lives is famous.} 3. \eng{Although the product is expensive, many people buy it.}

\textbf{C.} 1. \eng{The man I spoke to is a doctor.} / \eng{The man to whom I spoke is a doctor.} (\eng{whose} exige un nom). 2. \eng{Although he is sick, he works.} (jamais \eng{although... but} ensemble). 3. \eng{Here is what I think about pollution.} 4. \eng{My cousin lives in Bamenda.} (pas de reprise du sujet). 5. \eng{If it rains, the match will be cancelled.} (pas de futur après \eng{if}).
\end{corrige}

\begin{aretenir}[L'essentiel de la leçon]
\begin{itemize}
\item Pas de subjonctif en anglais : \eng{although} + indicatif ; jamais de futur après \eng{if}.
\item « Dont » = \eng{whose} + nom (possession) ou préposition en fin de clause (\eng{the man I talked about}).
\item « Ce qui / ce que » = \eng{what} ; jamais \eng{which} sans antécédent.
\item Jamais de reprise du sujet par un pronom ; une phrase complexe = 1 principale + 1 ou 2 subordonnées maximum.
\item Au Bac, la variété correcte des subordonnées est un critère de notation : entraînez-vous à combiner des phrases.
\end{itemize}
\end{aretenir}

\newpage

\coursec{Activité d'intégration}

\courssub{Objectifs de l'activité}

À l'issue de cette activité d'intégration, l'élève sera capable de :

\begin{itemize}
  \item produire un court discours oral et écrit (\eng{speech}) destiné à convaincre un auditoire réel, dans un contexte scolaire authentique ;
  \item mobiliser et combiner correctement les trois grands types de \cle{clauses} étudiés dans la leçon : \cle{relative clauses} (déterminatives et explicatives), \cle{noun clauses} (introduites par \eng{that} ou \eng{what}) et \cle{adverbial clauses} (condition, concession, temps, but) ;
  \item respecter la règle du \cle{non-futur} après \eng{if}, \eng{when}, \eng{unless}, \eng{although} ;
  \item repérer, nommer et justifier les clauses utilisées dans un texte produit par d'autres élèves (activité métalinguistique) ;
  \item travailler en groupe, répartir les rôles (rédacteur, secrétaire, porte-parole, vérificateur grammatical) et présenter une production collective devant la classe.
\end{itemize}

\courssub{Situation}

Ton lycée, le \emph{Government Bilingual High School} de Buea, possède un \eng{Environment Club} très actif. Chaque année, le club organise la \eng{Green Week}, une semaine de sensibilisation à la protection de l'environnement : nettoyage de la cour, plantation d'arbres derrière les salles de classe, tri des déchets plastiques, campagne contre le brûlage des ordures.

Cette année, le président du club a lancé un concours : chaque classe de Terminale doit présenter, lors de la cérémonie d'ouverture de la \eng{Green Week}, un court discours intitulé \eng{Save our school environment}. Le meilleur discours sera affiché dans la salle polyvalente et lu à la radio du lycée.

Voici l'affiche officielle du concours, publiée par le club :

\begin{textebox}[Environment Club — Official notice]
\textbf{GREEN WEEK SPEECH CONTEST}

Dear students,

The Environment Club invites every Form Five and Upper Sixth class to take part in our annual speech contest. Your task is to write and deliver a short speech entitled “Save our school environment”.

Your speech must be between eight and ten lines long. It must convince the audience that protecting our school environment is everybody's duty. Show us what you believe, tell us what we can do together, and explain why action is urgent.

The winning speech will be displayed in the multipurpose hall and read on the school radio during the opening ceremony on Monday.

Deadline: Friday at 4 p.m. — Good luck to all participants!

The Environment Club Committee
\end{textebox}

\begin{definition}[Glossaire utile]
\begin{itemize}
  \item \eng{speech} (n.) : discours ; \eng{to deliver a speech} : prononcer un discours.
  \item \eng{to take part in} : participer à ; \eng{contest} (n.) : concours.
  \item \eng{to convince} : convaincre ; \eng{audience} (n.) : auditoire, public.
  \item \eng{duty} (n.) : devoir ; \eng{urgent} (adj.) : urgent.
  \item \eng{deadline} (n.) : date limite ; \eng{to display} : afficher, exposer.
\end{itemize}
\end{definition}

Ton groupe de quatre élèves a décidé de participer. À vous de rédiger le discours gagnant !

\courssub{Consignes}

\textbf{Tâche 1 — Compréhension de la consigne officielle.} Lis attentivement l'affiche du \eng{Environment Club}. Réponds par écrit en anglais : (a) What is the title of the speech? (b) How long must it be? (c) What will happen to the winning speech? (d) When is the deadline?

\textbf{Tâche 2 — Remue-méninges lexical.} En groupe, dressez une liste d'au moins huit mots ou expressions liés à l'environnement scolaire que vous pourriez utiliser : \eng{litter}, \eng{to plant trees}, \eng{dustbins}, \eng{plastic waste}, \eng{to sweep the classrooms}, \eng{a clean environment}, \eng{to pollute}, \eng{to protect}, \eng{health}, \eng{diseases}. Vérifiez le sens de chaque mot dans le dictionnaire.

\textbf{Tâche 3 — Planification grammaticale.} Avant de rédiger, complétez la fiche de préparation suivante (elle vous oblige à prévoir chaque clause exigée) :

\begin{center}
\begin{tabularx}{\linewidth}{|l|X|}
\hline
\rowcolor{popblueL} Clause exigée & Idée de phrase (en anglais) \\
\hline
Relative déterminative (sans virgules) & \eng{Students who drop litter...} \\
\hline
Relative explicative (entre virgules) & \eng{Our school, which is near the market,...} \\
\hline
Noun clause avec \eng{that} ou \eng{what} & \eng{We believe that...} / \eng{What we need is...} \\
\hline
Condition avec \eng{if} ou \eng{unless} & \eng{If we act now,...} \\
\hline
Concession avec \eng{although} & \eng{Although the task is hard,...} \\
\hline
\end{tabularx}
\end{center}

\textbf{Tâche 4 — Rédaction collective.} Rédigez le discours de 8 à 10 lignes. Répartition des rôles : un élève propose les idées, un deuxième rédige, un troisième vérifie la grammaire (relatives, non-futur après \eng{if}/\eng{when}, virgules de la relative explicative), un quatrième prépare la présentation orale.

\textbf{Tâche 5 — Vérification grammaticale.} Relisez votre discours et cochez : au moins deux relatives (dont une explicative entre virgules), une noun clause, deux adverbial clauses (une condition, une concession), aucun futur après \eng{if}, \eng{when}, \eng{unless} ou \eng{although}.

\textbf{Tâche 6 — Présentation et analyse croisée.} Le porte-parole de chaque groupe présente le discours à la classe (ou l'affiche est accrochée au tableau). Pendant l'écoute, chaque élève note les clauses entendues et les nomme : \eng{relative}, \eng{noun clause}, \eng{adverbial clause of condition / concession}. La classe vote pour la meilleure affiche, qui sera exposée dans la salle.

\courssub{Ressources à mobiliser}

\begin{aretenir}[Rappel express : les trois types de clauses]
\begin{itemize}
  \item \textbf{Relative clause} : introduite par \eng{who} (personnes), \eng{which} (choses), \eng{that}, \eng{whose}, \eng{where}. \emph{Déterminative} : sans virgules, indispensable au sens (\eng{Students who litter must be punished.}). \emph{Explicative} : entre virgules, information supplémentaire, jamais \eng{that} (\eng{Our school, which is very old, needs new dustbins.}).
  \item \textbf{Noun clause} : fonctionne comme un nom (sujet ou complément), introduite par \eng{that} ou \eng{what} (\eng{We believe that change is possible.} / \eng{What we do today matters.}).
  \item \textbf{Adverbial clause} : introduite par \eng{if}, \eng{unless} (condition), \eng{although}, \eng{even though} (concession), \eng{because} (cause), \eng{so that} (but), \eng{when} (temps).
\end{itemize}
\end{aretenir}

\begin{attention}[Pièges à éviter dans votre discours]
\begin{itemize}
  \item Pas de futur après \eng{if}, \eng{when}, \eng{unless}, \eng{although} : \eng{If we \textbf{act} now...} et non \eng{If we will act...}.
  \item La relative explicative exige deux virgules et interdit \eng{that} : \eng{Buea, which is a green town,...}.
  \item Ne pas traduire « ce que » par \eng{that} : \eng{What we need is...} (et non \eng{That we need is...}).
  \item \eng{Unless} = \eng{if... not} : \eng{Unless we act, nothing will change.}
\end{itemize}
\end{attention}

Lexique utile du thème \emph{Environment, Well-being and Health} : \eng{to protect the environment}, \eng{to reduce waste}, \eng{to keep our school clean}, \eng{to plant trees}, \eng{fresh air}, \eng{a healthy environment}, \eng{to fall sick}, \eng{to take action}, \eng{together}, \eng{future generations}.

\courssub{Proposition de production et corrigé commenté}

\begin{exoresolu}[Modèle de discours — “Save our school environment”]
Rédige un discours de 8 à 10 lignes respectant toutes les contraintes du concours.
\tcblower
\textbf{Production modèle :}

\eng{Dear friends, what we do today will decide the future of our school. Our school, which is one of the oldest in Buea, is becoming dirty because some students who drop litter do not respect our common home. We all know that a dirty environment causes diseases. Although the task may seem difficult, we can succeed if we work together. If every student uses the dustbins, our classrooms and our playground will stay clean. Unless we act now, nobody will act for us. Let us plant trees so that our children can breathe fresh air. When the Green Week begins, let us be the first to take action. Save our school environment: it is our duty!}

\textbf{Commentaire des choix de langue :}
\begin{itemize}
  \item \eng{what we do today} : \textbf{noun clause} sujet, introduite par \eng{what} (« ce que ») — attention, pas de futur dans la noun clause ici car elle exprime une vérité générale ; le futur \eng{will decide} porte sur la principale.
  \item \eng{which is one of the oldest in Buea} : \textbf{relative explicative}, entre virgules, avec \eng{which} (et non \eng{that}) — conforme à la règle.
  \item \eng{who drop litter} : \textbf{relative déterminative}, sans virgules, \eng{who} pour les personnes.
  \item \eng{that a dirty environment causes diseases} : \textbf{noun clause} complément de \eng{know}.
  \item \eng{Although the task may seem difficult} : \textbf{adverbial clause de concession}, présent après \eng{although}.
  \item \eng{if we work together} et \eng{If every student uses the dustbins} : \textbf{adverbial clauses de condition}, présent après \eng{if}, futur dans la principale (\eng{will stay}) — règle du non-futur respectée.
  \item \eng{Unless we act now} : condition négative, présent obligatoire.
  \item \eng{so that our children can breathe fresh air} : \textbf{adverbial clause de but}.
  \item \eng{When the Green Week begins} : subordonnée de temps au présent, futur implicite dans l'impératif qui suit.
\end{itemize}
Le discours compte dix lignes, contient toutes les clauses exigées, et se termine par un appel à l'action (\eng{call to action}), procédé classique du discours persuasif.
\end{exoresolu}

\courssub{Bilan de l'activité}

Cette activité d'intégration a permis de réinvestir l'ensemble des savoir-faire de la leçon dans une tâche finale authentique et motivante : écrire et prononcer un discours réel pour un public réel. Sur le plan grammatical, l'élève a manipulé simultanément les trois grandes familles de clauses — relatives (déterminatives et explicatives), noun clauses et adverbial clauses — en veillant aux points critiques pour un francophone : les virgules de la relative explicative, l'interdiction de \eng{that} dans celle-ci, le choix entre \eng{that} et \eng{what}, et surtout le non-futur après les conjonctions de subordination. Sur le plan méthodologique, la planification guidée (fiche de préparation), la répartition des rôles et la vérification croisée ont développé l'autonomie et la coopération. Enfin, l'analyse métalinguistique des discours des autres groupes (repérer et nommer les clauses) prépare directement aux épreuves de grammaire et de rédaction du baccalauréat, où la capacité à articuler des phrases complexes est un critère majeur d'évaluation.

\newpage

\coursec{Méthodes et exercices résolus}

\coursec{Lesson 26 — Grammar: Revise clauses}

\courssub{Observation et notion de clause}

\begin{textebox}[Texte d'observation : Article de presse camerounais]
“Environmental experts warn that the rapid deforestation in the South Region threatens biodiversity. The Minister of Forestry, who attended the summit in Yaoundé last week, announced that new regulations will be enforced soon. Local communities, whose livelihoods depend on the forest, hope that these measures will not come too late. Unless illegal logging stops immediately, the Congo Basin could lose its most precious ecosystems.”
\end{textebox}

\begin{methode}[Qu'est-ce qu'une clause ?]
\begin{enumerate}
\item \textbf{Définition} : Une \cle{clause} (proposition) est un groupe de mots contenant un \emph{sujet} et un \emph{verbe conjugué}. Elle peut être \emph{principale} (indépendante) ou \emph{subordonnée} (dépendante d'une principale).
\item \textbf{Test de repérage} : Cherchez les verbes conjugués. Chaque verbe conjugué = une clause.
\item \textbf{Types de subordonnées} (toutes introduites par un mot-outil : pronom relatif, conjonction, mot interrogatif) :
\begin{itemize}
\item \cle{Relative clause} (proposition relative) : modifie un nom (antécédent). \eng{The forest \textbf{that covers the region} is vast.}
\item \cle{Noun clause} (proposition complétive) : fonctionne comme un nom (sujet, objet, complément). \eng{\textbf{What the minister said} is important.}
\item \cle{Adverbial clause} (proposition circonstancielle) : indique temps, cause, but, concession, condition… \eng{\textbf{Because the rain stopped}, we went out.}
\end{itemize}
\end{enumerate}
\end{methode}

\begin{exemplebox}[Analyse du texte d'observation]
\begin{tabular}{|l|l|l|}
\hline
\rowcolor{popblueL} \textbf{Extrait} & \textbf{Type de clause} & \textbf{Introducteur} \\
\hline
\eng{that the rapid deforestation… threatens biodiversity} & Noun clause (COD de \eng{warn}) & \eng{that} \\
\hline
\eng{who attended the summit in Yaoundé last week} & Relative clause (explicative, antécédent = \eng{The Minister}) & \eng{who} \\
\hline
\eng{that new regulations will be enforced soon} & Noun clause (COD de \eng{announced}) & \eng{that} \\
\hline
\eng{whose livelihoods depend on the forest} & Relative clause (explicative, possession) & \eng{whose} \\
\hline
\eng{that these measures will not come too late} & Noun clause (COD de \eng{hope}) & \eng{that} \\
\hline
\eng{Unless illegal logging stops immediately} & Adverbial clause (condition) & \eng{Unless} \\
\hline
\end{tabular}
\end{exemplebox}

\courssub{Fiche méthode 1 : Identifier et construire une relative clause}

\begin{methode}[Identifier et construire une relative clause]
Pour manipuler une proposition relative sans erreur, suis cette démarche en cinq étapes.
\begin{enumerate}
\item \textbf{Repérer l'antécédent} : c'est le nom placé juste avant le pronom relatif. Demande-toi : est-ce une personne (\cle{who}) ou une chose (\cle{which}) ? Est-ce un possessif (\cle{whose}) ou un lieu/temps (\cle{where}, \cle{when}) ?
\item \textbf{Choisir le pronom relatif} :
\begin{itemize}
\item Personne sujet $\rightarrow$ \eng{who} (\eng{The doctor \textbf{who} treats me})
\item Chose sujet $\rightarrow$ \eng{which} (\eng{The report \textbf{which} explains the crisis})
\item Complément d'objet (personne/chose) $\rightarrow$ \eng{who(m)} / \eng{which} (ou \textbf{zéro}) (\eng{The book (\textbf{which}) I read})
\item Possession $\rightarrow$ \eng{whose} (\eng{The farmer \textbf{whose} land was flooded})
\item Lieu $\rightarrow$ \eng{where} (= \eng{in/at which}) (\eng{The village \textbf{where} I was born})
\item Temps $\rightarrow$ \eng{when} (= \eng{on/in which}) (\eng{The day \textbf{when} we met})
\end{itemize}
\item \textbf{Vérifier si le pronom peut être omis} : \textbf{uniquement} s'il est \emph{complément d'objet} dans une relative \emph{déterminative}. \eng{The book (which) I bought is interesting.} \textbf{Jamais} s'il est sujet (\eng{The man *who* called} — omission impossible).
\item \textbf{Décider : déterminative ou explicative ?} Si l'information est indispensable pour identifier l'antécédent $\rightarrow$ pas de virgule (on \textbf{peut} utiliser \eng{that}). Si c'est un simple commentaire $\rightarrow$ virgules obligatoires, \textbf{jamais} \eng{that}, \textbf{jamais} d'omission.
\item \textbf{Placer la préposition correctement} :
\begin{itemize}
\item Style courant $\rightarrow$ préposition en fin de relative (\eng{the girl I talked \textbf{to}})
\item Style soutenu $\rightarrow$ préposition + \eng{whom/which} (\eng{the girl \textbf{to whom} I talked})
\end{itemize}
\end{enumerate}
Réflexe Bac : avant d'écrire, traduis mentalement le français « que/qui/dont/où » et vérifie la \emph{fonction grammaticale}, pas seulement le sens.
\end{methode}

\begin{exoresolu}[Relative clauses — type Bac]
\textbf{Énoncé.} Complete each sentence with the correct relative pronoun (\eng{who, which, whose, where, when}) or leave it out if possible. Justify your choice.

1. The nurse \ldots\ works at the Buea Regional Hospital is very kind.

2. The report \ldots\ the minister published yesterday concerns air pollution.

3. That is the village \ldots\ my grandmother was born.

4. The students \ldots\ project won the prize were invited to Yaoundé.

5. The man \ldots\ I spoke to is an environmental engineer.

\tcblower
\textbf{Solution détaillée.}

\textbf{1.} \eng{The nurse \textbf{who} works at the Buea Regional Hospital is very kind.}
Raisonnement : l'antécédent \eng{the nurse} est une \emph{personne}, et le pronom est \emph{sujet} du verbe \eng{works} $\rightarrow$ \eng{who} obligatoire (omission impossible car sujet).

\textbf{2.} \eng{The report \textbf{(which/that)} the minister published yesterday concerns air pollution.}
Raisonnement : antécédent = chose (\eng{the report}) ; le pronom est \emph{complément d'objet} de \eng{published} (le sujet est \eng{the minister}) $\rightarrow$ \eng{which} ou \eng{that}, et l'omission est possible : \eng{The report the minister published\ldots}

\textbf{3.} \eng{That is the village \textbf{where} my grandmother was born.}
Raisonnement : l'antécédent est un \emph{lieu} ; dans la relative, il joue le rôle de complément de lieu (\eng{she was born \emph{in} the village}) $\rightarrow$ \eng{where} (= \eng{in which}).

\textbf{4.} \eng{The students \textbf{whose} project won the prize were invited to Yaoundé.}
Raisonnement : il y a une relation de possession (\eng{the students' project}) $\rightarrow$ \eng{whose}, qui correspond au français « dont ». Attention : \eng{whose} s'emploie aussi pour les choses : \eng{a country whose forests are disappearing}.

\textbf{5.} \eng{The man \textbf{(who(m))} I spoke to is an environmental engineer.}
Raisonnement : le pronom est complément de la préposition \eng{to} ; style courant $\rightarrow$ \eng{who} ou omission, préposition en fin ; style soutenu $\rightarrow$ \eng{the man to whom I spoke}.

\textbf{Conclusion.} Pour choisir le pronom relatif, on analyse toujours la \emph{nature} de l'antécédent (personne/chose/lieu) et la \emph{fonction} du pronom dans la relative (sujet, objet, possession).
\end{exoresolu}

\courssub{Fiche méthode 2 : Distinguer relative déterminative et relative explicative}

\begin{methode}[Déterminative ou explicative ?]
\begin{enumerate}
\item \textbf{Poser la question-test} : si je supprime la relative, la phrase garde-t-elle un sens précis ? Si oui $\rightarrow$ explicative (virgules). Si non $\rightarrow$ déterminative (pas de virgule).
\item \textbf{Appliquer les deux règles d'écriture} :
\begin{itemize}
\item Explicative $\rightarrow$ virgule avant et après la relative, \eng{who/which} uniquement (\textbf{jamais} \eng{that}, \textbf{jamais} d'omission).
\item Déterminative $\rightarrow$ aucune virgule, \eng{that} possible, omission possible si le pronom est objet.
\end{itemize}
\item \textbf{Cas particuliers} : les noms propres et les antécédents uniques (\eng{my mother}, \eng{the Earth}, \eng{Mount Cameroon}) appellent presque toujours une explicative : \eng{My father, who is a doctor, works in Douala.}
\item \textbf{Traduire pour vérifier} : \eng{My brother who lives in Bamenda is a teacher} (j'ai plusieurs frères, celui de Bamenda) $\neq$ \eng{My brother, who lives in Bamenda, is a teacher} (je n'ai qu'un frère).
\end{enumerate}
\end{methode}

\begin{exoresolu}[Ponctuation et sens — transformation]
\textbf{Énoncé.} Rewrite each pair of sentences as one sentence containing a relative clause. Add commas where necessary and explain the difference in meaning.

1. My uncle lives in Kumba. He grows cocoa.

2. The students passed the exam. They had revised regularly. (Only those who had revised passed.)

3. Lake Nyos is a crater lake. It is situated in the North-West Region.

\tcblower
\textbf{Solution détaillée.}

\textbf{1.} \eng{My uncle, who lives in Kumba, grows cocoa.}
Justification : « mon oncle » est un antécédent unique et identifié ; l'information sur Kumba est un commentaire supplémentaire $\rightarrow$ relative \emph{explicative} entre virgules, pronom \eng{who} (personne, sujet), \eng{that} interdit.

\textbf{2.} \eng{The students who had revised regularly passed the exam.}
Justification : la consigne précise que seuls les élèves ayant révisé ont réussi ; la relative est indispensable pour identifier \emph{quels} élèves $\rightarrow$ relative \emph{déterminative}, sans virgules. \eng{that} serait aussi correct : \eng{The students that had revised\ldots} Comparaison : \eng{The students, who had revised regularly, passed the exam} signifierait que \emph{tous} les élèves avaient révisé et que tous ont réussi.

\textbf{3.} \eng{Lake Nyos, which is a crater lake, is situated in the North-West Region.}
Justification : \eng{Lake Nyos} est un nom propre, déjà parfaitement identifié $\rightarrow$ explicative entre virgules, pronom \eng{which} (chose, sujet).

\textbf{Conclusion.} La présence ou l'absence de virgules change le sens de la phrase : au Bac, une virgule mal placée dans une relative est sanctionnée comme une faute de grammaire.
\end{exoresolu}

\courssub{Fiche méthode 3 : Construire une noun clause (subordonnée complétive)}

\begin{methode}[Construire une subordonnée complétive (noun clause)]
\begin{enumerate}
\item \textbf{Identifier la fonction} : la noun clause peut être sujet (\eng{What he said is true}), COD (\eng{I know that she is ill}) ou complément de préposition/adjectif (\eng{I am worried about what will happen}).
\item \textbf{Choisir l'introducteur} :
\begin{itemize}
\item Fait affirmé $\rightarrow$ \eng{that} (souvent omis après verbes de parole/pensée : \eng{I think (that)…})
\item Question fermée rapportée $\rightarrow$ \eng{if} / \eng{whether} (\eng{whether} plus soutenu, obligatoire devant \eng{or not})
\item Question ouverte $\rightarrow$ \eng{what, where, when, why, how, who}
\end{itemize}
\item \textbf{Respecter l'ordre affirmatif} : dans une noun clause, \textbf{jamais} d'inversion sujet-verbe, \textbf{jamais} d'auxiliaire \eng{do/does/did}. \eng{I don't know where he lives.} (et non \eng{where does he live}).
\item \textbf{Accorder les temps} (discours indirect / \cle{backshifting}) : après un verbe introducteur au \emph{passé}, recul du temps :
\begin{center}
\begin{tabular}{|l|l|}
\hline
\rowcolor{popblueL} \textbf{Direct} & \textbf{Indirect (après passé)} \\
\hline
Present simple (\eng{go}) & Past simple (\eng{went}) \\
\hline
Present continuous (\eng{is going}) & Past continuous (\eng{was going}) \\
\hline
Present perfect (\eng{has gone}) & Past perfect (\eng{had gone}) \\
\hline
Past simple (\eng{went}) & Past perfect (\eng{had gone}) \\
\hline
\eng{will} & \eng{would} \\
\hline
\eng{can} & \eng{could} \\
\hline
\eng{must} & \eng{had to} \\
\hline
\end{tabular}
\end{center}
\item \textbf{Adapter les repères spatio-temporels} : \eng{now $\rightarrow$ then}, \eng{today $\rightarrow$ that day}, \eng{here $\rightarrow$ there}, \eng{this $\rightarrow$ that}, \eng{these $\rightarrow$ those}, \eng{tomorrow $\rightarrow$ the next/following day}, \eng{next week $\rightarrow$ the following week}, \eng{yesterday $\rightarrow$ the day before / the previous day}.
\item \textbf{Ne pas doubler le sujet} : \eng{What the government has decided is important} — un seul verbe principal pour la phrase entière.
\end{enumerate}
\end{methode}

\begin{exoresolu}[Noun clauses — questions rapportées et discours indirect]
\textbf{Énoncé.} Turn the following direct questions into noun clauses after the given beginnings.

1. “Where does the waste go?” — Nobody knows \ldots

2. “Will the new law protect the forest?” — We wonder \ldots

3. “Why are so many children malnourished?” — The report explains \ldots

4. “What did the doctor recommend?” — Can you tell me \ldots

5. “I have finished the report,” the officer said. — The officer said that \ldots

\tcblower
\textbf{Solution détaillée.}

\textbf{1.} \eng{Nobody knows where the waste goes.}
Raisonnement : question ouverte avec \eng{where} $\rightarrow$ on garde \eng{where}, mais on supprime l'inversion et l'auxiliaire \eng{does} : ordre affirmatif \eng{the waste goes} (le \eng{-s} de la 3\textsuperscript{e} personne réapparaît sur le verbe).

\textbf{2.} \eng{We wonder if (whether) the new law will protect the forest.}
Raisonnement : question fermée (oui/non) sans mot interrogatif $\rightarrow$ on introduit \eng{if} ou \eng{whether} ; ordre affirmatif conservé. \eng{Whether} est plus soutenu et obligatoire devant \eng{or not} : \eng{We wonder whether the law will protect the forest or not.}

\textbf{3.} \eng{The report explains why so many children are malnourished.}
Raisonnement : \eng{why} + ordre affirmatif \eng{so many children are} ; pas d'inversion \eng{are so many children}.

\textbf{4.} \eng{Can you tell me what the doctor recommended?}
Raisonnement : question directe au prétérit avec \eng{did} $\rightarrow$ dans la noun clause, \eng{did} disparaît et le verbe reprend la marque du passé : \eng{recommended}. La phrase principale reste interrogative (\eng{Can you tell me\ldots?}), donc le point d'interrogation final est justifié.

\textbf{5.} \eng{The officer said that he had finished the report.}
Raisonnement : discours indirect $\rightarrow$ \eng{that} (omisible) ; recul des temps : \eng{have finished} (present perfect) $\rightarrow$ \eng{had finished} (past perfect) ; changement de personne : \eng{I} $\rightarrow$ \eng{he}.

\textbf{Conclusion.} Le réflexe capital : dans une noun clause, on retrouve toujours l'ordre \emph{sujet + verbe}, comme dans une phrase affirmative.
\end{exoresolu}

\courssub{Fiche méthode 4 : Employer les adverbial clauses (subordonnées circonstancielles)}

\begin{methode}[Choisir et construire une subordonnée circonstancielle]
\begin{enumerate}
\item \textbf{Identifier la relation logique} et choisir le connecteur :
\begin{center}
\begin{tabular}{|l|l|}
\hline
\rowcolor{popblueL} \textbf{Relation} & \textbf{Connecteurs principaux} \\
\hline
Temps & \eng{when, while, as, as soon as, before, after, until, since} \\
\hline
Cause & \eng{because, since, as} (\eng{since/as} en tête de phrase, plus formel) \\
\hline
But & \eng{so that, in order that} + \eng{can/could/will/would} ; \eng{to / in order to / so as to} + BV (sujet identique) \\
\hline
Conséquence & \eng{so + adj/adv + that}, \eng{such + (a) + nom + that} \\
\hline
Concession & \eng{although, even though, though, while, whereas} \\
\hline
Condition & \eng{if, unless (= if not), provided (that), as long as, on condition that} \\
\hline
\end{tabular}
\end{center}
\item \textbf{Règle d'or du futur (Time/Condition clauses)} : après les conjonctions de \emph{temps} et de \emph{condition}, \textbf{jamais} de \eng{will} / \eng{shall} ; on utilise le \emph{present simple} (ou present perfect pour l'antériorité). \eng{When the rains \textbf{come}, the farmers will plant maize.} (et non \eng{when the rains will come}).
\item \textbf{Distinguer but et conséquence} :
\begin{itemize}
\item But $\rightarrow$ \eng{so that + sujet + can/could/will/would} ou \eng{in order to + BV}. \eng{He spoke slowly so that everyone could understand.}
\item Conséquence $\rightarrow$ \eng{so + adjectif/adverbe + that} ou \eng{such + (a) + nom + that}. \eng{It was so hot that we stopped.}
\end{itemize}
\item \textbf{Gérer la concession} : \eng{although} / \eng{even though} et \eng{but} / \eng{yet} ne se combinent \textbf{jamais} dans la même phrase (calque du français « bien que\ldots mais »). \eng{Although it was raining, the match continued.} OU \eng{It was raining, but the match continued.}
\item \textbf{Ponctuer correctement} : subordonnée en tête de phrase $\rightarrow$ virgule ; subordonnée après la principale $\rightarrow$ généralement pas de virgule (sauf concession parfois pour l'insistance).
\end{enumerate}
\end{methode}

\begin{exoresolu}[Adverbial clauses — choix du connecteur et temps]
\textbf{Énoncé.} Complete the sentences with a suitable conjunction (\eng{because, although, so that, unless, as soon as, when}) and put the verb in brackets in the correct tense.

1. \ldots\ the government (to pass) the new law, illegal logging will decrease.

2. \ldots\ the air in Douala is polluted, many inhabitants suffer from asthma.

3. The nurse explained the treatment slowly \ldots\ the patient could understand.

4. \ldots\ it was very hot, the workers continued planting trees.

5. You will not improve your health \ldots\ you stop smoking.

\tcblower
\textbf{Solution détaillée.}

\textbf{1.} \eng{\textbf{As soon as / When} the government \textbf{passes} the new law, illegal logging will decrease.}
Justification : relation de temps ; règle d'or $\rightarrow$ present simple \eng{passes} dans la subordonnée, \eng{will} uniquement dans la principale. \eng{As soon as} insiste sur l'immédiateté.

\textbf{2.} \eng{\textbf{Because} the air in Douala is polluted, many inhabitants suffer from asthma.}
Justification : relation de cause ; subordonnée en tête $\rightarrow$ virgule obligatoire. On aurait pu dire aussi \eng{Since/As the air is polluted\ldots}

\textbf{3.} \eng{The nurse explained the treatment slowly \textbf{so that} the patient could understand.}
Justification : relation de but avec un sujet différent (\eng{the patient}) $\rightarrow$ \eng{so that + can/could} ; le prétérit \eng{explained} appelle \eng{could}. Si le sujet était identique, on préférerait \eng{in order to} : \eng{She spoke slowly (in order) to be understood.}

\textbf{4.} \eng{\textbf{Although} it was very hot, the workers continued planting trees.}
Justification : concession $\rightarrow$ \eng{although} en tête ; jamais \eng{but} dans la principale (faute classique du francophone).

\textbf{5.} \eng{You will not improve your health \textbf{unless} you stop smoking.}
Justification : \eng{unless} = \eng{if\ldots not} (« à moins que ») ; present simple \eng{stop} après la conjonction de condition ; pas de virgule car la subordonnée suit la principale.

\textbf{Conclusion.} Le choix du connecteur dépend de la relation logique ; la forme du verbe dépend du type de subordonnée (jamais de \eng{will} après \eng{if, when, as soon as, unless}).
\end{exoresolu}

\courssub{Fiche méthode 5 : Transformer et combiner des phrases (entraînement Bac)}

\begin{methode}[Réussir les exercices de transformation]
\begin{enumerate}
\item \textbf{Lire la consigne et le début imposé} : le mot de départ donné (ex. \eng{Despite\ldots}, \eng{He said that\ldots}) détermine la structure attendue.
\item \textbf{Conserver le sens exact} : une transformation ne doit ni ajouter ni retirer d'information.
\item \textbf{Connaître les équivalences-clés} :
\begin{center}
\begin{tabular}{|l|l|}
\hline
\rowcolor{popblueL} \textbf{Structure phrase} & \textbf{Structure nominale (préposition + GN/V-ing)} \\
\hline
\eng{although / even though + sujet + verbe} & \eng{despite / in spite of + nom / V-ing} \\
\hline
\eng{because / since / as + sujet + verbe} & \eng{because of / due to / owing to + nom / V-ing} \\
\hline
\eng{so that + sujet + can/could} & \eng{in order to / so as to + BV} (sujet identique) \\
\hline
\eng{if + sujet + verbe} (condition) & \eng{in case of + nom} (précaution) \\
\hline
Question directe & Noun clause (\eng{that / if / whether / mot interrogatif} + ordre affirmatif) \\
\hline
\end{tabular}
\end{center}
\item \textbf{Vérifier la concordance des temps} après transformation en discours indirect (backshifting).
\item \textbf{Relire à voix haute} : ordre des mots, sujet unique, ponctuation des relatives, accords.
\end{enumerate}
\end{methode}

\begin{exoresolu}[Transformations — entraînement type épreuve]
\textbf{Énoncé.} Rewrite each sentence beginning as shown, without changing the meaning.

1. Although the hospital is far, many patients walk there. — Despite \ldots

2. “I will plant trees next week,” the mayor announced. — The mayor announced that \ldots

3. The water is dirty, so the children fall ill. — Because of \ldots

4. She eats fruit every day. She wants to stay healthy. — She eats fruit every day in order \ldots

5. “Don't drink that water,” the nurse told the child. — The nurse warned the child \ldots

\tcblower
\textbf{Solution détaillée.}

\textbf{1.} \eng{Despite the long distance to the hospital, many patients walk there.} (Ou : \eng{Despite the hospital being far, many patients walk there.})
Justification : \eng{despite} est une \emph{préposition} : elle est suivie d'un groupe nominal (GN) ou d'un gérondif (V-ing), \textbf{jamais} d'une phrase sujet + verbe. On transforme donc la proposition \eng{the hospital is far} en GN \eng{the long distance to the hospital} ou en gérondif \eng{the hospital being far}. Faute à éviter : \eng{Despite the hospital is far\ldots}

\textbf{2.} \eng{The mayor announced that he would plant trees the following week.}
Justification : discours indirect $\rightarrow$ noun clause introduite par \eng{that} ; recul des temps : \eng{will} $\rightarrow$ \eng{would} ; adaptation des repères temporels : \eng{next week} $\rightarrow$ \eng{the following week} ; \eng{I} $\rightarrow$ \eng{he}.

\textbf{3.} \eng{Because of the dirty water, the children fall ill.}
Justification : \eng{because + phrase} devient \eng{because of + groupe nominal} ; le verbe \eng{is dirty} se nominalise en \eng{the dirty water}. \eng{Due to / Owing to the dirty water} sont aussi corrects (plus formels).

\textbf{4.} \eng{She eats fruit every day in order to stay healthy.}
Justification : même sujet dans les deux phrases $\rightarrow$ \eng{in order to + base verbale} (ou \eng{so as to}) remplace \eng{so that she can stay healthy}.

\textbf{5.} \eng{The nurse warned the child not to drink that water.}
Justification : ordre/négation au discours indirect $\rightarrow$ \eng{warn/tell/ask + objet + not to + BV}. \eng{That} devient \eng{that} (démonstratif) ou \eng{the} selon le contexte.

\textbf{Conclusion.} Ces transformations (concession, discours indirect, cause, but, ordre) reviennent chaque année à l'épreuve : maîtriser les équivalences phrase $\leftrightarrow$ groupe nominal est la clé.
\end{exoresolu}

\courssub{Synthèse : Tableau récapitulatif des trois types de clauses}

\begin{center}
\begin{tabularx}{\linewidth}{|X|X|X|X|}
\hline
\rowcolor{popblueL} \textbf{Critère} & \textbf{Relative Clause} & \textbf{Noun Clause} & \textbf{Adverbial Clause} \\
\hline
\textbf{Fonction} & Modifie un nom (antécédent) & Fonctionne comme un nom (sujet, COD, compl. prép.) & Modifie le verbe principal (circonstance) \\
\hline
\textbf{Question} & Which one? What kind? & What? Who? & When? Why? How? Under what condition? \\
\hline
\textbf{Introducteurs} & \eng{who, which, that, whose, where, when} & \eng{that, if, whether, what, where, when, why, how, who} & \eng{when, because, although, if, so that, unless\ldots} \\
\hline
\textbf{Omission possible ?} & Oui, si pronom = objet (rel. déterminative) & Oui, \eng{that} après verbes de parole/pensée & Non (sauf \eng{while/when + V-ing} réduction) \\
\hline
\textbf{Ponctuation} & Virgules si explicative & Jamais de virgule après \eng{that} & Virgule si en tête de phrase \\
\hline
\textbf{Règle de temps} & Accord avec l'antécédent & Backshifting si verbe introducteur au passé & Present simple après \eng{if/when/unless\ldots} (pas de \eng{will}) \\
\hline
\textbf{Exemple} & \eng{The doctor who treated me is kind.} & \eng{I know that he is kind.} & \eng{Because he is kind, I trust him.} \\
\hline
\end{tabularx}
\end{center}

\begin{attention}[Les 5 erreurs qui coûtent des points au Bac]
\begin{enumerate}
\item \textbf{L'inversion dans les noun clauses} : \eng{I don't know where does he live.} \textbf{Faux !} Ordre affirmatif obligatoire : \eng{I don't know where he lives.} C'est l'erreur la plus fréquente des candidats francophones.
\item \textbf{\eng{Will} après \eng{if} ou \eng{when}} : \eng{When the rain will stop, we will go out.} \textbf{Faux !} Present simple dans la subordonnée : \eng{When the rain stops, we will go out.}
\item \textbf{\eng{Although\ldots but} dans la même phrase} : calque du français « bien que\ldots mais ». On choisit l'un ou l'autre : \eng{Although he is poor, he is happy} \textbf{ou} \eng{He is poor, but he is happy.}
\item \textbf{La virgule et \eng{that} dans les explicatives} : \eng{My father, that is a farmer, lives in Bafoussam.} Double faute : après une virgule, jamais \eng{that} $\rightarrow$ \eng{My father, who is a farmer, lives in Bafoussam.}
\item \textbf{Les faux amis et calques lexicaux (prépositions)} :
\begin{itemize}
\item \eng{actually} = « en fait » (pas « actuellement » $\rightarrow$ \eng{currently})
\item \eng{demand} = « exiger » (pas « demander » $\rightarrow$ \eng{ask})
\item \eng{eventually} = « finalement » (pas « éventuellement » $\rightarrow$ \eng{possibly})
\item \eng{sensible} = « raisonnable » (pas « sensible » $\rightarrow$ \eng{sensitive})
\item Prépositions figées : \eng{interested in}, \eng{listen to}, \eng{depend on}, \eng{good at}, \eng{afraid of}, \eng{capable of}, \eng{responsible for} — la préposition anglaise ne se devine jamais à partir du français.
\end{itemize}
\end{enumerate}
\end{attention}

\courssub{Exercices d'entraînement (Bac blanc)}

\begin{exercice}[Exercice 1 : Relative clauses — Choix et ponctuation]
Combine each pair of sentences using a relative clause. Decide whether the clause is defining or non-defining and punctuate accordingly.
\begin{enumerate}
\item The cocoa farmers protested. Their crops were destroyed by pests.
\item Mount Cameroon is an active volcano. It is the highest peak in West Africa.
\item The doctor prescribed a new medicine. It has no side effects.
\item My sister lives in Bamenda. She is a nurse.
\item The students cleaned the campus. They belong to the environment club.
\end{enumerate}
\end{exercice}

\begin{exercice}[Exercice 2 : Noun clauses — Discours indirect et questions rapportées]
Rewrite the following sentences as reported speech or indirect questions.
\begin{enumerate}
\item “The water in this well is not safe,” the expert said.
\item “Have you ever seen a gorilla in the wild?” the tourist asked the guide.
\item “Why did the bridge collapse?” the journalist wanted to know.
\item “I will submit the report tomorrow,” the officer promised.
\item “Don't forget to lock the door,” the manager told the guard.
\end{enumerate}
\end{exercice}

\begin{exercice}[Exercice 3 : Adverbial clauses — Connecteurs et temps]
Complete the sentences with the correct conjunction and verb form.
\begin{enumerate}
\item \ldots\ (you / finish) your homework, you can watch TV. (condition)
\item The match was cancelled \ldots\ (it / rain) heavily. (cause)
\item \ldots\ (he / be) tired, he continued working. (concession)
\item We will leave \ldots\ (the sun / rise). (temps)
\item She spoke very clearly \ldots\ (everyone / understand). (but)
\end{enumerate}
\end{exercice}

\begin{exercice}[Exercice 4 : Transformations mixtes]
Rewrite each sentence beginning as shown.
\begin{enumerate}
\item Because he was ill, he missed the exam. — Due to \ldots
\item “Where is the nearest health centre?” the stranger asked. — The stranger asked \ldots
\item Although the road is bad, trucks use it daily. — Despite \ldots
\item The farmer uses organic fertilizer. He wants to protect the soil. — The farmer uses organic fertilizer in order \ldots
\item “I have never been to Kribi,” my friend said. — My friend said that \ldots
\end{enumerate}
\end{exercice}

\courssub{Corrigés des exercices d'entraînement}

\begin{corrige}[Exercice 1]
\begin{enumerate}
\item \eng{The cocoa farmers whose crops were destroyed by pests protested.} (Déterminative : identifie quels fermiers. Pas de virgule. \eng{whose} = possession.)
\item \eng{Mount Cameroon, which is an active volcano, is the highest peak in West Africa.} (Explicative : nom propre déjà identifié. Virgules obligatoires. \eng{which} sujet, \eng{that} interdit.)
\item \eng{The doctor prescribed a new medicine which has no side effects.} (Déterminative : précise quel médicament. \eng{which} sujet $\rightarrow$ pas d'omission. \eng{that} possible.)
\item \eng{My sister, who lives in Bamenda, is a nurse.} (Explicative : antécédent unique « my sister ». Virgules. \eng{who} sujet.)
\item \eng{The students who belong to the environment club cleaned the campus.} (Déterminative : identifie quels étudiants. Pas de virgule. \eng{who} sujet. \eng{that} possible.)
\end{enumerate}
\end{corrige}

\begin{corrige}[Exercice 2]
\begin{enumerate}
\item \eng{The expert said that the water in that well was not safe.} (Backshifting : \eng{is} $\rightarrow$ \eng{was} ; \eng{this} $\rightarrow$ \eng{that}.)
\item \eng{The tourist asked the guide if/whether he had ever seen a gorilla in the wild.} (Question fermée $\rightarrow$ \eng{if/whether} ; backshifting : \eng{have seen} $\rightarrow$ \eng{had seen} ; \eng{you} $\rightarrow$ \eng{he}.)
\item \eng{The journalist wanted to know why the bridge had collapsed.} (Question ouverte \eng{why} ; ordre affirmatif ; backshifting : \eng{did collapse} $\rightarrow$ \eng{had collapsed}.)
\item \eng{The officer promised that he would submit the report the next day.} (\eng{will} $\rightarrow$ \eng{would} ; \eng{tomorrow} $\rightarrow$ \eng{the next day} / \eng{the following day}.)
\item \eng{The manager told the guard not to forget to lock the door.} (Ordre négatif $\rightarrow$ \eng{tell + objet + not to + BV}.)
\end{enumerate}
\end{corrige}

\begin{corrige}[Exercice 3]
\begin{enumerate}
\item \eng{If / Provided (that) / As long as you finish your homework, you can watch TV.} (Condition $\rightarrow$ present simple \eng{finish}.)
\item \eng{The match was cancelled because / since / as it was raining heavily.} (Cause ; \eng{was raining} = past continuous pour action en cours.)
\item \eng{Although / Even though he was tired, he continued working.} (Concession ; pas de \eng{but} dans la principale.)
\item \eng{We will leave when / as soon as the sun rises.} (Temps $\rightarrow$ present simple \eng{rises} ; pas de \eng{will rise}.)
\item \eng{She spoke very clearly so that everyone could understand.} (But, sujet différent $\rightarrow$ \eng{so that + could} ; \eng{spoke} au passé appelle \eng{could}.)
\end{enumerate}
\end{corrige}

\begin{corrige}[Exercice 4]
\begin{enumerate}
\item \eng{Due to his illness, he missed the exam.} (Ou \eng{Due to the fact that he was ill\ldots} — mais la forme nominale simple est attendue : \eng{his illness}.)
\item \eng{The stranger asked where the nearest health centre was.} (Noun clause : \eng{where} + ordre affirmatif \eng{the nearest health centre was} ; backshifting \eng{is} $\rightarrow$ \eng{was}.)
\item \eng{Despite the bad road, trucks use it daily.} (Ou \eng{Despite the road being bad\ldots} ; \eng{Despite} + GN/V-ing.)
\item \eng{The farmer uses organic fertilizer in order to protect the soil.} (Même sujet $\rightarrow$ \eng{in order to + BV}.)
\item \eng{My friend said that he had never been to Kribi.} (Backshifting : \eng{have never been} $\rightarrow$ \eng{had never been} ; \eng{I} $\rightarrow$ \eng{he}.)
\end{enumerate}
\end{corrige}

\begin{savaistu}[Le Cameroun et la langue anglaise : contexte bilingue]
Le Cameroun est le seul pays d'Afrique (avec le Canada) à avoir deux langues officielles \emph{de droit} : le français et l'anglais. L'anglais y est parlé dans les deux régions du Nord-Ouest et du Sud-Ouest (ancien Southern Cameroons britannique). Au Bac, les sujets puisent souvent dans ce contexte : environnement (forêt du bassin du Congo, lac Nyos, Mont Cameroun), santé (paludisme, nutrition, hôpitaux régionaux), agriculture (cacao, café, banane). Maîtriser les clauses permet de rédiger des compositions structurées sur ces thèmes : \eng{Although Cameroon has rich biodiversity, deforestation threatens it.} / \eng{The government should ensure that communities whose land is affected are compensated.}
\end{savaistu}

\begin{aretenir}[Check-list finale avant l'épreuve]
\begin{itemize}
\item \textbf{Relative} : Antécédent (qui/quoi ?) $\rightarrow$ Fonction (sujet/objet/possessif/lieu/temps) $\rightarrow$ Pronom (\eng{who/which/whose/where/when/that/Ø}) $\rightarrow$ Virgules ? (Explicative = oui, \eng{that} non).
\item \textbf{Noun clause} : Introducteur (\eng{that/if/whether/wh-}) $\rightarrow$ \textbf{Ordre affirmatif} (sujet + verbe) $\rightarrow$ Backshifting si passé $\rightarrow$ Adaptation repères.
\item \textbf{Adverbial} : Connecteur logique $\rightarrow$ Temps du verbe (présent après \eng{if/when/unless}) $\rightarrow$ Pas \eng{although\ldots but} $\rightarrow$ Virgule si en tête.
\item \textbf{Transformation} : Préposition (\eng{despite/in spite of/because of}) + GN/V-ing $\leftrightarrow$ Conjonction + phrase.
\end{itemize}
\end{aretenir}

\newpage

\coursec{Exercices}

\courssub{Je vérifie mes connaissances}

\begin{exercice}[QCM : choose the correct answer]
Choose the correct answer (a, b or c) to complete each sentence.
\begin{enumerate}
\item I don't know where \ldots{} \quad a) is the market \quad b) the market is \quad c) does the market be
\item The nurse \ldots{} works at the Buea Regional Hospital is very kind. \quad a) which \quad b) who \quad c) whose
\item \ldots{} it rains heavily, the match will be postponed. \quad a) If \quad b) Unless \quad c) Although
\item She asked me \ldots{} I had taken my medicine. \quad a) that \quad b) if \quad c) what
\item This is the village \ldots{} my grandfather was born. \quad a) which \quad b) who \quad c) where
\item He went to school \ldots{} he was feeling sick. \quad a) because \quad b) although \quad c) so that
\item The fact \ldots{} water is polluted worries everybody. \quad a) that \quad b) who \quad c) where
\item Do you know \ldots{} the hospital closes on Sundays? \quad a) when does \quad b) when \quad c) what time does
\end{enumerate}
\end{exercice}

\begin{exercice}[True or false — justify your answer]
Say whether each statement is \textbf{true} or \textbf{false}, and justify your answer in one sentence.
\begin{enumerate}
\item In a defining relative clause, the relative pronoun can never be omitted.
\item In English, we say \eng{Although she was tired, she continued working}, without \eng{but}.
\item After \eng{if} in a conditional sentence, we use \eng{will} to talk about the future.
\item \eng{Whose} is used to express possession.
\item In an indirect question, the word order is the same as in a direct question.
\end{enumerate}
\end{exercice}

\begin{exercice}[Vocabulary : match the clause to its type]
Match each underlined clause (1--5) to its type (a--c). One type may be used more than once.
\begin{enumerate}
\item The doctor \textbf{who vaccinated the children} comes from Bamenda.
\item \textbf{Because the river is polluted}, the villagers boil their drinking water.
\item I believe \textbf{that regular exercise keeps us healthy}.
\item My uncle, \textbf{who lives in Douala}, is a pharmacist.
\item She wears a mask \textbf{so that she does not breathe the dust}.
\end{enumerate}
\begin{center}
a) relative clause \qquad b) noun clause \qquad c) adverbial clause
\end{center}
\end{exercice}

\begin{exercice}[Complete the sentences]
Complete each sentence with the correct word from the box. Use each word only once.
\begin{center}
\begin{tabular}{|l|l|l|l|l|}
\hline
\rowcolor{popblueL} who & whose & unless & because & that \\
\hline
\end{tabular}
\end{center}
\begin{enumerate}
\item The girl \ldots{} father is a surgeon wants to study medicine.
\item People fall sick \ldots{} they drink untreated water.
\item You will not stay healthy \ldots{} you eat a balanced diet.
\item The man \ldots{} spoke to us is an environmental inspector.
\item Everybody knows \ldots{} smoking is dangerous.
\end{enumerate}
\end{exercice}

\courssub{Je m'entraîne}

\begin{exercice}[Fill in the blanks]
Fill in the blanks with \eng{who}, \eng{which}, \eng{whose}, \eng{where} or \eng{that}. (Theme : health and environment in Cameroon.)
\begin{enumerate}
\item The hospital \ldots{} was built last year has a modern laboratory.
\item Farmers \ldots{} use too many chemicals pollute the soil.
\item This is the market \ldots{} my mother buys fresh vegetables.
\item The students \ldots{} project won the prize cleaned the school yard.
\item Malaria is a disease \ldots{} kills many children in Africa.
\item The woman \ldots{} sells fruit near the post office is my aunt.
\item Yaoundé is a city \ldots{} traffic jams cause a lot of pollution.
\item The trees \ldots{} we planted last year are growing well.
\item The boy \ldots{} bicycle was stolen reported the case to the police.
\item Clean water, \ldots{} is essential for life, is scarce in some villages.
\end{enumerate}
\end{exercice}

\begin{exercice}[Combine the sentences]
Combine each pair of sentences using a suitable relative pronoun or conjunction, as in the example.\\[4pt]
\emph{Example :} The man is a doctor. He treated my mother. $\rightarrow$ \eng{The man who treated my mother is a doctor.}
\begin{enumerate}
\item The nurse is very patient. She looks after my grandmother.
\item The factory pollutes the Wouri river. It was built in 2010.
\item I bought a water filter. It was not expensive.
\item The girl won the health quiz. Her mother teaches biology.
\item We visited a health centre. It is situated near the stadium.
\item He did not sleep well. He had a terrible headache.
\end{enumerate}
\end{exercice}

\begin{exercice}[Correct the mistakes]
The following sentences were written by francophone learners. Each sentence contains \textbf{one mistake}. Correct it.
\begin{enumerate}
\item If it will rain, we will stay at home.
\item The boy who he came late is my friend.
\item Although she is tired, but she continues to work.
\item I don't know where does she live.
\item The house where I was born there is very old.
\item Because he was ill, so he did not come to school.
\item This is the man which sold me the medicine.
\end{enumerate}
\end{exercice}

\begin{exercice}[Rewrite the sentences]
Rewrite each sentence beginning as shown, keeping the same meaning.
\begin{enumerate}
\item Despite the rain, they went out. $\rightarrow$ Although \ldots
\item She was tired, so she went to bed early. $\rightarrow$ Because \ldots
\item He saved money in order to buy a bicycle. $\rightarrow$ He saved money so that \ldots
\item If you don't wash your hands, you will catch diseases. $\rightarrow$ Unless \ldots
\item “Where is the clinic?” the tourist asked. $\rightarrow$ The tourist asked \ldots
\item It was such a hot day that we stayed indoors. $\rightarrow$ The day was so \ldots
\end{enumerate}
\end{exercice}

\courssub{Je me prépare au Bac}

\begin{exercice}[Reading comprehension and language work]
Read the following text carefully and answer the questions that follow.

\begin{textebox}[Keeping Buea clean]
Every Saturday morning, the young people of Buea, who belong to an association called Green Mountain, meet at the market square. They collect the plastic bags and bottles which traders have thrown on the ground during the week. Because the town lies on the slopes of Mount Cameroon, the rain carries this rubbish into the streams that supply the villages below. “If we do not act now, our children will drink polluted water,” says Miriam, whose family has lived in Buea for three generations. Although the work is hard and the volunteers receive no money, they believe that a clean town is a healthy town. The mayor, who visited them last month, promised that the council would provide gloves and wheelbarrows so that the work could be done more safely. Many residents say that they do not know where they should put their waste, so Green Mountain also teaches families how to sort rubbish before it is collected.
\end{textebox}

\textbf{A. Comprehension questions.} Answer in complete sentences.
\begin{enumerate}
\item What do the members of Green Mountain do every Saturday?
\item Why is the rubbish dangerous for the villages below the town?
\item What did the mayor promise the volunteers?
\item According to the text, why does Green Mountain also teach families?
\end{enumerate}

\textbf{B. Language work.}
\begin{enumerate}
\item Find in the text : (a) one defining relative clause, (b) one non-defining relative clause, (c) one noun clause, (d) two adverbial clauses. Copy them and name the connector used.
\item Rewrite : “The work is hard. The volunteers continue.” (Use \eng{although}.)
\item Rewrite : “Where should we put our waste?” the residents asked. (Use an indirect question.)
\end{enumerate}
\end{exercice}

\begin{exercice}[Translation]
Traduisez le texte suivant en anglais.\\[4pt]
« L'hôpital qui a été construit dans notre village l'année dernière aide beaucoup de gens. Parce que les routes sont mauvaises, les malades arrivaient autrefois trop tard à la ville. Le médecin, dont la femme est infirmière, dit que la santé de la population s'améliore. Si le gouvernement construit plus de centres de santé, moins d'enfants mourront de maladies qu'on peut soigner. »
\end{exercice}

\begin{exercice}[Guided writing : composition]
Your school magazine has asked students to write about health and the environment. Write a composition of \textbf{150 to 180 words} entitled :\\[4pt]
\begin{center}
\textbf{“How to keep our town clean and healthy”}
\end{center}
Your composition must contain :
\begin{itemize}
\item at least \textbf{one relative clause} (with \eng{who}, \eng{which}, \eng{whose} or \eng{where}) ;
\item at least \textbf{one noun clause} (after \eng{I think that}, \eng{everybody knows that}, etc.) ;
\item at least \textbf{one adverbial clause} (with \eng{because}, \eng{if}, \eng{although}, \eng{so that}, etc.) ;
\item a clear introduction, two developed ideas and a short conclusion.
\end{itemize}
\textbf{Self-assessment grid} — tick when finished :
\begin{center}
\begin{tabular}{|l|l|}
\hline
\rowcolor{popblueL} Criteria & Done \\
\hline
I used at least three different types of clauses. & $\square$ \\
\hline
My relative pronouns are correct (no \eng{which} for people). & $\square$ \\
\hline
No \eng{but} after \eng{although}, no \eng{will} after \eng{if}. & $\square$ \\
\hline
My word order in indirect questions is correct. & $\square$ \\
\hline
My text has between 150 and 180 words. & $\square$ \\
\hline
\end{tabular}
\end{center}
\end{exercice}

\newpage

\coursec{Corrigés des exercices}

\begin{corrige}[Exercice 1 — QCM : choose the correct answer]
\begin{enumerate}
\item \textbf{b) the market is} — \eng{I don't know where the market is.} Dans une question indirecte, on garde l'ordre sujet + verbe (pas d'inversion). « Je ne sais pas où est le marché. »
\item \textbf{b) who} — \eng{The nurse who works at the Buea Regional Hospital is very kind.} \eng{Who} pour une personne (\eng{which} est réservé aux choses et aux animaux).
\item \textbf{a) If} — \eng{If it rains heavily, the match will be postponed.} Condition logique : « Si ». \eng{Unless} (« sauf si ») donnerait un sens absurde ici.
\item \textbf{b) if} — \eng{She asked me if I had taken my medicine.} Question indirecte de type oui/non : on introduit par \eng{if} (« si »).
\item \textbf{c) where} — \eng{This is the village where my grandfather was born.} \eng{Where} pour un lieu.
\item \textbf{b) although} — \eng{He went to school although he was feeling sick.} Concession : « bien qu'il se sente malade ». \eng{Because} et \eng{so that} donneraient un sens contraire au contexte.
\item \textbf{a) that} — \eng{The fact that water is polluted worries everybody.} Subordonnée complétive (noun clause) introduite par \eng{that}.
\item \textbf{b) when} — \eng{Do you know when the hospital closes on Sundays?} Question indirecte : ordre sujet + verbe, sans \eng{does}.
\end{enumerate}
\end{corrige}

\begin{corrige}[Exercice 2 — True or false — justify your answer]
\begin{enumerate}
\item \textbf{False.} The relative pronoun can be omitted when it is the \emph{object} of the defining clause : \eng{The book (that) I bought is interesting.}
\item \textbf{True.} English never uses \eng{although} and \eng{but} together in the same sentence : \eng{Although she was tired, she continued working.}
\item \textbf{False.} After \eng{if} we use the present simple, not \eng{will} : \eng{If it rains, we will stay at home.}
\item \textbf{True.} \eng{Whose} expresses possession : \eng{The girl whose father is a surgeon wants to study medicine.}
\item \textbf{False.} In an indirect question the word order is subject + verb, without inversion or \eng{do/does} : \eng{Do you know where she lives?}
\end{enumerate}
\end{corrige}

\begin{corrige}[Exercice 3 — Vocabulary : match the clause to its type]
\begin{enumerate}
\item \textbf{a) relative clause} — \eng{who vaccinated the children} qualifie le nom \eng{the doctor} (relative déterminative, pronom \eng{who}).
\item \textbf{c) adverbial clause} — \eng{Because the river is polluted} exprime la cause (conjonction \eng{because}).
\item \textbf{b) noun clause} — \eng{that regular exercise keeps us healthy} est complément d'objet du verbe \eng{believe} (introduite par \eng{that}).
\item \textbf{a) relative clause} — \eng{who lives in Douala} : relative explicative (entre virgules) qui apporte une information supplémentaire sur \eng{my uncle}.
\item \textbf{c) adverbial clause} — \eng{so that she does not breathe the dust} exprime le but (conjonction \eng{so that}).
\end{enumerate}
\textbf{Astuce de reconnaissance :} une relative suit toujours un nom ; une noun clause joue le rôle de sujet ou de complément d'un verbe ; une adverbial clause répond aux questions \emph{quand ? pourquoi ? dans quel but ? à quelle condition ?}
\end{corrige}

\begin{corrige}[Exercice 4 — Complete the sentences]
\begin{enumerate}
\item \textbf{whose} — \eng{The girl whose father is a surgeon wants to study medicine.} Possession : « la fille dont le père est chirurgien ».
\item \textbf{because} — \eng{People fall sick because they drink untreated water.} Cause : « parce qu'ils boivent de l'eau non traitée ».
\item \textbf{unless} — \eng{You will not stay healthy unless you eat a balanced diet.} Condition négative : « à moins que tu ne manges équilibré ».
\item \textbf{who} — \eng{The man who spoke to us is an environmental inspector.} Personne, sujet du verbe de la relative.
\item \textbf{that} — \eng{Everybody knows that smoking is dangerous.} Noun clause complément de \eng{knows}.
\end{enumerate}
\end{corrige}

\begin{corrige}[Exercice 5 — Fill in the blanks]
\begin{enumerate}
\item \textbf{which/that} — \eng{The hospital which was built last year has a modern laboratory.} (chose, sujet)
\item \textbf{who/that} — \eng{Farmers who use too many chemicals pollute the soil.} (personnes, sujet)
\item \textbf{where} — \eng{This is the market where my mother buys fresh vegetables.} (lieu)
\item \textbf{whose} — \eng{The students whose project won the prize cleaned the school yard.} (possession)
\item \textbf{which/that} — \eng{Malaria is a disease which kills many children in Africa.} (chose, sujet)
\item \textbf{who/that} — \eng{The woman who sells fruit near the post office is my aunt.} (personne, sujet)
\item \textbf{where} — \eng{Yaoundé is a city where traffic jams cause a lot of pollution.} (lieu)
\item \textbf{which/that} (ou omis) — \eng{The trees (which) we planted last year are growing well.} (chose, complément d'objet : le pronom peut être omis)
\item \textbf{whose} — \eng{The boy whose bicycle was stolen reported the case to the police.} (possession)
\item \textbf{which} — \eng{Clean water, which is essential for life, is scarce in some villages.} (relative explicative : seul \eng{which} est possible, jamais \eng{that} après une virgule)
\end{enumerate}
\end{corrige}

\begin{corrige}[Exercice 6 — Combine the sentences]
\begin{enumerate}
\item \eng{The nurse who looks after my grandmother is very patient.} (« L'infirmière qui s'occupe de ma grand-mère est très patiente. »)
\item \eng{The factory which was built in 2010 pollutes the Wouri river.} (chose, sujet : \eng{which/that})
\item \eng{I bought a water filter which was not expensive.} ou \eng{The water filter (which) I bought was not expensive.}
\item \eng{The girl whose mother teaches biology won the health quiz.} (possession : \eng{whose} = « dont »)
\item \eng{We visited a health centre which is situated near the stadium.}
\item \eng{He did not sleep well because he had a terrible headache.} (subordonnée de cause avec \eng{because})
\end{enumerate}
\textbf{Point de vigilance :} ne jamais répéter le sujet dans la relative — pas \eng{The nurse who she looks after\ldots} ; le pronom relatif remplace le nom.
\end{corrige}

\begin{corrige}[Exercice 7 — Correct the mistakes]
\begin{enumerate}
\item \eng{If it will rain} $\rightarrow$ \eng{If it \textbf{rains}, we will stay at home.} Après \eng{if}, le futur se traduit par le présent simple en anglais.
\item \eng{The boy who he came late} $\rightarrow$ \eng{The boy \textbf{who} came late is my friend.} On supprime \eng{he} : le pronom relatif est déjà le sujet.
\item \eng{Although she is tired, but\ldots} $\rightarrow$ \eng{Although she is tired, she continues to work.} On supprime \eng{but} : \eng{although} et \eng{but} ne se combinent jamais.
\item \eng{where does she live} $\rightarrow$ \eng{I don't know where \textbf{she lives}.} Question indirecte : ordre sujet + verbe, sans auxiliaire.
\item \eng{The house where I was born there} $\rightarrow$ \eng{The house \textbf{where} I was born is very old.} On supprime \eng{there} : \eng{where} exprime déjà le lieu.
\item \eng{Because he was ill, so\ldots} $\rightarrow$ \eng{Because he was ill, he did not come to school.} On supprime \eng{so} : \eng{because} et \eng{so} ne s'emploient pas ensemble.
\item \eng{the man which} $\rightarrow$ \eng{This is the man \textbf{who} sold me the medicine.} \eng{Who} pour une personne.
\end{enumerate}
\textbf{Remarque :} toutes ces erreurs viennent du calque du français (« bien que\ldots mais », « parce que\ldots donc », « l'homme que »). Traduisez la structure, pas les mots.
\end{corrige}

\begin{corrige}[Exercice 8 — Rewrite the sentences]
\begin{enumerate}
\item \eng{Although it rained, they went out.} (ou \eng{Although it was raining, they went out.}) — concession, sans \eng{but}.
\item \eng{Because she was tired, she went to bed early.} — cause ; attention à ne pas ajouter \eng{so}.
\item \eng{He saved money so that he could buy a bicycle.} — but ; \eng{so that} + \eng{could} au passé.
\item \eng{Unless you wash your hands, you will catch diseases.} — \eng{unless} = \eng{if\ldots not} ; on supprime la négation.
\item \eng{The tourist asked where the clinic was.} — question indirecte : ordre sujet + verbe, \eng{is} devient \eng{was} (concordance des temps).
\item \eng{The day was so hot that we stayed indoors.} — \eng{so} + adjectif + \eng{that}.
\end{enumerate}
\end{corrige}

\begin{corrige}[Exercice 9 — Reading comprehension and language work]
\textbf{A. Comprehension questions} (réponses modèles ; toute formulation correcte est acceptée).
\begin{enumerate}
\item \eng{Every Saturday, the members of Green Mountain meet at the market square and collect the plastic bags and bottles which traders have thrown on the ground during the week.}
\item \eng{The rubbish is dangerous because the rain carries it into the streams that supply the villages below, so the villagers may drink polluted water.}
\item \eng{The mayor promised that the council would provide gloves and wheelbarrows so that the work could be done more safely.}
\item \eng{Green Mountain also teaches families because many residents do not know where they should put their waste, so the association shows them how to sort rubbish before it is collected.}
\end{enumerate}
\textbf{B. Language work.}
\begin{enumerate}
\item Exemples attendus (plusieurs réponses possibles) :
\begin{itemize}
\item (a) Defining relative clause : \eng{which traders have thrown on the ground} (connector : \eng{which}) — ou \eng{that supply the villages below} (\eng{that}).
\item (b) Non-defining relative clause : \eng{whose family has lived in Buea for three generations} (connector : \eng{whose}) — ou \eng{who visited them last month} (\eng{who}).
\item (c) Noun clause : \eng{that a clean town is a healthy town} (connector : \eng{that}) — ou \eng{where they should put their waste} (\eng{where}).
\item (d) Adverbial clauses : \eng{Because the town lies on the slopes of Mount Cameroon} (\eng{because}, cause) ; \eng{Although the work is hard} (\eng{although}, concession) ; \eng{so that the work could be done more safely} (\eng{so that}, but) ; \eng{before it is collected} (\eng{before}, temps).
\end{itemize}
\item \eng{Although the work is hard, the volunteers continue.}
\item \eng{The residents asked where they should put their waste.} (question indirecte : sujet + verbe, pas de point d'interrogation)
\end{enumerate}
\textbf{Barème indicatif (sur 20) :} A : 4 questions $\times$ 2 pts = 8 pts (1 pt sens, 1 pt langue) ; B.1 : 5 clauses $\times$ 1 pt = 5 pts ; B.2 : 2 pts ; B.3 : 2 pts ; présentation générale : 3 pts.
\textbf{Points de vigilance :} recopier la clause entière, pas seulement le connecteur ; dans B.3, ne pas écrire \eng{where should they put} ; ne pas mettre de point d'interrogation à la question indirecte.
\end{corrige}

\begin{corrige}[Exercice 10 — Translation]
\textbf{Traduction de référence :}\\[4pt]
\eng{The hospital which was built in our village last year helps many people. Because the roads are bad, sick people used to arrive in town too late. The doctor, whose wife is a nurse, says that the health of the population is improving. If the government builds more health centres, fewer children will die of diseases which can be treated.}\\[4pt]
\textbf{Points de vigilance :}
\begin{itemize}
\item « qui a été construit » $\rightarrow$ \eng{which was built} (passif : \eng{be} + participe passé).
\item « arrivaient autrefois » $\rightarrow$ \eng{used to arrive} (habitude passée révolue).
\item « dont la femme » $\rightarrow$ \eng{whose wife} ; relative explicative entre virgules.
\item « s'améliore » $\rightarrow$ \eng{is improving} (present continuous pour un changement en cours).
\item « Si le gouvernement construit\ldots mourront » $\rightarrow$ \eng{If the government builds\ldots will die} : présent après \eng{if}, futur dans la principale.
\item « moins d'enfants » $\rightarrow$ \eng{fewer children} (dénombrable), pas \eng{less children}.
\item « qu'on peut soigner » $\rightarrow$ \eng{which can be treated} (passif).
\end{itemize}
\textbf{Barème indicatif (sur 10) :} 5 segments $\times$ 2 pts (1 pt lexique, 1 pt grammaire) ; toute erreur de pronom relatif ou de temps coûte 0,5 pt.
\end{corrige}

\begin{corrige}[Exercice 11 — Guided writing : composition]
\textbf{Proposition de rédaction modèle (environ 165 mots) :}\\[4pt]
\eng{Our town is our home, and everybody knows that a dirty town is a sick town. Keeping it clean and healthy is the duty of every citizen.}\\[4pt]
\eng{First of all, we must manage our waste properly. Many diseases spread because people throw plastic bags and bottles into the streets and gutters. If every family sorted its rubbish and used dustbins, the town would be much cleaner. The council should also provide more bins so that residents can dispose of their waste easily.}\\[4pt]
\eng{Secondly, we must protect our water and our air. The streams which supply our neighbourhoods must not be used as dumping grounds. Although some people think that one person cannot change anything, I believe that small actions, when they are repeated by thousands of citizens, make a big difference. Young people who join clean-up campaigns, like the Green Mountain volunteers in Buea, show us the way.}\\[4pt]
\eng{To conclude, a clean town is a healthy town. If we all act together, our children will grow up in a safer and more beautiful environment.}\\[4pt]
\textbf{Analyse du modèle :}
\begin{itemize}
\item Noun clauses : \eng{everybody knows that a dirty town is a sick town} ; \eng{I believe that small actions\ldots make a big difference}.
\item Relative clauses : \eng{The streams which supply our neighbourhoods} ; \eng{Young people who join clean-up campaigns}.
\item Adverbial clauses : \eng{because people throw plastic bags\ldots} (cause) ; \eng{If every family sorted its rubbish\ldots} (condition) ; \eng{so that residents can dispose of their waste} (but) ; \eng{Although some people think\ldots} (concession).
\end{itemize}
\textbf{Barème indicatif (sur 20) :} respect du sujet et de la longueur : 4 pts ; richesse et correction des clauses exigées : 6 pts ; grammaire générale (temps, accords, ordre des mots) : 5 pts ; vocabulaire thématique : 3 pts ; organisation (introduction, deux idées, conclusion, connecteurs) : 2 pts.
\textbf{Points de vigilance :} vérifier qu'il n'y a ni \eng{but} après \eng{although}, ni \eng{will} après \eng{if} ; employer \eng{who} pour les personnes ; compter ses mots (150--180) ; relire pour éliminer les calques du français (\eng{the people} pour « les gens » en général se dit simplement \eng{people}).
\end{corrige}

\newpage

\coursec{Fiche bilan}

\begin{popfigure}
\begin{tikzpicture}[mindmap, grow cyclic, every node/.style={concept, font=\small},
root concept/.append style={concept color=popblue, font=\bfseries},
level 1/.append style={level distance=3.4cm, sibling angle=60},
level 2/.append style={level distance=2.2cm, sibling angle=45, font=\scriptsize}]
\node[root concept]{CLAUSES}
  child[concept color=poporange]{ node{Relative clauses}
    child{ node{who, which, that} }
    child{ node{whose, where, when} }
  }
  child[concept color=popgreen]{ node{Defining vs non-defining}
    child{ node{commas = extra info} }
    child{ node{no \emph{that} after comma} }
  }
  child[concept color=popgold]{ node{Noun clauses}
    child{ node{that, if, whether} }
    child{ node{wh- words} }
  }
  child[concept color=poppurple]{ node{Adverbial clauses}
    child{ node{time, cause, purpose} }
    child{ node{condition, concession} }
  }
  child[concept color=poppink]{ node{Word order}
    child{ node{no inversion} }
    child{ node{no future after when/if} }
  }
  child[concept color=popblue!60]{ node{Francophone traps}
    child{ node{preposition + whom/which} }
  };
\end{tikzpicture}
\legende{Carte mentale : les trois grandes familles de subordonnées et leurs pièges.}
\end{popfigure}

\begin{definition}[L'essentiel de la leçon : lexique clé (English -- Français)]
\begin{itemize}
\item \eng{clause} : proposition (sujet + verbe conjugué)
\item \eng{main clause} : proposition principale ; \eng{subordinate clause} : proposition subordonnée
\item \eng{relative pronoun} : pronom relatif ; \eng{antecedent} : antécédent (le nom auquel renvoie le pronom relatif)
\item \eng{defining relative clause} : relative déterminative ; \eng{non-defining relative clause} : relative explicative
\item \eng{noun clause} : subordonnée complément (qui fonctionne comme un nom)
\item \eng{adverbial clause} : subordonnée circonstancielle (temps, cause, but, concession, condition)
\item \eng{linkword, connective} : mot de liaison
\end{itemize}
\end{definition}

\begin{aretenir}[Règles à connaître par cœur]
\begin{center}
\begin{tabularx}{\linewidth}{|l|X|X|}
\hline
\rowcolor{popblueL} Type & Marqueurs & Exemple \\
\hline
Relative (personne) & who, that, whom & \eng{The nurse who treated me was kind.} \\
\hline
Relative (chose) & which, that & \eng{The report which she wrote is clear.} \\
\hline
Possession & whose & \eng{The boy whose father is a doctor...} \\
\hline
Lieu / temps & where, when & \eng{The hospital where I was born.} \\
\hline
Explicative & , who/which , (jamais \emph{that}) & \eng{Lake Nyos, which is in Cameroon, is dangerous.} \\
\hline
Noun clause & that, if, whether, wh- & \eng{I know that pollution kills.} \\
\hline
Temps & when, while, as soon as, until & \eng{When the rain stops, we will plant.} \\
\hline
Cause & because, since, as & \eng{Since the water is dirty, boil it.} \\
\hline
But & so that, in order to & \eng{He trains so that he can win.} \\
\hline
Concession & although, even though, whereas & \eng{Although he is ill, he works.} \\
\hline
Condition & if, unless, provided that & \eng{If we recycle, the city stays clean.} \\
\hline
\end{tabularx}
\end{center}

\textbf{Précisions importantes :}
\begin{itemize}
\item \eng{whose} peut renvoyer à une chose, pas seulement à une personne : \eng{a country whose forests are disappearing} « un pays dont les forêts disparaissent ».
\item \eng{so that} est suivi d'une clause (sujet + verbe) ; \eng{in order to} est suivi d'un infinitif : \eng{He trains in order to win.}
\item \eng{despite} et \eng{in spite of} sont suivis d'un nom ou d'un gérondif, jamais d'une clause complète : \eng{Despite his illness, he works.} / \eng{Despite being ill, he works.}
\end{itemize}
\end{aretenir}

\begin{methode}[Méthodes express]
\begin{enumerate}
\item Repérer le verbe conjugué : chaque verbe conjugué = une clause. Dans \eng{The nurse who treated me was kind}, il y a deux verbes conjugués (\eng{treated}, \eng{was}), donc deux clauses.
\item Pour une relative : identifier l'antécédent (personne $\rightarrow$ \eng{who} ; chose $\rightarrow$ \eng{which} ; les deux $\rightarrow$ \eng{that}).
\item Information indispensable = pas de virgule ; information supplémentaire = virgules + jamais \eng{that}.
\item Après \eng{when, if, as soon as, unless, until} : présent pour exprimer le futur (\eng{When he arrives}, jamais \emph{when he will arrive}).
\item Ordre des mots dans une noun clause : sujet + verbe, sans inversion (\eng{I don't know where he lives.}).
\end{enumerate}
\end{methode}

\begin{attention}[Pièges fréquents des francophones]
\begin{itemize}
\item Ne pas traduire « que » systématiquement par \eng{that} : pour une personne sujet, préférer \eng{who} (\eng{The man who lives next door}, et non \emph{the man which lives next door}).
\item Ne jamais mettre \eng{that} après une virgule dans une relative explicative : \eng{My uncle, who lives in Buea, is a teacher.} (jamais \emph{my uncle, that lives...}).
\item Ne pas répéter le complément après le pronom relatif : \eng{The book which I bought is expensive.} (jamais \emph{which I bought it}).
\item Question indirecte : pas d'inversion, pas d'auxiliaire \eng{do} : \eng{Tell me where the clinic is.} (jamais \emph{where is the clinic} dans la subordonnée).
\item \eng{Despite} n'accepte pas \eng{of} seul : on dit \eng{despite the rain} ou \eng{in spite of the rain}, jamais \emph{despite of the rain}.
\end{itemize}
\end{attention}

\begin{aretenir}[Checklist avant le Bac]
\begin{itemize}
\item[$\square$] Je sais distinguer une main clause d'une subordinate clause.
\item[$\square$] Je choisis correctement \eng{who}, \eng{which}, \eng{that}, \eng{whose}, \eng{where}, \eng{when}.
\item[$\square$] Je n'utilise jamais \eng{that} dans une relative explicative (entre virgules).
\item[$\square$] Je sais que \eng{whose} peut renvoyer à une chose : \eng{a country whose forests are disappearing}.
\item[$\square$] Je place la préposition correctement : \eng{the man to whom I spoke} (soutenu) / \eng{the man (who) I spoke to} (courant).
\item[$\square$] Je mets le présent (et non le futur) après \eng{when}, \eng{if}, \eng{as soon as}, \eng{until}.
\item[$\square$] Je n'inverse pas sujet et verbe dans une question indirecte : \eng{Tell me where the clinic is.}
\item[$\square$] Je distingue \eng{because} (cause), \eng{so that} (but) et \eng{although} (concession).
\item[$\square$] Je sais que \eng{despite} et \eng{in spite of} sont suivis d'un nom, pas d'une clause.
\item[$\square$] Je peux transformer deux phrases simples en une phrase complexe avec une relative.
\item[$\square$] Je relie mes idées dans un paragraphe du Bac avec au moins trois connecteurs variés.
\end{itemize}
\end{aretenir}

\findecours

\end{document}
